\documentclass{amsart}
\usepackage{amssymb,mathtools,array}
\usepackage[hidelinks]{hyperref}
\hypersetup{pdftitle={Cubic sieves, a primitive-ratio cut and a 20/37 bound for the abc exceptional set},pdfauthor={Ruiyi Zhang}}

\newtheorem{theorem}{Theorem}
\newtheorem{theoremA}{Theorem}
\renewcommand{\thetheoremA}{\Alph{theoremA}}
\newtheorem{lemma}{Lemma}[section]
\newtheorem{proposition}[lemma]{Proposition}
\newtheorem{corollary}[lemma]{Corollary}
\theoremstyle{definition}
\newtheorem{definition}[lemma]{Definition}
\theoremstyle{plain}
\numberwithin{equation}{section}
\numberwithin{table}{section}

\newcommand{\rad}{\operatorname{rad}}
\newcommand{\Z}{\mathbb{Z}}
\newcommand{\N}{\mathbb{N}}
\newcommand{\pos}[1]{\max\{#1,0\}}

\newcommand{\certLeaves}{394\,633}
\newcommand{\certInternal}{394\,243}
\newcommand{\certNodes}{788\,876}
\newcommand{\certDepth}{49}
\newcommand{\certInfeasible}{0}
\newcommand{\certLambda}{586/201}
\newcommand{\certLambdaFrac}{\frac{586}{201}}
\newcommand{\certLambdaLeaf}{95\,263}
\newcommand{\certLambdaRho}{\frac{3516}{67}}
\newcommand{\certLambdaNineteen}{\frac{11134}{201}}
\newcommand{\certBudget}{\frac{11335}{201}}
\newcommand{\certTight}{1\,358}
\newcommand{\certMaxTerms}{28}
\newcommand{\certUNused}{17\,256}
\newcommand{\certMut}{46}
\newcommand{\certWitMut}{19}
\newcommand{\certFamilyTable}{%
linear rows of $\mathcal S$ & & & 29 & 2\,061\,265\\
nonnegativity & & & 21 & 2\,478\,333\\
\eqref{eq:ARCH} & 4\,972 & 1 & 2 & 26\,673\\
(G1)--(G13) & 27\,223 & 13 & 26 & 675\,941\\
\eqref{eq:SPor} & 15\,468 & 30 & 60 & 479\,423\\
\eqref{eq:cutor} & 236\,038 & 17\,256 & 34\,116 & 2\,417\,940\\
\eqref{eq:mo:CS4or} & 48\,765 & 10\,649 & 21\,220 & 644\,485\\
\eqref{eq:mo:PS23or} & 61\,110 & 155 & 310 & 673\,495\\
\eqref{eq:pw:POW3} & 387 & 75 & 149 & 1\,924\\
\eqref{eq:pw:G12}, \eqref{eq:pw:G1212} & 280 & 27 & 89 & 2\,679\\}
\newcommand{\shaManifestI}{7ba2202c473e0204baf7752a8fd41923}
\newcommand{\shaManifestII}{82111ca4cd189a6db584e3912a3bcf43}
\newcommand{\shaReportXzI}{f0412ad682389d463d410b96737144e6}
\newcommand{\shaReportXzII}{65afba13ef26c29f64c7a13f128e24e0}
\newcommand{\shaReportI}{9cf43afd90bb683897290c18a9ebfb58}
\newcommand{\shaReportII}{697d0baafe8c9d2c6ae6ff71e8f5e2f7}
\newcommand{\shaNormI}{f38667a2fe45bda186c0ef0ae827f11c}
\newcommand{\shaNormII}{8c68cfac0976da043f7eba59bbd9c701}
\newcommand{\shaCleanI}{9c301c46a2561e2d3c7c27a2c258186d}
\newcommand{\shaCleanII}{66703d89a52f86d74c2f1e6604ad2d0d}
\newcommand{\shaTreeI}{d7abb356e4098127e775c4f003ef4bd8}
\newcommand{\shaTreeII}{2660cbd734547154ab5f8b3f1b394f4a}
\newcommand{\shaWitI}{b795b735b3ec81177b57f6402b746d26}
\newcommand{\shaWitII}{2b5fae0a102eadc5f6b07777b9944cf3}
\newcommand{\shaDiffI}{39ece2ea82da2cd5e221624707de6f0b}
\newcommand{\shaDiffII}{4d5e7cd18a6599b3e53f02452f9072e3}
\newcommand{\shaAuditI}{ea40e6d7cbac17b5a2662b31e4d2b96a}
\newcommand{\shaAuditII}{2ee3a1805ae137892d947311d15df929}
\newcommand{\shaRunI}{dacd0fe7868341f22c4ed5b94a284492}
\newcommand{\shaRunII}{df3c8c45dc4f9da578800ba8b0db215b}
\newcommand{\shaStdoutI}{3c7494813ddb89721817b77f48bdac05}
\newcommand{\shaStdoutII}{5f9bef2f298b3ca7291799e4565df947}

\newif\ifarxiv
\arxivtrue

\begin{document}

\title[The $20/37$ bound for the $abc$ exceptional set]{Cubic sieves, a primitive-ratio cut and a $20/37$ bound for the $abc$ exceptional set}
\author{Ruiyi Zhang}
\address{Independent Researcher, Columbia, Maryland, USA}
\email{ryzfsu@gmail.com}
\subjclass[2020]{Primary 11D75; Secondary 11N25, 11N36, 11L40, 11D25, 11Y16}
\keywords{abc conjecture, exceptional set, radical, anatomy of integers, power sieve, character sums, Farey fractions, linear programming certificate}

\begin{abstract}
We prove $N_1(X)\ll_\varepsilon X^{20/37+\varepsilon}$, where $N_\lambda(X)$
counts the $abc$ triples of exponent $\lambda$ in $[1,X]^3$. The previous full
bound was $X^{13/20+\varepsilon}$. We use the
anatomic reduction of Bernert, Browning, Lichtman and Ter\"av\"ainen. Five
families add new rows. The first two are an archimedean cover and a
row from Heath-Brown's uniform bound for diagonal cubic curves. Spacing gives $30$ two-arm clauses.
Cubic characters and a fourth moment of Dirichlet $L$-functions give the
cubic-sieve rows. Cochrane and Shi's fourth moment gives the moment rows.
An exact Farkas certificate with $\certLeaves$ leaves gives exponent $20/37$
for the box count. On each leaf the multipliers of all rows except
nonnegativity sum to at most $\certLambda$. Exact rows are included in this
sum. This pays for the losses at $\lambda=1$.
\end{abstract}

\maketitle

\section{Introduction}\label{sec:intro}

For this paper $\N=\{1,2,\ldots\}$. For $n\ne0$ put
$\rad(n)=\prod_{p\mid n}p$, and write $\omega(n)=\#\{p:p\mid n\}$,
with $\omega(1)=0$. Following \cite{BBLT}, call
$(a,b,c)\in\N^3$ with $\gcd(a,b,c)=1$ an \emph{$abc$ triple of exponent
$\lambda$} if
\begin{equation}\label{eq:triple}
 a+b=c,\qquad \rad(abc)<c^{\lambda}.
\end{equation}
Let $N_\lambda(X)$ count these triples in $[1,X]^3$. The $abc$ conjecture of
Masser and Oesterl\'e implies that $N_\lambda(X)=O_\lambda(1)$ for $\lambda<1$,
so $\lambda=1$ is the endpoint of interest.

The bound $N_\lambda(X)\ll_\varepsilon X^{2\lambda/3+\varepsilon}$ is
classical \cite[Proposition~1.1]{BBLT}; it comes from de Bruijn's count
\cite{dB} of integers with restricted radical. For $\lambda\ge2$, Kane's result
\cite{Ka} gives $N_\lambda(X)\ll_\varepsilon X^{\lambda-1+\varepsilon}$
\cite[Section~1]{BBLT}. On $2\le\lambda\le3$ it matches
$N_\lambda(X)\gg X^{\lambda-1-\varepsilon}$. Granville's unpublished work gives the lower bound for $1<\lambda\le3$, as
reported in \cite[Section~1]{BBLT}. The question is Mazur's. It can hold only for
$\lambda\le3$, since
$N_\lambda(X)\le\#\{(a,b,c)\in[1,X]^3:a+b=c\}=O(X^2)$. Browning, Lichtman and
Ter\"av\"ainen \cite[Theorem~1.2]{BLT} gave the first power saving below
$2\lambda/3$, for $0.99<\lambda<1.001$. Bernert \cite[Theorem~1]{Be} proved
\begin{equation}\label{eq:bernertbound}
 N_\lambda(X)\ll_\varepsilon X^{(23\lambda+3)/40+\varepsilon}
 \qquad(0<\lambda<2),
\end{equation}
so $N_1(X)\ll_\varepsilon X^{13/20+\varepsilon}$. In the merged paper,
\cite[Theorem~1.2]{BBLT} states \eqref{eq:bernertbound} for $0<\lambda\le2$.
Its Theorem~1.3 gives $0.6$ for $0<\lambda<1$, with a constant depending on
$\lambda$. At $\lambda=1$: a remark at the end of
\cite[Section~4]{Be} sketches $8/13=0.6153\ldots$; Li \cite{Li} obtains
$56/85=0.6588\ldots$ from the method of \cite{BLT}, its limit without further
optimisation by his account.

\begin{theorem}\label{thm:main}
$N_1(X)\ll_\varepsilon X^{20/37+\varepsilon}$, where $20/37=0.5405\ldots$
\end{theorem}

Since $N_\mu\le N_1$ for $\mu\le1$, Theorem~\ref{thm:main} also improves
\cite[Theorem~1.3]{BBLT} on $0<\mu<1$, with a constant independent of $\mu$.
Among the estimates cited here, the resulting exponent is
$\min\{2\mu/3,20/37\}$ for $0<\mu\le1$. Thus the new exponent is best known
for $30/37<\mu\le1$; it ties the classical exponent at $\mu=30/37$.
Its crossing with the Bernert exponent alone is $\mu=689/851$.

Section~\ref{sec:red} recalls the anatomic reduction of \cite{BLT,Be,BBLT}.
Fix $\varepsilon_0>0$ and $d=d(\varepsilon_0)$. Every triple lies in a box.
In a box, each side is a fixed coefficient times
$\prod_{i\le d}\mathfrak s_i^{\,i}$. The roots $\mathfrak s_i$ have fixed
dyadic sizes. The three side values are pairwise coprime. Let $\mathcal B_d$
count the root-factorization tuples in a box, with distinct factorizations
counted separately, and put $D=\log_X\mathcal B_d$. The proof is the chain.
Section~\ref{sec:red} works at a generic scale $X$ and exponent $\lambda$;
Lemma~\ref{lem:endpoint} applies it at the scale $Y$ of \eqref{eq:dyadic},
with $\lambda:=\mu$ and base-$Y$ logarithms. The quantity $V_s$ is the
weighted exponent of the root product, excluding its fixed coefficient, as
defined in Section~\ref{sec:arch}; $h_d=d(d+1)/2$. The quantity $\rho$ is the
largest finite right-hand-side shift in Lemma~\ref{lem:endpoint}. The bound
on $\rho$ holds for $Y\ge Y_0(\varepsilon_0,d)$. Smaller $Y$ contribute
$O_{\varepsilon_0,d}(1)$, and $\mu\le0$ contributes nothing
(Section~\ref{sec:proof}).
\[
 \begin{gathered}
 N_1(X)\ll(\log X)^2\max_{Y\le X,\ \mu\le1}N^*_\mu(Y),\qquad
 N^*_\mu(Y)\ll Y^{\varepsilon_0}\max\mathcal B_d,\\
 D\le\frac{20}{37}+\Lambda_U\rho,\qquad \Lambda_U=\certLambdaFrac,
 \qquad \rho\le19\varepsilon_0,
 \end{gathered}
\]
from \eqref{eq:dyadic}, \eqref{eq:reduction}, Theorem~\ref{thm:cert} with
Lemma~\ref{lem:stab}, and Lemma~\ref{lem:endpoint}. Since
$19\Lambda_U+1=\certBudget$, choosing $\varepsilon_0$ with
$\certBudget\,\varepsilon_0<\varepsilon/2$ proves Theorem~\ref{thm:main}.

The inequalities for $D$ are linear in $28$ exponents. They are the root and
tail coordinates of the three sides, $D$, and three auxiliary Fourier
variables. BBLT's Theorem~1.3 gives the $3/5$ bound. Its Proposition~6.1 lists
the $22$ inherited inequalities used here. We add the tail
rows and the cut. Tail roots
(Section~\ref{sec:tail}) of index at least seven cost $X^{\tau_s}$ to
enumerate. Here $\tau_s=s-\sum_{i\le6}s_i$. The tail product is at least
$X^{v_s}$, where $v_s\ge7\tau_s$. The cut (Section~\ref{sec:cut}) fixes all
but two coprime products $u,v$ on two sides:
\[
 \gcd(u,v)=1,\qquad u\equiv\xi v\pmod M,\qquad
 0<u\le U_{\rm cut},\quad 0<v\le V_{\rm cut}.
\]
Here $M$ is built from roots on the third side, and $\xi$ lies in one of
$X^{o(1)}$ classes. Lemma~\ref{lem:farey} counts these by
$1+2U_{\rm cut}V_{\rm cut}/M$. There are
\begin{equation}\label{eq:cutcount}
 3\,(64^2-1)\,127=1\,560\,195
\end{equation}
cuts.

The paper adds five further families.
\begin{itemize}
\item \emph{Archimedean} (Section~\ref{sec:arch}). One of the side values
 $n_a,n_b$ is at least $n_c/2$. In the limiting coordinates, $V_a,V_b\le V_c$, and
 $V_c\le V_a$ or $V_c\le V_b$. At finite scale,
 $V_a,V_b\le V_c+\varepsilon_0/2$, and one of
 $V_c\le V_a+3\varepsilon_0/2+\log_Y2^{1+h_d}$ or
 $V_c\le V_b+3\varepsilon_0/2+\log_Y2^{1+h_d}$ holds.
\item \emph{Cubic} (Section~\ref{sec:hb3}). Fix all roots but the three
 cube roots. Heath-Brown's bound \cite[Theorem~1.3]{HBcube} for primitive
 points on $\alpha x^3+\beta y^3=\gamma z^3$, uniform in the coefficients,
 gives $D\le L-p_3$. Here $L$ and $p_i$ are as in Section~\ref{sec:red}.
\item \emph{Spacing} (Section~\ref{sec:spacing}). Consecutive $\ell$th powers
 are far apart: $30$ two-arm clauses.
\item \emph{Cubic sieve} (Section~\ref{sec:pow3}). Write one side as
 $Mx^3$. Sieve with cubic characters modulo primes $p\equiv1\pmod3$. The
 error term needs a smooth fourth moment of character sums. It comes from the
 fourth moment of Dirichlet $L$-functions in \cite[Lemma~4]{GMRR}: $1774$ clauses.
\item \emph{Moments} (Section~\ref{sec:moment}). Expand the congruence of
 the cut modulo $M$ by characters. The fourth moment of Cochrane and Shi
 \cite[Theorem~1]{CS} bounds the nonprincipal part: $188\,976$ clauses.
\end{itemize}

Section~\ref{sec:cert} proves $D\le20/37$ from these rows. The certificate is
a finite trie. Each internal node splits on one clause, one child per arm.
Each leaf carries nonnegative rational Farkas multipliers. There are
$\certLeaves$ leaves, depth $\certDepth$, and no infeasible leaf. On a leaf
the multipliers of all rows except nonnegativity sum to at most
$\Lambda_U=\certLambda$, whence the term $\Lambda_U\rho$. This includes exact
rows. A point of the
installed system attains $20/37$; the qualification on this optimum appears in
Section~\ref{sec:witness}.

Section~\ref{sec:end} treats $\mu=1$. In the proof of Theorem~\ref{thm:main},
the dyadic reduction only requires $0<\mu\le1$. At this endpoint the row
$L\le1$ of \cite{BBLT} is unavailable: the proof of
\cite[Theorem~1.3]{BBLT} uses $\varepsilon_0\le(1-\mu)/3$. We use
$L\le1+3\varepsilon_0/4$ instead. The $L$-native Fourier and geometry rows cost
$18\varepsilon_0$ and $6\varepsilon_0$. After the choice of $Y_0$ in
Lemma~\ref{lem:endpoint}, every new row costs at most $2\varepsilon_0$.
Hence $\rho=18\varepsilon_0+o_{\varepsilon_0,d}(1)$ and
$\Lambda_U\rho=\certLambdaRho\,\varepsilon_0+o_{\varepsilon_0,d}(1)$.

The proof in \cite{BBLT} does not use the determinant method; \cite{BLT} used
it together with Thue equations. Lemma~\ref{lem:farey} is an elementary
lattice-point count of the kind used in \cite[Proposition~4.1]{BBLT}. The cut
applies it to two sides at once, with modulus $M$ built from roots on the third
side.

Imported statements are cited where used. Citations of \cite{BBLT} refer to
arXiv:\allowbreak2410.12234v2 (to appear in C.~R.~Math.); each imported
statement was checked against it. The dyadic reduction \eqref{eq:dyadic} is
only stated in \cite[Section~2]{BBLT}; Section~\ref{sec:red} proves it.
Appendix~\ref{app:windows} reproves the estimates of
\cite[Propositions~2.3, 3.1 and~4.1]{BBLT} for our boxes. They use the
divisor bound, the count of solutions of $x^j-y^j=m$
(Lemma~\ref{lem:app:power}) and the lattice-point bound \cite[Lemma~3]{HB}
(Theorem~\ref{thm:app:HB}). The new rows use
the cubic bound \cite[Theorem~1.3]{HBcube} and the fourth moments
\cite[Lemma~4]{GMRR} and \cite[Theorem~1]{CS}. The cubic-sieve character expansion
(Section~\ref{sec:pow3}) also uses the classical Jacobi-sum bound \cite{IR}.
No conjecture is assumed. Theorem~\ref{thm:cert}
has multiplier bound $\Lambda_U$. It is computer-assisted
and uses a certificate checked by the standard-library programs in
Appendix~\ref{app:ver}. The analytic reductions and row families are proved in
the text.


\section{The anatomic reduction}\label{sec:red}

Throughout this section $0<\lambda\le2$, $\varepsilon>0$ is arbitrary,
$\varepsilon_0>0$ is fixed before $X$, and $d=d(\varepsilon_0)$. We take
$\varepsilon_0\le4/7$, so that the folded index $k=\lceil4/\varepsilon_0\rceil$
of \cite[Lemma~2.2]{BBLT} is at least $7$ (Section~\ref{sec:fold}).
Implied constants depend only on $\varepsilon,\varepsilon_0,d$; $\lambda$ occurs
only through the displayed $\ll_\lambda$ in \eqref{eq:dyadic}. They never depend on the box,
the coefficients or $X$.

Write $N^*_\lambda(X)$ for the number of triples counted by $N_\lambda(X)$ that
satisfy in addition $c\ge X/2$ and $\rad(abc)\ge c^{\lambda}/2$. A dyadic
decomposition and the pigeonhole principle give
\begin{equation}\label{eq:dyadic}
 N_\lambda(X)\ll_\lambda(\log X)^2\max_{1\le Y\le X}\ \max_{\mu\le\lambda}N^*_\mu(Y),
\end{equation}
as in \cite[Section~2]{BBLT}. In the application $\lambda=1$, and the bounds
below for $N^*_\mu(Y)$ are uniform in $0<\mu\le1$
(Lemma~\ref{lem:endpoint}).

\begin{proof}[Proof of \eqref{eq:dyadic}]
Let $(a,b,c)$ be counted by $N_\lambda(X)$. Then $c\ge2$ and $a,b<c\le X$.
Choose $l\ge1$ with $c\in(X2^{-l},X2^{1-l}]$ and put
$Y=X2^{1-l}\le X$. Thus $Y/2<c\le Y$. Put
\[
\mu_i=\lambda-i/\log_2X\qquad
 (0\le i\le\lceil\lambda\log_2X\rceil).
\]
Here $\log_2X=(\log X)/(\log2)$.
There are $O(\log X)$ such exponents, all at most $\lambda$, and the last is
at most $0$. With $r=\rad(abc)$, we have
$r<c^{\mu_0}$ and $r\ge1\ge c^{\mu_m}$ for the last index $m$.
Choose the largest $i$ with $r<c^{\mu_i}$. Then
$r\ge c^{\mu_{i+1}}\ge c^{\mu_i}/2$, since $c\le X$.
Hence the triple is counted by $N^*_{\mu_i}(Y)$. There are $O(\log X)$ choices
of $Y$ and $O(\log X)$ choices of $\mu_i$, which proves the claim.
\end{proof}

For $\mu\le0$, $N^*_{\mu}(Y)=0$: its defining condition would require
$1\le\rad(abc)<c^\mu\le1$. Thus the maximum in \eqref{eq:dyadic} may be
restricted to $0<\mu\le\lambda$.

For $n\ge1$ write $[n]=\{1,\ldots,n\}$. Here $x\sim A$ means $A\le x<2A$, as
in \cite{Be}; \cite{BBLT} writes $x\sim A$ for $x\in(A,2A]$.

Let $d\in\N$ and $(a_0,b_0,c_0)\in\N^3$. Let $\mathbf A,\mathbf B,\mathbf C\in\N^d$
have entries powers of two, $1$ included. Put $\mathfrak a_0=a_0$,
$\mathfrak b_0=b_0$, $\mathfrak c_0=c_0$ and
\begin{equation}\label{eq:box}
 \mathcal B_d=\#\Bigl\{(\mathfrak a_i,\mathfrak b_i,\mathfrak c_i)_{i=1}^d\in\N^{3d}:
 \begin{aligned}
  &\mathfrak a_i\sim A_i,\ \mathfrak b_i\sim B_i,\ \mathfrak c_i\sim C_i\quad(1\le i\le d),\\
  &a_0\textstyle\prod_{i\le d}\mathfrak a_i^{\,i}+b_0\prod_{i\le d}\mathfrak b_i^{\,i}
    =c_0\prod_{i\le d}\mathfrak c_i^{\,i},\\
  &\gcd\bigl(\textstyle\prod_{i=0}^{d}\mathfrak a_i,
    \prod_{i=0}^{d}\mathfrak b_i,\prod_{i=0}^{d}\mathfrak c_i\bigr)=1
 \end{aligned}\Bigr\}.
\end{equation}
For a side write $S_j\in\{A_j,B_j,C_j\}$ for its dyadic root size. When the
side $s$ must be named, write $S_{s,j}$. Recall $h_d=d(d+1)/2$.
Equation~\eqref{eq:box} counts root-factorization tuples, not distinct
side-value triples; different factorizations are counted separately. It is
the $B_d$ of \cite{BBLT}, under our dyadic convention, with no condition within
a side. We write $\mathcal B_d$ to avoid a clash with the size $B_d$.

Every interval $[A,2A)\cap\N$ has exactly $A$ integers, including
$[1,2)\cap\N=\{1\}$. We keep every root coordinate, including those equal to
$1$. The statements of \cite[Propositions~2.3, 3.1 and~4.1]{BBLT} assume
windows $(A,2A]$ with $A\ge1$, which exclude roots equal to $1$. So we do not
cite them for our boxes. We rerun the window-dependent steps of the reduction
below.
Appendix~\ref{app:windows} proves the pair, Fourier and geometry bounds in our
convention (Lemmas~\ref{lem:app:pair}, \ref{lem:app:fourier}
and~\ref{lem:app:geom}). Their constants depend only on $\varepsilon_0$
and $d$.

Roots equal to $1$ occur: in
\cite[Lemma~2.2]{BBLT}, $m_j=\prod_{p^{j}\|n}p$ and $n_j=m_j$ for $j\ne k$, the
folded index. So $n_j=1$ unless some prime has multiplicity exactly $j$. As
$n$ has at most $\omega(n)$ distinct multiplicities, all but at most
$\omega(n)$ indices $j\le d$, $j\ne k$, give $n_j=1$. Our convention assigns
these roots the size $A_i=1$ and keeps all $d$ indices in \eqref{eq:box}.
By the proof of \cite[Proposition~2.1]{BBLT}, each $\varepsilon_0>0$ has a
$d=d(\varepsilon_0)$ with
\begin{equation}\label{eq:reduction}
 N^*_\lambda(X)\ll_{\varepsilon_0}X^{\varepsilon_0}\max \mathcal B_d,
\end{equation}
the maximum over pairwise coprime $1\le a_0,b_0,c_0\le X^{\varepsilon_0}$ and
dyadic $A_i,B_i,C_i\in[1,X]$ with
\begin{equation}\label{eq:boxconstraints}
 \begin{gathered}
  2^{-O_d(1)}X^{\lambda}\le\prod_{j\le d}A_jB_jC_j\le X^{\lambda+3\varepsilon_0/4},\\
  \prod_{j\le d}A_j^{\,j}\le X,\qquad
  \prod_{j\le d}B_j^{\,j}\le X,\qquad
 X^{1-\varepsilon_0/2}\le\prod_{j\le d}C_j^{\,j}\le X ,
 \end{gathered}
\end{equation}
The power of two depends only on $d$. We rederive the lower bounds
below. The upper bounds and the pigeonhole step of the source proof hold
verbatim under $[A,2A)$, since $A_i\le\mathfrak a_i<2A_i$. The lower bound uses the
nontrivial branch of \cite[Lemma~2.2]{BBLT}: $\varepsilon_0<4$, so the
parameter $\varepsilon_0/4<1$ and $k=\lceil4/\varepsilon_0\rceil\le d$.
Since $c\ge X/2$,
\[
 \prod_{j\le d}A_jB_jC_j\ge2^{-3d}\rad(abc)
 \ge2^{-3d-\lambda-1}X^\lambda\ge2^{-3d-3}X^\lambda,
\]
since $\lambda\le2$. Each of the $3d$ roots loses at most a factor $2$; the
source proof prints $2^{-d}$ here. The weighted lower bound for the $c$-roots is
\[
 \prod_{j\le d}C_j^j\ge2^{-h_d}\prod_{j\le d}\mathfrak c_j^j.
\]
The factorization uses parameter $\varepsilon_0/4$, so
$c_0\le X^{\varepsilon_0/4}$ and
$\prod_{j\le d}\mathfrak c_j^{\,j}=c/c_0\ge2^{-1}X^{1-\varepsilon_0/4}$.
Since $C_j\ge\mathfrak c_j/2$,
\[
 \prod_{j\le d}C_j^j\ge2^{-h_d-1}X^{1-\varepsilon_0/4}
 \ge X^{1-\varepsilon_0/2}
 \quad\text{if }X\ge Y_{\rm red}:=2^{4(h_d+1)/\varepsilon_0}.
\]
This gives the lower bound in \eqref{eq:boxconstraints}. The global
threshold $Y_0(\varepsilon_0,d)\ge Y_{\rm red}$ is fixed in
Lemma~\ref{lem:endpoint}.
This suffices, since $d$ is fixed before $X$. Fix such a box. For a side $s\in\{a,b,c\}$ set
\begin{equation}\label{eq:coord}
 s_i=\log_XS_i\ (1\le i\le d),\qquad
 s=\log_X\!\prod_{i\le d}S_i=\sum_{i\le d}\log_XS_i .
\end{equation}
Index $0$ is not a coordinate: $a_0,b_0,c_0$ are the integer coefficients.
From here on $a,b,c$ denote coordinates, not the triple.
Boxes with $\mathcal B_d=0$ do not attain the maximum in \eqref{eq:reduction}
unless $N^*_{\lambda}(X)=0$. So $\mathcal B_d\ge1$, and we set
\begin{equation*}
 D=\log_X\mathcal B_d\ge0.
\end{equation*}
For the tail put
\begin{equation}\label{eq:tailcoord}
 \tau_s=s-\sum_{i\le6}s_i=\sum_{7\le j\le d}\log_XS_j,
 \qquad
 v_s=\sum_{7\le j\le d}j\log_XS_j .
\end{equation}
Put $L=a+b+c$ and $p_i=a_i+b_i+c_i$; a bare $p$ is a prime. Moving a
dyadic endpoint at scale $X$ changes the root coordinates by
\begin{equation}\label{eq:dyadicshift}
 \Delta s_i,\ \Delta s=O_d\!\left(\frac1{\log X}\right).
\end{equation}

The left inequality of \eqref{eq:boxconstraints} uses
$\prod_{j\le d}n_j\ge\rad(n)$ with no loss, while in this direction
\cite[Lemma~2.2]{BBLT} gives only $\rad(n)\le n^{\varepsilon_0}\prod_jn_j$. In the proof of that
lemma
$m_j=\prod_{p^{j}\|n}p$, and, for the folded index $k$, $n_j=m_j$ for
$j\ne k$ and $n_k=m_k\prod_{j>d}m_j^{\lfloor j/k\rfloor}$. Since $d\ge k$,
$\lfloor j/k\rfloor\ge1$ for $j>d$, so
\[
 \prod_{j\le d}n_j\ \ge\ \prod_{j\le d}m_j\prod_{j>d}m_j\ =\ \rad(n).
\]

\begin{lemma}\label{lem:coprime}
The three side values of a box are pairwise coprime. Hence every factor of one
side is a unit modulo every factor of another side.
\end{lemma}

\begin{proof}
 A prime dividing two side values divides the third by the equation in
\eqref{eq:box}. If it divides
$a_0\prod_{i=1}^d\mathfrak a_i^{\,i}$, it divides $\prod_{i=0}^d\mathfrak a_i$;
likewise on the other sides. This contradicts the gcd condition in
\eqref{eq:box}.
\end{proof}

The index $i=0$ in this gcd condition is needed. With $d=1$, the choice
\[
 a_0=b_0=1,\quad \mathfrak a_1=\mathfrak b_1=2,\quad c_0=4,\quad \mathfrak c_1=1
\]
satisfies the condition with $i\ge1$: $\gcd(2,2,1)=1$. Yet its side values
$2$, $2$, $4$ share the prime $2$. The
full condition excludes it: $\gcd(2,2,4)=2$.

Nothing is assumed within a side. Roots of one side may share primes with
$a_0$. Unenumerated roots are recovered by a divisor bound, not by
squarefreeness.

\subsection{The inherited rows}
For these limiting rows take $0<\lambda\le1$. The finite box bound is
$L\le\lambda+3\varepsilon_0/4\le1+3\varepsilon_0/4$. We use this endpoint
envelope throughout, so no choice of $\varepsilon_0$ depends on the gap
$1-\lambda$.
Lemmas~\ref{lem:app:pair} and~\ref{lem:app:fourier} and
Corollary~\ref{cor:app:geom} are our forms of
\cite[Propositions~2.3, 3.1 and 4.1]{BBLT}. In the coordinates \eqref{eq:coord}
they yield the $22$ inequalities of \cite[Proposition~6.1]{BBLT}. We display them in exact limiting
form, as Section~\ref{sec:cert} uses them. At finite resolution each
right-hand side rises by the amount in Lemma~\ref{lem:endpoint} and
Table~\ref{tab:rho}; \eqref{eq:root} and \eqref{eq:63} hold exactly.
The reduction has constants depending only on $\varepsilon_0,d$. So do the
bounds of Appendix~\ref{app:windows}. They use $\lambda$ only through
$ABC\le X^{\lambda+3\varepsilon_0/4}$ in Corollary~\ref{cor:app:geom}. The
condition $\varepsilon_0\le(1-\mu)/3$ occurs only in the final step of the
proof of \cite[Theorem~1.3]{BBLT}, where it ensures
$ABC\le X^{\mu+3\varepsilon_0}\le X$.

The radical bound in \eqref{eq:boxconstraints} gives \cite[(6.1)]{BBLT},
\begin{equation}\label{eq:L}
 L\le1,
\end{equation}
derived there for $\lambda<1$. At $\lambda=1$ only $L\le1+3\varepsilon_0/4$
holds (Lemma~\ref{lem:endpoint}). Roots sum into their side,
\cite[(6.2)]{BBLT}:
\begin{equation}\label{eq:root}
 \sum_{i\le6}s_i\le s\qquad(s=a,b,c).
\end{equation}
The side-value bound in \eqref{eq:boxconstraints} gives \cite[(6.3)]{BBLT},
\begin{equation}\label{eq:63}
 7s-1\le6s_1+5s_2+4s_3+3s_4+2s_5+s_6\qquad(s=a,b,c).
\end{equation}
Two sides determine the third up to a divisor factor
(Lemma~\ref{lem:app:pair}). This is \cite[(6.4)]{BBLT}:
\begin{equation}\label{eq:pair}
 D\le\min(a+b,\ a+c,\ b+c).
\end{equation}
Geometry of numbers (Corollary~\ref{cor:app:geom}) gives the thirteen
min-disjunctions \cite[(6.5)--(6.9)]{BBLT}, labelled (G1)--(G13) in
Appendix~\ref{app:geom}. The Fourier estimate (Lemma~\ref{lem:app:fourier})
gives \cite[(6.10)]{BBLT},
\begin{equation}\label{eq:fourier}
 \sum_{s=a,b,c}\max(s_2+s_4+s_6,\ s_3+s_6,\ s_5)\le4L-6D .
\end{equation}
One argument per maximum gives $27$ linear rows. Section~\ref{sec:cert}
encodes them with three auxiliary variables.

The \texttt{Python} listing of \cite[Appendix~A]{BBLT} transcribes the first
line of \cite[(6.9)]{BBLT} as \texttt{p1+a2+b3+a3}, where the displayed formula
reads $p_1+a_2+b_3+c_3$. We use the displayed formula. It is what
Corollary~\ref{cor:app:geom} gives with the printed second argument
$1-a_2-2b_3-2c_3$. For \eqref{eq:63} we use
\begin{equation}\label{eq:63proof}
 7s-1\le7s-\sum_{i\le d}i\,s_i
 =\sum_{i\le6}(7-i)s_i+\sum_{7\le i\le d}(7-i)s_i
 \le\sum_{i\le6}(7-i)s_i,
\end{equation}
since $7-i\le0$ for $i\ge7$. In our system \eqref{eq:63} is the sum of the two
tail rows of Section~\ref{sec:tail}, so it is redundant.

\section{The multiplicity-at-least-seven tail}\label{sec:tail}

The tail rows we use are
\begin{equation}\label{eq:tailrows}
 \tau_s\ge0,\qquad v_s\ge7\tau_s,\qquad \sum_{i\le6}i\,s_i+v_s\le1
 \qquad(s=a,b,c).
\end{equation}
The first is \eqref{eq:root}. The second: $\log_XS_j\ge0$, with weight
$j\ge7$ in $v_s$ and weight $1$ in $\tau_s$. The third is the side-value
bound in \eqref{eq:boxconstraints}. Adding the last two,
\[
 (v_s-7\tau_s)+(1-\sum_{i\le6}i s_i-v_s)\ge0,
\]
that is, $\sum_{i\le6}i s_i+7\tau_s\le1$; with $\tau_s=s-\sum_{i\le6}s_i$
this is \eqref{eq:63}. With the
upper row $v_s\le d\tau_s$, these rows describe the real relaxation of the
pairs $(\tau_s,v_s)$ exactly. We omit that row, so the system $\mathcal S$ of
Section~\ref{sec:cert} is free of $d$.

\begin{lemma}\label{lem:tail}
Fix a side $s\in\{a,b,c\}$ and let
$M_{\mathrm{tail}}=\prod_{7\le j\le d}\mathfrak s_j^{\,j}$ with
$\mathfrak s_j\sim S_j$. Then
\begin{equation}\label{eq:tailcount}
 \begin{gathered}
  \#\{(\mathfrak s_7,\dots,\mathfrak s_d)\in\N^{d-6}:\mathfrak s_j\sim S_j\}\le\prod_{7\le j\le d}S_j=X^{\tau_s},\\
  M_{\mathrm{tail}}\ge\prod_{7\le j\le d}S_j^{\,j}=X^{v_s},
 \end{gathered}
\end{equation}
and every prime $p$ dividing $M_{\mathrm{tail}}$ does so with
$\operatorname{ord}_p(M_{\mathrm{tail}})\ge7$, where $\operatorname{ord}_p$
denotes the $p$-adic valuation.
\end{lemma}

\begin{proof}
Each interval $[S_j,2S_j)$ contains exactly $S_j$ integers, which gives the
count, and $\mathfrak s_j\ge S_j$ gives the modulus bound. If $p\mid\mathfrak s_j$
with $j\ge7$ then $\operatorname{ord}_p(M_{\mathrm{tail}})\ge j\ge7$.
\end{proof}

The count is exact for $[N,2N)$ and for $(N,2N]$. A closed $[N,2N]$ would
overcount by a factor of up to $2^{d-6}$, attained when all $S_j=1$.

So the tail costs $X^{\tau_s}$ to enumerate. The bound is on roots, not on
$\rad(M_{\mathrm{tail}})$, and uses no squarefreeness. Let $\sigma$ and $T$ be the modulus side and
index set of Section~\ref{sec:cut}. Then $M_{\mathrm{tail}}$ divides the side
value and is at least $X^{v_\sigma}$.
Hence $T$ may contain the tail.

\subsection{Folded roots}\label{sec:fold}
For the rest of the proof fix $0<\varepsilon_0\le4/7$. This is allowed, as
$\varepsilon_0$ is arbitrarily small at the endpoint step.
In \cite[Lemma~2.2]{BBLT} the very high multiplicities are folded into one
index $k$. The folded root need not be squarefree; the remainder goes into
$n_0$. The proof of \cite[Proposition~2.1]{BBLT} uses that lemma with
parameter $\varepsilon_0/4$. So
\[
 k=\lceil4/\varepsilon_0\rceil,\qquad d=6\lceil(4/\varepsilon_0)^{2}\rceil .
\]
Then $k\ge7$, and
$d\ge6(4/\varepsilon_0)^2>4/\varepsilon_0+1>k$. So the folded root lies in the
tail, $M_{\mathrm{tail}}$ still divides the side value, and
$\operatorname{ord}_p(n_k^{\,k})=k\operatorname{ord}_p(n_k)\ge7$ for
$p\mid n_k$. The overlap of $n_0$ and $n_k$ is within one side, so
Lemma~\ref{lem:coprime} is untouched. Unenumerated roots are recovered by
Lemma~\ref{lem:divisor}. The folding costs the factor
$X^{3\varepsilon_0/4}$ in the upper radical bound of \eqref{eq:boxconstraints}.
It also costs the choice of $a_0,b_0,c_0$ in the $X^{\varepsilon_0}$ of
\eqref{eq:reduction}. Both are already charged. It changes neither
Lemma~\ref{lem:tail} nor \eqref{eq:tailrows}.




\section{Three lattice and divisor lemmas}\label{sec:lat}

\begin{lemma}\label{lem:farey}
Let $M\ge1$, let $\xi$ be a unit modulo $M$ and let $U,V\ge1$. Then
\begin{equation}\label{eq:farey}
 \#\{(u,v)\in\N^2:u\le U,\ v\le V,\ \gcd(u,v)=1,\ u\equiv\xi v\ (M)\}
 \le1+\frac{2UV}{M}.
\end{equation}
\end{lemma}

\begin{proof}
Put $\mathcal L_\xi=\{(u,v)\in\Z^2:u\equiv\xi v\ (M)\}$, of index $M$ in
$\Z^2$. Any two of its points have determinant in $M\Z$. Positive collinear
integer points with $\gcd(u,v)=1$ coincide. Thus the points of $\mathcal L_\xi$
with $\gcd(u,v)=1$ in $[1,U]\times[1,V]$ have distinct slopes; order them by
slope. Consecutive
points span with $0$ a triangle of area $\ge M/2$. These triangles have
disjoint interiors and lie in $[0,U]\times[0,V]$. With $P$ points,
$(P-1)M/2\le UV$.
\end{proof}

The proof does not use that $\xi$ is a unit. The constant $2$ is sharp. For
odd $n\ge3$ take $U=n$, $V=n+4$, $M=n^2+2n-4$, $\xi=n-1$. The primitive points
are $(n-1,1)$, $(n,n+4)$, $(1,n+3)$, with consecutive determinants $M$, and
$(P-1)M/(UV)\to2$.

\begin{lemma}\label{lem:qroot}
Let $1\le q\le6$, let $M\ge1$ and let $K$ be a unit modulo $M$. The congruence
$\eta^{q}\equiv K\pmod M$ has at most $8^{\omega(M)}$ solutions in units $\eta$,
hence $X^{o(1)}$ solutions uniformly for $M\le X^{O(1)}$.
\end{lemma}

\begin{proof}
The solution set is empty or a coset of the kernel of $\eta\mapsto\eta^q$.
Odd prime power: the unit group is cyclic, so the kernel has order at most
$q\le6$. Modulus $2^e$, $e\le2$: the unit group has order
at most $2$. Modulus $2^e$, $e\ge3$:
$(\mathbb Z/2^e\mathbb Z)^\times\cong C_2\times C_{2^{e-2}}$, so the kernel has
order $\gcd(q,2)\gcd(q,2^{e-2})\le8$; the $2$-part of $q\le6$ is at most
$4$. Multiply by the Chinese remainder theorem. As
$\omega(M)\ll\log M/\log\log M$ for $M\ge3$,
\[
 8^{\omega(M)}=\exp\bigl(O(\log M/\log\log M)\bigr)=X^{O(1/\log\log X)}
 \qquad(M\le X^{O(1)}).
\]
The form $C^{\omega(M)}\ll_{C,\delta}M^{\delta}$, for fixed $C\ge1$ and $\delta>0$, alone gives no rate. For
$M\le2$ there is one class. No bound $B^{\omega(M)}$ with $B<8$ holds. For
$q=4$, $e\ge4$, the group $C_2\times C_{2^{e-2}}$ has
$\gcd(4,2)\gcd(4,2^{e-2})=8$ fourth roots of unity. And $q=4$ occurs in
\eqref{eq:q}. Local kernel size $8$ occurs only for $2^e$, $e\ge4$, $q=4$. So
for $\omega(M)\ge2$ the product of local kernel sizes is at most
$8\cdot6^{\omega(M)-1}<8^{\omega(M)}$.
\end{proof}

\begin{lemma}\label{lem:divisor}
Write $\operatorname{div}(N)=\#\{m\in\N:m\mid N\}$ for $N\in\N$. For
$r\ge1$ and positive integers $e_1,\dots,e_r$,
\[
 \#\{(x_1,\dots,x_r)\in\N^r:x_1^{e_1}\cdots x_r^{e_r}=N\}
 \le \operatorname{div}(N)^{r-1}\ll_{r,\delta}N^{\delta}\qquad(\delta>0),
\]
uniformly in $e_1,\dots,e_r$. For $N\ge3$ the maximal-order bound gives the
rate
\[
 \operatorname{div}(N)^{r-1}=\exp\!\bigl(O_r(\log N/\log\log N)\bigr),
\]
so for $N\le X^{O_d(1)}$ and $r\le d$ the count is $X^{O_d(1/\log\log X)}$.
\end{lemma}

\begin{proof}
Each $x_i$ divides $N$. The equation determines $x_r$ from
$x_1,\dots,x_{r-1}$, for any $e_i$. Apply the divisor bound and, for the rate,
the maximal order of $\operatorname{div}(N)$.
\end{proof}

Thus an unfactored product of roots determines those roots up to $X^{o(1)}$.

\section{The cut}\label{sec:cut}

Fix a box. Let $(y,z)$ be the other two sides in alphabetical order:
$(b,c)$, $(a,c)$, $(a,b)$ for $\sigma=a,b,c$. Write
$[6]=\{1,2,3,4,5,6\}$. Choose a nonempty
\begin{equation}\label{eq:T}
 \begin{gathered}
  T\subseteq\{1,2,3,4,5,6,\mathrm{tail}\},\\
  \kappa_T=\sum_{t\in T\cap[6]}\sigma_t+\mathbf 1_{\mathrm{tail}\in T}\tau_\sigma,
  \qquad
  \Omega_T=\sum_{t\in T\cap[6]}t\,\sigma_t+\mathbf 1_{\mathrm{tail}\in T}v_\sigma .
 \end{gathered}
\end{equation}
Here $\kappa_T$ is the selected root mass for this section. The symbols
$u,v,M,\xi$ below are local to the cut and are not the corresponding symbols
in Sections~\ref{sec:pow3} and~\ref{sec:moment}; $v$ is also distinct from
the tail coordinate $v_\sigma$.
Choose free sets $R_y,R_z\subseteq[6]$, not both empty, and put
\begin{equation}\label{eq:q}
 q=\gcd(\gcd R_y,\gcd R_z),
 \qquad
 F=\sum_{i\in R_y}i\,y_i+\sum_{j\in R_z}j\,z_j,
\end{equation}
with $\gcd\emptyset=0$. For such a choice
$(\sigma,T,R_y,R_z)$, the quantities in \eqref{eq:T}--\eqref{eq:q}
are fixed by the box and the index sets.
As $R_y\cup R_z\ne\emptyset$, $1\le q\le6$ and $q$ divides every element of
$R_y\cup R_z$. By the choice of $d$ in Section~\ref{sec:fold},
$d=6\lceil(4/\varepsilon_0)^2\rceil\ge294$,
so the tail notation applies. Write $\pos{t}=\max\{t,0\}$.

\begin{proposition}\label{prop:cut}
Let a box satisfy the constraints after \eqref{eq:reduction}, with
$\mathcal B_d\ge1$. Then
\begin{equation}\label{eq:cut}
 D\le\Bigl(y-\sum_{i\in R_y}y_i\Bigr)+\Bigl(z-\sum_{j\in R_z}z_j\Bigr)
  +\kappa_T+\pos{\frac Fq-\Omega_T}+o_{\varepsilon_0,d}(1),
\end{equation}
uniformly over the finitely many choices of $(\sigma,T,R_y,R_z)$.
\end{proposition}

For later use write
\[
  \mathsf A_{\rm cut}=\Bigl(y-\sum_{i\in R_y}y_i\Bigr)+\Bigl(z-\sum_{j\in R_z}z_j\Bigr)+\kappa_T,
  \qquad \mathsf B_{\rm cut}=F/q-\Omega_T.
\]
The limiting form of \eqref{eq:cut} is
$D\le\mathsf A_{\rm cut}+\max(0,\mathsf B_{\rm cut})$.

\begin{proof}
\emph{Sizes.} The box fixes $n_{a,0}=a_0$, $n_{b,0}=b_0$ and $n_{c,0}=c_0$.
Write $n_{\sigma,0},n_{y,0},n_{z,0}$ for these coefficients and
$S_{\sigma,i},S_{y,i},S_{z,i}$ for the dyadic root sizes of the $\sigma,y,z$ sides.
Besides \eqref{eq:box}, the proof uses
the dyadic ranges and $\prod_iA_i^i$, $\prod_iB_i^i$, $\prod_iC_i^i\le X$. So all
products and cofactors are at most $X^{O_d(1)}$, and
Lemmas~\ref{lem:qroot} and~\ref{lem:divisor} cost $X^{o(1)}$. The
coefficients enter no size estimate.

\emph{Enumeration.} Write $\mathfrak m_i,\mathfrak y_i,\mathfrak z_i$ for the
roots of the $\sigma$-, $y$- and $z$-sides, so
$\mathfrak m_i\sim S_{\sigma,i}$, $\mathfrak y_i\sim S_{y,i}$,
$\mathfrak z_i\sim S_{z,i}$. Enumerate: the roots of
$y$ outside $R_y$; of $z$ outside $R_z$; the selected roots of $\sigma$. Tails
go through Lemma~\ref{lem:tail}. The $y$- and $z$-tails are never in
$R_y,R_z\subseteq[6]$, so they are enumerated too. The cost is
\begin{equation}\label{eq:enum}
 X^{(y-\sum_{i\in R_y}y_i)+(z-\sum_{j\in R_z}z_j)+\kappa_T+o_{\varepsilon_0,d}(1)} .
\end{equation}

\emph{Modulus.} The tail roots of $\sigma$ are $\mathfrak m_j$, $j\ge7$. Put
\begin{equation}\label{eq:M}
  M=\prod_{t\in T\cap[6]}\mathfrak m_t^{\,t}\,
  M_{\mathrm{tail}}^{\mathbf 1_{\mathrm{tail}\in T}},
 \qquad M_{\mathrm{tail}}=\prod_{7\le j\le d}\mathfrak m_j^{\,j}.
\end{equation}
Then $M$ divides the $\sigma$-side value. Also
$\mathfrak m_t\ge S_{\sigma,t}=X^{\sigma_t}$
(left-closed intervals), and Lemma~\ref{lem:tail} bounds the tail. So
\begin{equation}\label{eq:Msize}
 M\ge X^{\Omega_T} .
\end{equation}
Put
\begin{equation}\label{eq:uv}
 u=\prod_{i\in R_y}\mathfrak y_i^{\,i/q},
 \qquad
 v=\prod_{j\in R_z}\mathfrak z_j^{\,j/q},
\end{equation}
an empty product being $1$; the exponents are integers by \eqref{eq:q}.

\emph{Congruence.} By Lemma~\ref{lem:coprime}, every factor of the
$y$- and $z$-sides is a unit modulo $M$. This covers the coefficients, the
enumerated roots, $u$ and $v$. The equation of
\eqref{eq:box} modulo $M$ gives
\begin{equation}\label{eq:congr}
 (uv^{-1})^{q}\equiv K\pmod M
\end{equation}
with
\begin{equation}\label{eq:Ksign}
 K\equiv\epsilon_{\rm cut} E_zE_y^{-1}\pmod M,\qquad
 \epsilon_{\rm cut}=\begin{cases}-1,&\sigma=c,\\+1,&\sigma\ne c,\end{cases}
\end{equation}
where
\[
\begin{aligned}
 E_y&=n_{y,0}\prod_{i\in[6]\setminus R_y}\mathfrak y_i^{\,i}
 \prod_{7\le j\le d}\mathfrak y_j^{\,j},\\
 E_z&=n_{z,0}\prod_{j\in[6]\setminus R_z}\mathfrak z_j^{\,j}
 \prod_{7\le j\le d}\mathfrak z_j^{\,j}.
\end{aligned}
\]
So $K$ is a unit fixed by the enumerated data.

\emph{Count.} By Lemma~\ref{lem:qroot}, $\xi=uv^{-1}$ lies in one of
$X^{o(1)}$ classes. By Lemma~\ref{lem:coprime}, $\gcd(u,v)=1$. Also
\[
 u\le2^{\sum_{i\in R_y}i/q}\prod_{i\in R_y}S_{y,i}^{\,i/q}=:U_{\rm cut},
 \qquad
 v\le2^{\sum_{j\in R_z}j/q}\prod_{j\in R_z}S_{z,j}^{\,j/q}=:V_{\rm cut},
\]
with $U_{\rm cut},V_{\rm cut}\ge1$; the powers of $2$ cost $O_d(1/\log X)$, as
$\sum_{r\le6}r=21$. Hence
\[
 U_{\rm cut}=X^{\sum_{i\in R_y}i\,y_i/q+O_d(1/\log X)},
 \qquad
 V_{\rm cut}=X^{\sum_{j\in R_z}j\,z_j/q+O_d(1/\log X)},
\]
so $U_{\rm cut}V_{\rm cut}=X^{F/q+O_d(1/\log X)}$.
Lemma~\ref{lem:farey} and \eqref{eq:Msize} bound the pairs $(u,v)$ in one
class by
\begin{equation}\label{eq:pairs}
 1+\frac{2U_{\rm cut}V_{\rm cut}}{M}\ll_dX^{\pos{F/q-\Omega_T}+O_d(1/\log X)} .
\end{equation}
\emph{Recovery.} Lemma~\ref{lem:divisor} recovers the free roots from $u,v$ up
to $X^{o(1)}$. Then the $y$- and $z$-side values are known, and the equation fixes
the $\sigma$-side value. Divide by $n_{\sigma,0}$ and the selected roots:
\[
 Q_\sigma=\prod_{i\in[6]\setminus T}\mathfrak m_i^{\,i}\,
 \left(\prod_{7\le j\le d}\mathfrak m_j^{\,j}\right)^{1-\mathbf 1_{\mathrm{tail}\in T}}.
\]
Lemma~\ref{lem:divisor} recovers these unselected roots up to $X^{o(1)}$.
Multiplying \eqref{eq:enum} by $X^{o(1)}$ and \eqref{eq:pairs} gives
\eqref{eq:cut}.
\end{proof}

If one free set is empty, $U_{\rm cut}=1$ or $V_{\rm cut}=1$, and Lemma~\ref{lem:farey} still
applies. The exponent $q$ depends only on the index sets $R_y,R_z$, not on
which roots equal $1$. Counting coefficients is charged to \eqref{eq:reduction}:
inside a box $a_0,b_0,c_0$ are fixed.

Count: $64$ subsets of $[6]$ per counted side, $127$ sets $T$, three sides
$\sigma$. Without the pair of empty free sets this is
$3(64^2-1)127=1\,560\,195$, as in \eqref{eq:cutcount}.


\section{Archimedean rows}\label{sec:arch}

Throughout Sections~\ref{sec:arch}--\ref{sec:spacing}, assume
$0<\mu\le1$. Take $\varepsilon_0\le4/7$ and $d=d(\varepsilon_0)$ as fixed
in Section~\ref{sec:fold}. Fix a box admissible in \eqref{eq:reduction} at
scale $Y$, with $\lambda$ replaced by $\mu$ and $\mathcal B_d\ge1$. Of the constraints \eqref{eq:boxconstraints} only the side-value bounds are
used below, with the coefficient bound $1\le n_{s,0}\le Y^{\varepsilon_0}$. We
also use the equation, the gcd condition and the windows of \eqref{eq:box}. For the
finite estimates take $Y\ge Y_0$, the threshold of Lemma~\ref{lem:endpoint};
it depends only on $\varepsilon_0,d$, and $Y_0\ge\exp(\exp(1))$. All coordinates
$s,\tau_s,v_s,L,p_i$ from Section~\ref{sec:red} are now read at base $Y$.
The integer side values are denoted by $n_s$; $V_s$ is the exponent
coordinate of the root product, not of $n_s$. The root mass of side $s$ is
$s=\sum_{i\le d}s_i$.
Put $s_i=\log_YS_{s,i}$ and
$D=\log_Y\mathcal B_d$.
\begin{equation}\label{eq:V}
 V_s=\log_Y\prod_{j\le d}S_{s,j}^{\,j}=\sum_{i\le6}i\,s_i+v_s\qquad(s=a,b,c).
\end{equation}
For a tuple of the box write
$n_s=n_{s,0}\prod_{j\le d}\mathfrak s_j^{\,j}$ for its side values, so
$n_a+n_b=n_c$. Here $n_{s,0}$ is the coefficient $a_0,b_0$ or $c_0$ on
side $s$. Since $S_j\le\mathfrak s_j<2S_j$ and $1\le n_{s,0}\le Y^{\varepsilon_0}$,
\begin{equation}\label{eq:nsize}
 Y^{V_s}\le n_s<2^{h_d}\,n_{s,0}\,Y^{V_s}\le2^{h_d}\,Y^{V_s+\varepsilon_0}.
\end{equation}
The side-value constraints of \eqref{eq:boxconstraints} read
\begin{equation}\label{eq:Vbox}
 V_a\le1,\qquad V_b\le1,\qquad 1-\frac{\varepsilon_0}2\le V_c\le1 .
\end{equation}

\begin{lemma}\label{lem:archorder}
\[
 V_a\le V_c+\frac{\varepsilon_0}2,\qquad V_b\le V_c+\frac{\varepsilon_0}2 .
\]
\end{lemma}

\begin{proof}
By \eqref{eq:Vbox}, $V_a\le1\le V_c+\varepsilon_0/2$. Likewise for $V_b$.
\end{proof}

\begin{lemma}\label{lem:archcover}
At least one of
\[
 V_c\le V_a+\frac{3\varepsilon_0}2+\log_Y2^{1+h_d},
 \qquad
 V_c\le V_b+\frac{3\varepsilon_0}2+\log_Y2^{1+h_d}
\]
holds.
\end{lemma}

\begin{proof}
As $\mathcal B_d\ge1$, choose one tuple of the box. Then $n_a,n_b\ge1$ and
$n_a+n_b=n_c$, so $\max(n_a,n_b)\ge n_c/2$. Suppose $n_a\ge n_c/2$. By
\eqref{eq:Vbox} and \eqref{eq:nsize},
\[
 Y^{1-\varepsilon_0/2}\le Y^{V_c}\le n_c\le2n_a
 <2^{1+h_d}\,Y^{V_a+\varepsilon_0}.
\]
Take $\log_Y$:
\[
 1\le V_a+\frac{3\varepsilon_0}2+\log_Y2^{1+h_d}.
\]
With $V_c\le1$ this is the first inequality. If $n_b\ge n_c/2$, the same
argument gives the second.
\end{proof}

The choice in Lemma~\ref{lem:archcover} uses one tuple. The conclusion
concerns the dyadic endpoints, so it holds for the whole box.

\begin{corollary}\label{cor:arch}
Put $\rho_{\mathrm{arch}}=3\varepsilon_0/2+\log_Y2^{1+h_d}$. Every nonempty
admissible box satisfies
\begin{equation}\label{eq:ARCHorder}
 V_a\le V_c+\rho_{\mathrm{arch}},\qquad V_b\le V_c+\rho_{\mathrm{arch}},
\end{equation}
and at least one arm of
\begin{equation}\label{eq:ARCH}
 V_c\le V_a+\rho_{\mathrm{arch}}\quad\text{or}\quad V_c\le V_b+\rho_{\mathrm{arch}} .
\end{equation}
\end{corollary}

\begin{proof}
Lemmas~\ref{lem:archorder} and~\ref{lem:archcover}, and
$\varepsilon_0/2\le\rho_{\mathrm{arch}}$.
\end{proof}

In the limit $\rho_{\mathrm{arch}}\to0$, the first arm of \eqref{eq:ARCH} with
\eqref{eq:ARCHorder} gives $V_a=V_c\ge V_b$, the second $V_b=V_c\ge V_a$.
A cover node of the certificate carries both arms.

\section{The cubic row}\label{sec:hb3}

We use the following theorem of Heath-Brown \cite[Theorem~1.3]{HBcube}, with
its counting function $\rho$ renamed $\mathcal R$. In this theorem $B$ is a
height and $a,b,c$ are coefficients.

\begin{theoremA}[{\cite[Theorem~1.3]{HBcube}}]\label{thm:HBcubic}
Let $\mathcal R(B;a,b,c)$ be the number of integer solutions $x,y,z$ to the equation
\[
 ax^3+by^3+cz^3=0
\]
with $\gcd(x,y,z)=1$ and $\max(|x|,|y|,|z|)\le B$. Then
\[\mathcal R(B;a,b,c)\ll_\varepsilon B^\varepsilon\]
for any $\varepsilon>0$, uniformly for all non-zero integer $a,b,c$.
\end{theoremA}

The implied constant depends on $\varepsilon$ alone. The coefficients need
not be coprime or positive.

\begin{proposition}\label{prop:HB3}
For every $\varepsilon_1>0$,
\begin{equation}\label{eq:HB3}
 D\le\sum_{s=a,b,c}(s-s_3)+\varepsilon_1+O_{\varepsilon_1}\Bigl(\frac1{\log Y}\Bigr)
 =L-p_3+\varepsilon_1+O_{\varepsilon_1}\Bigl(\frac1{\log Y}\Bigr).
\end{equation}
The implied constant depends only on $\varepsilon_1$.
\end{proposition}

\begin{proof}
\emph{Enumeration.} Send a tuple of the box to its roots
$(\mathfrak a_i,\mathfrak b_i,\mathfrak c_i)_{i\ne3}$ and to
$(x_a,x_b,x_c)=(\mathfrak a_3,\mathfrak b_3,\mathfrak c_3)$. This map is injective.
As $\#\{\mathfrak s:\mathfrak s\sim S\}=S$, the first component takes at most
\[
 \prod_{s=a,b,c}\ \prod_{i\le d,\ i\ne3}S_{s,i}=Y^{\sum_s(s-s_3)}
\]
values. The tail roots $i\ge7$ are enumerated one by one, as in
Lemma~\ref{lem:tail}.

\emph{Equation.} Fix the first component. Put
\[
 \alpha=a_0\prod_{i\ne3}\mathfrak a_i^{\,i},\qquad
 \beta=b_0\prod_{i\ne3}\mathfrak b_i^{\,i},\qquad
 \gamma=c_0\prod_{i\ne3}\mathfrak c_i^{\,i}.
\]
These are positive integers. The equation of \eqref{eq:box} becomes
\[
\alpha x_a^3+\beta x_b^3+(-\gamma)x_c^3=0 .
\]
\emph{Primitivity.} $x_a\mid n_a$, $x_b\mid n_b$, $x_c\mid n_c$, and the side values
are pairwise coprime by Lemma~\ref{lem:coprime}. Hence
$\gcd(x_a,x_b,x_c)=1$.

\emph{Height.} $1\le x_a<2A_3\le2\bigl(\prod_jA_j^{\,j}\bigr)^{1/3}\le2Y$ by
\eqref{eq:Vbox}. Likewise $x_b,x_c<2Y$.

\emph{Count.} By Theorem~\ref{thm:HBcubic} with $B=2Y$ and coefficients
$(\alpha,\beta,-\gamma)$, the number of $(x_a,x_b,x_c)$ is at most
$\mathcal R(2Y;\alpha,\beta,-\gamma)\le C(\varepsilon_1)(2Y)^{\varepsilon_1}$,
uniformly in $\alpha,\beta,\gamma$. Hence
\[
 \mathcal B_d\le C(\varepsilon_1)\,2^{\varepsilon_1}\,Y^{\sum_s(s-s_3)+\varepsilon_1}.
\]
Take $\log_Y$.
\end{proof}

The bound is uniform in the coefficients $a_0,b_0,c_0$ and in the enumerated
roots. A constant depending on $\alpha,\beta,\gamma$ would not give a row.

\section{Spacing rows}\label{sec:spacing}

Choose $e\in\{a,b,c\}$, either ordering $(y,z)$ of the other two sides, and
$2\le\ell\le6$. Put
\begin{equation}\label{eq:SPAB}
 \mathsf A_{\rm sp}=y+(z-z_\ell),\qquad
 \mathsf B_{\rm sp}=V_e-V_z+z_\ell .
\end{equation}
Here $y$ and $z$ also denote the root masses of these sides, as in
\eqref{eq:coord}.

For side $e$, write $n_{e,0}$ for its coefficient and $\mathfrak e_j$ for
its roots (Section~\ref{sec:arch}). Choose $K_d$ so that
\[
 \operatorname{div}(N)^{d-1}\le Y^{K_d/\log\log Y}
 \quad\text{if }Y\ge\exp(\exp(1)),\quad N\le2^{h_d}Y;
\]
such a constant exists by the divisor bound.

\begin{proposition}\label{prop:spacing}
\begin{equation}\label{eq:SP}
 D\le\mathsf A_{\rm sp}+\pos{\mathsf B_{\rm sp}}+\varepsilon_0
 +O_d\Bigl(\frac1{\log\log Y}\Bigr).
\end{equation}
\end{proposition}

\begin{proof}
\emph{Signs.} From $n_a+n_b=n_c$,
\[
 n_e=\varsigma_yn_y+\varsigma_zn_z,\qquad
 (\varsigma_y,\varsigma_z)=
 \begin{cases}
  (1,1),&e=c,\\
  (-1,1),&e\ne c,\ z=c,\\
  (1,-1),&e\ne c,\ y=c.
 \end{cases}
 \]
This ordering need not be the alphabetical order used in Section~\ref{sec:cut}.

\emph{Enumeration.} Enumerate all roots $\mathfrak y_j$ $(j\le d)$ of $y$ and
all roots $\mathfrak z_j$ $(j\le d,\ j\ne\ell)$ of $z$. As
$\#\{\mathfrak s:\mathfrak s\sim S\}=S$, there are
\begin{equation}\label{eq:SPenum}
 \prod_{j\le d}S_{y,j}\prod_{j\ne\ell}S_{z,j}=Y^{y+z-z_\ell}
\end{equation}
choices. Fix one. Then $n_y$ is fixed, and
\[
 n_z=\Pi_z\,x^\ell,\qquad \Pi_z=n_{z,0}\prod_{j\ne\ell}\mathfrak z_j^{\,j},\qquad
 x=\mathfrak z_\ell\in[S_{z,\ell},2S_{z,\ell}).
\]

\emph{Interval.} By \eqref{eq:nsize}, $n_e$ lies in an interval $I_e$ of length
$|I_e|<2^{h_d}Y^{V_e+\varepsilon_0}$. Then
\[
 x^\ell=\varsigma_z\,\frac{n_e-\varsigma_yn_y}{\Pi_z}\in J,
 \qquad J=\varsigma_z\,\frac{I_e-\varsigma_yn_y}{\Pi_z},
\]
an interval of length $|I_e|/\Pi_z$. As
$\Pi_z\ge\prod_{j\ne\ell}S_{z,j}^{\,j}=Y^{V_z-\ell z_\ell}$,
\[
 |J|<2^{h_d}\,Y^{V_e-V_z+\ell z_\ell+\varepsilon_0}.
\]

\emph{Spacing.} Let $\mathcal X$ be the integers
$x\in[S_{z,\ell},2S_{z,\ell})$ for which $x^\ell\in J$, and put
$m=|\mathcal X|$. If $m=0$, this enumeration choice contributes nothing.
Otherwise write $x_1<\dots<x_m$ for $\mathcal X$. For $i\ge2$,
\[
 x_i^\ell-x_{i-1}^\ell\ge \ell\,x_{i-1}^{\ell-1}\ge S_{z,\ell}^{\ell-1}=Y^{(\ell-1)z_\ell}.
\]
So $(m-1)Y^{(\ell-1)z_\ell}\le x_m^\ell-x_1^\ell<|J|$, and
\begin{equation}\label{eq:SPcount}
 m<1+2^{h_d}Y^{\mathsf B_{\rm sp}+\varepsilon_0}
 \le2^{1+h_d}\,Y^{\pos{\mathsf B_{\rm sp}}+\varepsilon_0}.
\end{equation}
This includes $S_{z,\ell}=1$, where $x=1$ and $m\le1$.

\emph{Recovery.} Given $x$, $n_z$ and then $n_e=\varsigma_yn_y+\varsigma_zn_z$ are
fixed. Put $N=n_e/n_{e,0}=\prod_{j\le d}\mathfrak e_j^{\,j}<2^{h_d}Y$. By
Lemma~\ref{lem:divisor} the roots of $e$ take at most
$Y^{K_d/\log\log Y}$ values.

\emph{Total.} Multiply \eqref{eq:SPenum}, \eqref{eq:SPcount} and the recovery
factor:
\[
 \mathcal B_d\le2^{1+h_d}\,
 Y^{\mathsf A_{\rm sp}+\pos{\mathsf B_{\rm sp}}+\varepsilon_0+O_d(1/\log\log Y)} .
\]
Take $\log_Y$; $2^{1+h_d}$ costs $O_d(1/\log Y)$.
\end{proof}

\begin{corollary}\label{cor:spacing}
Put
\[
 \rho_{\mathrm{sp}}=\varepsilon_0+\frac{K_d}{\log\log Y}+\log_Y2^{1+h_d}.
\]
This is the error in \eqref{eq:SP}. For each of the
$3\cdot2\cdot5=30$ choices of $(e,y,z,\ell)$, a
nonempty admissible box satisfies at least one arm of
\begin{equation}\label{eq:SPor}
 D\le\mathsf A_{\rm sp}+\rho_{\mathrm{sp}}\quad\text{or}\quad
 D\le\mathsf A_{\rm sp}+\mathsf B_{\rm sp}+\rho_{\mathrm{sp}} .
\end{equation}
\end{corollary}

\begin{proof}
$\mathsf A_{\rm sp}+\pos{\mathsf B_{\rm sp}}
=\max(\mathsf A_{\rm sp},\mathsf A_{\rm sp}+\mathsf B_{\rm sp})$.
\end{proof}

All these errors are $\mu$-free. With $\varepsilon_1=\varepsilon_0$, each is at
most $2\varepsilon_0$ for $Y\ge Y_0$. Table~\ref{tab:rho}
lists them with the others.

\section{Cubic-sieve rows}\label{sec:pow3}

\subsection{Setting}\label{sec:pow3:setting}

Fix $\varepsilon_0\le4/7$ and $d=d(\varepsilon_0)$ as in
Section~\ref{sec:fold}. Assume
$Y\ge Y_0$, the threshold of Lemma~\ref{lem:endpoint}. Fix parameters for a nonempty
box counted by $\mathcal B_d$ in \eqref{eq:reduction} at scale $Y$, and put
$D=\log_Y\mathcal B_d$. For each side $s\in\{a,b,c\}$, write $n_{s,0}$ for its
 coefficient and $\mathfrak r_{s,j}\sim S_{s,j}$ $(1\le j\le d)$ for its roots.
The symbols $M,x,u,v,U,V,\xi$ below are local to this section. Its $M$
is the canonical coefficient of \eqref{eq:pw:Mx}. The
modulus $Q$ is local to each analytic-input lemma and to Steps~5--7 of
Proposition~\ref{prop:pw:master}. Here
$v_s$ is the tail coordinate, while $v$ is a represented-side label. The
integer $V$ (as in $V_{\mathfrak o}$, $V(\mathbf u)$, $V_\iota$) is distinct
from the coordinate $V_s$.
The side values
\[
 n_s=n_{s,0}\prod_{j\le d}\mathfrak r_{s,j}^{\,j}
\]
 satisfy $n_a+n_b=n_c$ and are pairwise coprime by Lemma~\ref{lem:coprime}. The
coordinates are
\begin{equation}\label{eq:pw:coord}
 \begin{gathered}
 s_{s,j}=\log_YS_{s,j},\qquad
 r_s=\sum_{j\le d}s_{s,j},\qquad
 t_s=\sum_{7\le j\le d}s_{s,j},\\
 v_s=\sum_{7\le j\le d}j\,s_{s,j},\qquad
 V_s=\sum_{j\le6}j\,s_{s,j}+v_s .
 \end{gathered}
\end{equation}
In the notation of Section~\ref{sec:red}, $r_s=s$ and $t_s=\tau_s$; the
coordinate $s_{s,j}$ is the earlier $s_j$. By
\eqref{eq:boxconstraints}, $V_s\le1$. Hence $n_s\le n_{s,0}\prod_j(2S_{s,j})^j\le
2^{d^2}Y^{1+\varepsilon_0}\le Y^2$ for large $Y$.

Fix a \emph{powered side} $\sigma$. Its two \emph{represented sides} $u,v$ and
the sign $\epsilon_{\rm pw}$ are
\begin{equation}\label{eq:pw:orient}
 (\sigma;u,v;\epsilon_{\rm pw})\in\{(a;b,c;-1),\ (b;a,c;-1),\ (c;a,b;+1)\},
 \qquad n_v+\epsilon_{\rm pw}\,n_u=n_\sigma>0 .
\end{equation}
So $\sigma\in\{a,b\}$ uses the positive difference $n_c-n_u$. The side
$\sigma=c$ uses the sum $n_a+n_b$. Put
\begin{equation}\label{eq:pw:Mx}
 M=n_{\sigma,0}\,\mathfrak r_{\sigma,1}\mathfrak r_{\sigma,2}^{\,2}
   \mathfrak r_{\sigma,4}\mathfrak r_{\sigma,5}^{\,2}
   \prod_{7\le j\le d}\mathfrak r_{\sigma,j}^{\,j},
 \qquad
 x=\mathfrak r_{\sigma,3}\mathfrak r_{\sigma,4}\mathfrak r_{\sigma,5}
   \mathfrak r_{\sigma,6}^{\,2}.
\end{equation}
Then $n_\sigma=Mx^3$. We call $M$ the canonical coefficient. Its labels are
the tuple
\[
 (\mathfrak r_{\sigma,j})_{j\in\{1,2,4,5,7,\dots,d\}}.
\]
Put
\begin{equation}\label{eq:pw:CWR}
 \mathcal C_\sigma=r_\sigma-s_{\sigma,3}-s_{\sigma,6},\qquad
 W_\sigma=v_\sigma+s_{\sigma,1}+2s_{\sigma,2}+s_{\sigma,4}+2s_{\sigma,5},\qquad
 R_\sigma=r_u+r_v .
\end{equation}
Here $R_\sigma$ is the combined root mass of the represented sides.
Put $\zeta_\sigma=\log_Y n_{\sigma,0}$. The coefficient range in
\eqref{eq:reduction} gives $0\le\zeta_\sigma\le\varepsilon_0$. Since each
labelled root lies in $[S_{\sigma,j},2S_{\sigma,j})$, the finite-box modulus
satisfies
\begin{equation}\label{eq:pw:Mscale}
 W_\sigma+\zeta_\sigma\le\log_Y M
 <W_\sigma+\zeta_\sigma+
 \frac{\bigl(6+\sum_{7\le j\le d}j\bigr)\log2}{\log Y}
 =W_\sigma+\zeta_\sigma+O_d\!\left(\frac1{\log Y}\right).
\end{equation}

\begin{lemma}\label{lem:pw:ledger}
\begin{enumerate}
\item[(i)] There are at most $Y^{\mathcal C_\sigma}$ label tuples.
\item[(ii)] Every label tuple gives $M\ge Y^{W_\sigma}$.
\item[(iii)] $W_\sigma-\mathcal C_\sigma=s_{\sigma,2}+s_{\sigma,5}+(v_\sigma-t_\sigma)\ge0$.
\item[(iv)] A given integer $M\le Y^2$ arises from $Y^{o_d(1)}$ label tuples.
\end{enumerate}
\end{lemma}

\begin{proof}
(i) $\#\{\mathfrak r:\mathfrak r\sim S\}=S$, and
$\prod_{j\ne3,6}S_{\sigma,j}=Y^{r_\sigma-s_{\sigma,3}-s_{\sigma,6}}$.
(ii) $n_{\sigma,0}\ge1$ and $\mathfrak r_{\sigma,j}\ge S_{\sigma,j}$, so
$M\ge Y^{W_\sigma}$. Also
\[
 M\le n_{\sigma,0}\prod_{j\le d}\mathfrak r_{\sigma,j}^{\,j}
 <2^{h_d}Y^{V_\sigma+\varepsilon_0}\le Y^2
\]
for large $Y$, since $V_\sigma\le1$ and $\varepsilon_0<1$.
(iii) Expand \eqref{eq:pw:CWR}; $v_\sigma\ge7t_\sigma$ by
\eqref{eq:pw:coord}. (iv) $n_{\sigma,0}$ is fixed; apply
Lemma~\ref{lem:divisor} to $M/n_{\sigma,0}$ with $r=d-2$.
\end{proof}

All $o(1)$ below tend to $0$ as $Y\to\infty$, for fixed $\varepsilon_0$, $d$
and the parameter $\theta$ of the sieve. They are uniform in the box.

\subsection{Analytic inputs}\label{sec:pow3:analytic}

Fix once and for all $w\in C^\infty(\mathbb R)$ with $0\le w\le1$,
$\operatorname{supp}w\subset[1/2,4]$ and $w=1$ on $[1,2]$. For a Dirichlet
character $\chi$, extended by $0$ on nonunits, and $N\ge1$ put
\begin{equation}\label{eq:pw:F}
 F_N(\chi)=\sum_{n\ge1}\chi(n)\,w(n/N).
\end{equation}

We quote the all-character estimate of \cite[Lemma~4]{GMRR}. In this section
$Q$ denotes the modulus of an analytic input, or the modulus
$M\mathfrak p_1\mathfrak p_2$ of Steps~5--7 of
Proposition~\ref{prop:pw:master}. It is local to that input or step. It differs
from the cut cofactor $Q_\sigma$ and the geometry forms $Q_j$.
The $L(s,\chi)$ below is the Dirichlet $L$-function, not the root mass $L$;
$s$ in this subsection is a complex variable.

\begin{lemma}[{\cite[Lemma~4]{GMRR}, ``Hybrid fourth moment estimate''}]\label{lem:pw:gmrr}
Let $H\ge2$ and $q_0\ge2$. Then
\[
 \sum_\chi\int_{|t|\le H}\bigl|L(\tfrac12+it,\chi)\bigr|^4dt\ll H\varphi(q_0)\log(Hq_0)^4,
\]
 where the sum is over all characters modulo $q_0$.
\end{lemma}

The all-character estimate in Lemma~\ref{lem:pw:gmrr} is the only input needed
for the next lemma. Since $\varphi(Q)\le Q$, it implies the weaker form used below.

\begin{lemma}\label{lem:pw:conductor}
For $H\ge2$ and $Q\ge2$,
\[
 \sum_{\chi\bmod Q}\int_{|t|\le H}
 |L(\tfrac12+it,\chi)|^4\,dt
 \ll HQ\log(HQ)^4.
\]
\end{lemma}

\begin{proof}
Apply Lemma~\ref{lem:pw:gmrr} with $q_0=Q\ge2$ and use
$\varphi(Q)\le Q$. In our applications $Q\le Y^{O_{d,\theta}(1)}$. The factor
$Q/\varphi(Q)$ lost here is paid in Step~7 of
Proposition~\ref{prop:pw:master}.
\end{proof}

\begin{lemma}\label{lem:pw:S4}
For all integers $Q\ge2$ and reals $N\ge1$,
\begin{equation}\label{eq:pw:S4}
 \sum_{\substack{\chi\bmod Q\\ \chi\ne\chi_0}}|F_N(\chi)|^4
 \ll Q\,N^2\log(2Q)^4 .
\end{equation}
The implied constant depends only on $w$.
\end{lemma}

\begin{proof}
Let $\widetilde w(s)=\int_0^\infty w(y)y^{s-1}dy$. It is entire. Four
integrations by parts give, for $\tfrac12\le\Re s\le2$,
\[
 \widetilde w(s)=\frac{1}{s(s+1)(s+2)(s+3)}\int_0^\infty w^{(4)}(y)\,y^{s+3}dy
 \ll(1+|s|)^{-4}.
\]
By Mellin inversion,
\[
 F_N(\chi)=\frac1{2\pi i}\int_{(2)}L(s,\chi)\,\widetilde w(s)\,N^s\,ds .
\]
Let $\chi\ne\chi_0$. Its partial sums $A_\chi(y)=\sum_{n\le y}\chi(n)$ satisfy
$|A_\chi(y)|\le\varphi(Q)$. Partial summation gives
$L(s,\chi)=s\int_1^\infty A_\chi(y)y^{-s-1}dy$. Thus $L(s,\chi)$ is
holomorphic in $\Re s>0$, and $|L(s,\chi)|\le2Q(1+|s|)$ for
$\Re s\ge\frac12$. The
integrand is therefore $\ll QN^2(1+|s|)^{-3}$ on $\tfrac12\le\Re s\le2$.
Move the line to $\Re s=\frac12$; there is no pole. With
$k_w(t)=|\widetilde w(\tfrac12+it)|\ll(1+|t|)^{-4}$,
\[
 |F_N(\chi)|\le\frac{N^{1/2}}{2\pi}\int_{\mathbb R}
 \bigl|L(\tfrac12+it,\chi)\bigr|\,k_w(t)\,dt .
\]
Hölder with respect to $k_w(t)\,dt$ gives
\[
 \begin{aligned}
 |F_N(\chi)|^4&\le\frac{N^2}{(2\pi)^4}\Bigl(\int k_w\Bigr)^3
 \int_{\mathbb R}\bigl|L(\tfrac12+it,\chi)\bigr|^4k_w(t)\,dt\\
 &\ll N^2\int_{\mathbb R}\bigl|L(\tfrac12+it,\chi)\bigr|^4(1+|t|)^{-4}dt .
 \end{aligned}
\]
Sum over $\chi\ne\chi_0$ and enlarge to all $\chi\bmod Q$. Split
$|t|\le2$ and $2^{j-1}<|t|\le2^j$ $(j\ge2)$. By
Lemma~\ref{lem:pw:conductor} with modulus $Q$ and $H=2^j$,
\[
 \sum_{\chi\ne\chi_0}|F_N(\chi)|^4
 \ll N^2\sum_{j\ge1}2^{-4(j-1)}\,2^jQ\log(2^jQ)^4
 \ll N^2Q\log(2Q)^4\sum_{j\ge1}2^{4-3j}j^4 .
\]
The last sum converges.
\end{proof}

\begin{lemma}\label{lem:pw:G4}
For $N',N''\ge1$ and a Dirichlet character $\chi$, put
\[
 G(\chi)=F_{N'}(\chi)F_{N''}(\chi^2),\qquad
 Z_0=4^6(N'N''^{\,2})^2.
\]
For every integer $Q\ge1$ and $\delta>0$,
\begin{equation}\label{eq:pw:G4}
 \sum_{\chi\bmod Q}|G(\chi)|^4\ll_\delta\varphi(Q)\,(N'N'')^2\,Z_0^{\delta}
 \Bigl(1+\frac {Z_0}Q\Bigr).
\end{equation}
The sum includes the principal character.
\end{lemma}

\begin{proof}
Complete multiplicativity gives
\[
 G(\chi)=\sum_{n\ge1}A(n)\chi(n),\qquad
 A(n)=\sum_{n_1n_2^2=n}w(n_1/N')\,w(n_2/N'')\ge0 .
\]
The coefficients do not depend on $\chi$. Then
$G(\chi)^2=\sum_mB(m)\chi(m)$ with $B=A*A\ge0$. Now
$B(m)=0$ unless $m\le(4N'(4N'')^2)^2=Z_0$, and
\[
 B(m)\le\#\{(n_1,n_2,n_3,n_4):n_1n_2^2n_3n_4^2=m\}\ll_\delta m^{\delta}
\]
by Lemma~\ref{lem:divisor}. Also
\[
 \sum_mB(m)=\Bigl(\sum_nA(n)\Bigr)^2
 =\Bigl(\sum_{n_1}w(n_1/N')\Bigr)^2\Bigl(\sum_{n_2}w(n_2/N'')\Bigr)^2\le(16N'N'')^2 .
\]
Orthogonality over all characters modulo $Q$ gives
\[
 \sum_{\chi\bmod Q}|G(\chi)|^4=\sum_{\chi\bmod Q}\Bigl|\sum_mB(m)\chi(m)\Bigr|^2
 =\varphi(Q)\sum_{\substack{m\equiv m'\ (\mathrm{mod}\ Q)\\(mm',Q)=1}}B(m)B(m').
\]
For fixed $m$, at most $1+Z_0/Q$ integers $m'\in[1,Z_0]$ satisfy
$m'\equiv m\pmod Q$. Hence the right side is at most
\[
 \varphi(Q)\sum_mB(m)\cdot\max_{m'\le Z_0}B(m')\cdot\Bigl(1+\frac {Z_0}Q\Bigr).
\]
Insert the two bounds.
\end{proof}

\begin{lemma}\label{lem:pw:jacobi}
Let $\ell$ be prime, $\psi$ a nonprincipal and $\nu$ any character modulo
$\ell$, and $\epsilon\in\{\pm1\}$. Then
\[
 \Bigl|\sum_{t\in(\mathbb Z/\ell)^\times}\nu(t)\,\psi(1+\epsilon t)\Bigr|\le\sqrt\ell .
\]
\end{lemma}

\begin{proof}
If $\nu$ is principal the sum is $\sum_{y\in\mathbb F_\ell}\psi(y)-\psi(1)=-1$.
Otherwise substitute $y=-\epsilon t$. The sum becomes
$\nu(-\epsilon)\mathcal J(\nu,\psi)$ with
$\mathcal J(\nu,\psi)=\sum_{y\in\mathbb F_\ell}\nu(y)\psi(1-y)$, both characters
nonprincipal. Then $|\mathcal J(\nu,\psi)|=\sqrt\ell$ if $\nu\psi$ is
nonprincipal and $|\mathcal J(\nu,\psi)|=1$ otherwise \cite[Ch.~8, \S3]{IR}.
\end{proof}

\begin{lemma}\label{lem:pw:FI}
Let $M\ge1$, and let $\mathfrak p_1\ne\mathfrak p_2$ be primes not dividing $M$. Put
$Q=M\mathfrak p_1\mathfrak p_2$. Let $\xi$ be a character modulo
$\mathfrak p_1\mathfrak p_2$. Assume its components modulo $\mathfrak p_1$ and
$\mathfrak p_2$ are both nonprincipal. Let $\tilde\xi$ be the character modulo $Q$
induced by $\xi$, and fix $\epsilon\in\{\pm1\}$. There are coefficients
$c_\epsilon(\alpha)\in\mathbb C$ for characters $\alpha$ modulo $Q$, with
\begin{equation}\label{eq:pw:jac}
 |c_\epsilon(\alpha)|\le\sqrt{\mathfrak p_1\mathfrak p_2},
\end{equation}
such that for all integers $U,V$
\begin{equation}\label{eq:pw:FI}
 \mathbf 1_{(UV,Q)=1}\,\mathbf 1_{M\mid V+\epsilon U}\;
 \xi\Bigl(\frac{V+\epsilon U}{M}\Bigr)
 =\frac{\overline{\xi(M)}}{\varphi(Q)}\sum_{\alpha\bmod Q}
 c_\epsilon(\alpha)\,\alpha(U)\,(\tilde\xi\bar\alpha)(V).
\end{equation}
The left side is read as $0$ when $M\nmid V+\epsilon U$.
\end{lemma}

\begin{proof}
If $(UV,Q)>1$, then $\alpha(U)(\tilde\xi\bar\alpha)(V)=0$ for every $\alpha$,
and both sides vanish. Let $(UV,Q)=1$. By the Chinese remainder theorem each
$\alpha$ is $\alpha=\eta\nu$ with unique $\eta\bmod M$ and
$\nu\bmod \mathfrak p_1\mathfrak p_2$; then $\tilde\xi\bar\alpha=\bar\eta\cdot\xi\bar\nu$.
Define
\[
 c_\epsilon(\alpha)=\eta(-\epsilon)\,\hat f(\nu),\qquad
 \hat f(\nu)=\sum_{y\in(\mathbb Z/\mathfrak p_1\mathfrak p_2)^\times}\bar\nu(y)\,\xi(1+\epsilon y).
\]
Orthogonality on the finite group $(\mathbb Z/M)^\times$, for any $M$, gives
\[
 \frac1{\varphi(M)}\sum_{\eta\bmod M}\eta(-\epsilon U)\bar\eta(V)
 =\mathbf 1_{-\epsilon U\equiv V\ (\mathrm{mod}\ M)}=\mathbf 1_{M\mid V+\epsilon U}.
\]
Fourier inversion of $f(y)=\xi(1+\epsilon y)$ on $(\mathbb Z/\mathfrak p_1\mathfrak p_2)^\times$, at
$y=UV^{-1}$, gives
\[
 \frac1{\varphi(\mathfrak p_1\mathfrak p_2)}\sum_{\nu\bmod \mathfrak p_1\mathfrak p_2}\hat f(\nu)\,\nu(U)\,(\xi\bar\nu)(V)
 =\xi(V)\,\xi(1+\epsilon UV^{-1})=\xi(V+\epsilon U).
\]
Multiply; $\varphi(Q)=\varphi(M)\varphi(\mathfrak p_1\mathfrak p_2)$. The right side of
\eqref{eq:pw:FI} becomes
\[
 \overline{\xi(M)}\,\mathbf 1_{M\mid V+\epsilon U}\,\xi(V+\epsilon U).
\]
If
$M\mid V+\epsilon U$, then $\xi(V+\epsilon U)=\xi(M)\,\xi((V+\epsilon U)/M)$
and $|\xi(M)|=1$. This is \eqref{eq:pw:FI}. Write
$\xi=\xi_{\mathfrak p_1}\xi_{\mathfrak p_2}$ and
$\nu=\nu_{\mathfrak p_1}\nu_{\mathfrak p_2}$. The Chinese remainder theorem
factors $\hat f(\nu)$ into two sums of the shape of
Lemma~\ref{lem:pw:jacobi}, with $\psi=\xi_{\mathfrak p_1}$,
$\xi_{\mathfrak p_2}$. This gives
\eqref{eq:pw:jac}.
\end{proof}

No division by $x$ occurs in \eqref{eq:pw:FI}, and $M$ need not be squarefree,
odd or prime to $x$. For $\epsilon=+1$ this is the sum identity; for
$\epsilon=-1$ the difference identity.

\begin{lemma}\label{lem:pw:sieve}
Let $K_{\rm sv}\ge1$ and $\theta>0$. Let $M\le Y^{K_{\rm sv}}$ be a positive integer and
\[
 \mathcal P_M=\{p\ \text{prime}:\ p\equiv1\ (\mathrm{mod}\ 3),\
 Y^\theta\le p\le2Y^\theta,\ p\nmid M\}.
\]
For $p\in\mathcal P_M$ fix a character $\chi_p\bmod p$ of exact order $3$. Let
$(\varpi_\iota,U_\iota,V_\iota,n_\iota)_{\iota\in\mathcal I}$ be a finite
family with $\varpi_\iota\ge0$ and nonzero integers
$U_\iota,V_\iota,n_\iota$ satisfying $|U_\iota V_\iota n_\iota|\le Y^{K_{\rm sv}}$. If
$Y\ge Y_{\rm sv}(K_{\rm sv},\theta)$, then $|\mathcal P_M|\ge Y^\theta/(3\theta\log Y)$ and
\begin{equation}\label{eq:pw:sieve}
 \sum_{\iota:\,n_\iota\in\mathrm{Cub}}\varpi_\iota
 \le\frac4{|\mathcal P_M|}\sum_\iota\varpi_\iota
 +\frac4{|\mathcal P_M|^2}\sum_{\substack{\mathfrak p_1,\mathfrak p_2\in\mathcal P_M\\\mathfrak p_1\ne \mathfrak p_2}}
 \Bigl|\sum_\iota\varpi_\iota\,\mathbf 1_{(\mathfrak p_1\mathfrak p_2,U_\iota V_\iota)=1}\,
 \chi_{\mathfrak p_1}(n_\iota)\overline{\chi_{\mathfrak p_2}(n_\iota)}\Bigr| ,
\end{equation}
where $\mathrm{Cub}=\{m^3:m\in\mathbb Z\}$ is the set of integer cubes.
\end{lemma}

\begin{proof}
By the prime number theorem in progressions modulo $3$,
\[
 \#\{p\in[Y^\theta,2Y^\theta]:p\equiv1\pmod3\}
 \ge\frac{Y^\theta}{2.5\,\theta\log Y}
\]
for large $Y$. At most $K_{\rm sv}/\theta$ of them divide $M$. This gives the bound
for $|\mathcal P_M|$. For $\iota\in\mathcal I$ put
\[
 \mathcal Z_\iota=\sum_{\substack{p\in\mathcal P_M\\ p\nmid U_\iota V_\iota n_\iota}}\chi_p(n_\iota).
\]
Since $0<|U_\iota V_\iota n_\iota|\le Y^{K_{\rm sv}}$, at most $K_{\rm sv}/\theta$ primes of $\mathcal P_M$
divide $U_\iota V_\iota n_\iota$. If $n_\iota=m^3$ and $p\nmid n_\iota$, then
$\chi_p(n_\iota)=\chi_p(m)^3=1$. So $n_\iota\in\mathrm{Cub}$ gives
$\mathcal Z_\iota\ge|\mathcal P_M|-K_{\rm sv}/\theta\ge|\mathcal P_M|/2$ for large $Y$. Hence
\[
 \mathbf 1_{n_\iota\in\mathrm{Cub}}\le\frac{4|\mathcal Z_\iota|^2}{|\mathcal P_M|^2}
\qquad(\iota\in\mathcal I).
\]
Expand $|\mathcal Z_\iota|^2$. The diagonal terms $\mathfrak p_1=\mathfrak p_2$ equal
$1$, and there are at most $|\mathcal P_M|$ of them. For
$\mathfrak p_1\ne\mathfrak p_2$, the conditions
$\mathfrak p_1,\mathfrak p_2\nmid n_\iota$ may be dropped, because
$\chi_{\mathfrak p_1}(n_\iota)=0$ when $\mathfrak p_1\mid n_\iota$. Multiply by
$\varpi_\iota$, sum over $\iota\in\mathcal I$, and use the triangle inequality in
$(\mathfrak p_1,\mathfrak p_2)$.
\end{proof}

\subsection{Descriptors and the master estimate}\label{sec:pow3:master}

A \emph{slot} is a pair $(s,j)$ with $s\in\{u,v\}$ and $1\le j\le6$. Its mass is
$s_{s,j}$.

\begin{definition}\label{def:pw:desc}
 A \emph{descriptor} is a triple $\mathfrak D=(\sigma,\mathcal G,\mathcal U)$.
Here $\sigma$ is a powered side; $\mathcal G$ is a set of \emph{groups}, each
of the form $g_s=\{(s,1),(s,2)\}$ with $s\in\{u,v\}$; $\mathcal U$ is a set of
\emph{singleton} slots disjoint from $\bigcup\mathcal G$; and
$|\mathcal G|+|\mathcal U|=4$. Its retained set, retained mass, grouped
lengths and generic exponent are
\begin{equation}\label{eq:pw:PSL}
 J=\mathcal U\cup\bigcup\mathcal G,\qquad
 P_J=\sum_{(s,j)\in J}s_{s,j},\qquad
 \mathcal L_{g_s}=s_{s,1}+2s_{s,2},\qquad
 E_{\mathfrak D}=\mathcal C_\sigma+R_\sigma-\tfrac12P_J .
\end{equation}
Write $\widehat{\mathcal U}=\{\{(s,j)\}:(s,j)\in\mathcal U\}$ and
$\mathcal F_{\mathfrak D}=\mathcal G\cup\widehat{\mathcal U}$ for the four
analytic factors.
\end{definition}

So $|J|=4$, $5$ or $6$. There are always four analytic factors, one per
group and one per singleton. Aggregate tails are never
slots. $\mathcal G=\emptyset$ is ordinary $\mathrm{POW3}$.

\begin{proposition}\label{prop:pw:master}
Let $\mathfrak D$ be a descriptor and $\theta>0$ with
\begin{equation}\label{eq:pw:guard}
 W_\sigma+2\theta\ge2\mathcal L_g\qquad(g\in\mathcal G).
\end{equation}
Then
\begin{equation}\label{eq:pw:master}
 D\le\max\bigl(E_{\mathfrak D}+\theta,\ R_\sigma-\theta\bigr)+o(1).
\end{equation}
\end{proposition}

\begin{proof}
Write $C=\mathcal C_\sigma$, $W=W_\sigma$, $R=R_\sigma$ and $\mathfrak r_{s,j}$,
$S_{s,j}$ for the roots and sizes of the represented sides.

\emph{Step 1. Outer configurations.} An outer configuration $\mathfrak o$
consists of the labels of $M$ and all roots $\mathfrak r_{s,j}$ $(s\in\{u,v\},\
1\le j\le d,\ (s,j)\notin J)$. By Lemma~\ref{lem:pw:ledger}(i) and
$\#\{\mathfrak r\sim S\}=S$,
\begin{equation}\label{eq:pw:nout}
 \#\{\mathfrak o\}\le Y^{C}\prod_{(s,j)\notin J}S_{s,j}=Y^{C+R-P_J}.
\end{equation}
Fix $\mathfrak o$. For $(s,j)\in J$ write $\mathfrak r_{s,j}=\mathfrak u_{s,j}$
and $\mathbf u=(\mathfrak u_{s,j})_{(s,j)\in J}$. Then
\begin{equation}\label{eq:pw:UV}
 n_u=U_{\mathfrak o}\prod_{(u,j)\in J}\mathfrak u_{u,j}^{\,j},\qquad
 n_v=V_{\mathfrak o}\prod_{(v,j)\in J}\mathfrak u_{v,j}^{\,j},
\end{equation}
with $U_{\mathfrak o},V_{\mathfrak o}$ fixed. Write $U(\mathbf u),V(\mathbf u)$
for the right sides. Given $\mathfrak o$ and $\mathbf u$,
\eqref{eq:pw:orient} fixes $n_\sigma=V+\epsilon_{\rm pw} U$ and \eqref{eq:pw:Mx}
fixes $x=(n_\sigma/M)^{1/3}$. By Lemma~\ref{lem:divisor},
$\mathfrak r_{\sigma,3}\mathfrak r_{\sigma,6}^2=x/(\mathfrak r_{\sigma,4}\mathfrak r_{\sigma,5})$
has $Y^{o(1)}$ solutions. Pairwise coprimality of
the sides gives $(M,UV)=1$. Hence
\begin{equation}\label{eq:pw:step1}
 \begin{gathered}
 \mathcal B_d\le Y^{o(1)}\sum_{\mathfrak o}N(\mathfrak o),\\
 N(\mathfrak o)=\#\Bigl\{\mathbf u:\
 \begin{aligned}
 &\mathfrak u_{s,j}\sim S_{s,j},\ (M,UV)=1,\ M\mid V+\epsilon_{\rm pw} U,\\
 &(V+\epsilon_{\rm pw} U)/M\ \text{is a positive integer cube}
 \end{aligned}\Bigr\}.
 \end{gathered}
\end{equation}

\emph{Step 2. Positive ambientization.} For $\mathbf u\in\mathbb N^J$ put
\begin{equation}\label{eq:pw:omega}
 \begin{aligned}
 \varpi_{\mathfrak o}(\mathbf u)&=\prod_{(s,j)\in J}w\bigl(\mathfrak u_{s,j}/S_{s,j}\bigr)\\
 &\quad\cdot\mathbf 1_{(M,UV)=1}\,\mathbf 1_{V+\epsilon_{\rm pw} U\ne0}\,
 \mathbf 1_{M\mid V+\epsilon_{\rm pw} U},\\
 n(\mathbf u)&=\frac{V+\epsilon_{\rm pw} U}{M}.
 \end{aligned}
\end{equation}
Since $w\ge0$ and $w=1$ on $[1,2)$, pointwise
\begin{equation}\label{eq:pw:amb}
 N(\mathfrak o)\le\sum_{\mathbf u\in\mathbb N^J}\varpi_{\mathfrak o}(\mathbf u)\,
 \mathbf 1_{n(\mathbf u)\in\mathrm{Cub}}.
\end{equation}
We drop the condition $\gcd(U,V)=1$ and the sharp dyadic windows of the
retained roots, and keep $(M,UV)=1$. We also weaken the positive-cube
condition to: $n(\mathbf u)$ is a nonzero integer cube. These changes only
increase the count. We retain $M$, its labels, every unretained represented root and every
tail in $\mathfrak o$. On the support,
$\mathfrak u_{s,j}\le4S_{s,j}$, so $|UVn|\le Y^{8}$ for large $Y$.

\emph{Step 3. Sieve.} Apply Lemma~\ref{lem:pw:sieve} with $K_{\rm sv}=8$ to the family
$(\varpi_{\mathfrak o}(\mathbf u),U,V,n(\mathbf u))$ restricted to its support.
Then
\begin{equation}\label{eq:pw:split}
 \begin{gathered}
 N(\mathfrak o)\le\frac4{|\mathcal P_M|}\sum_{\mathbf u}\varpi_{\mathfrak o}(\mathbf u)
 +4\max_{\substack{\mathfrak p_1,\mathfrak p_2\in\mathcal P_M\\\mathfrak p_1\ne \mathfrak p_2}}|\mathcal E_{\mathfrak p_1,\mathfrak p_2}(\mathfrak o)|,\\
 \mathcal E_{\mathfrak p_1,\mathfrak p_2}(\mathfrak o)=\sum_{\mathbf u}\varpi_{\mathfrak o}(\mathbf u)
 \mathbf 1_{(\mathfrak p_1\mathfrak p_2,UV)=1}\,\xi(n(\mathbf u)),
 \end{gathered}
\end{equation}
with $\xi=\chi_{\mathfrak p_1}\bar\chi_{\mathfrak p_2}$, a character modulo
$\mathfrak p_1\mathfrak p_2$ whose components are $\chi_{\mathfrak p_1}$ and
$\bar\chi_{\mathfrak p_2}$, both nonprincipal. Also
$|\mathcal P_M|^{-1}\le Y^{-\theta+o(1)}$.

\emph{Step 4. Diagonal.} Split $\mathfrak o=(\mathfrak m,\mathfrak y)$ into the
labels $\mathfrak m$ of $M$ and the represented part $\mathfrak y$. For fixed
$(\mathfrak y,\mathbf u)$ the integers $U,V$ are fixed and
$V+\epsilon_{\rm pw} U\ne0$ on the support. Then
$M\mid V+\epsilon_{\rm pw} U$ leaves at most
$\operatorname{div}(|V+\epsilon_{\rm pw} U|)$ values of $M$, each with $Y^{o(1)}$ label
tuples by Lemma~\ref{lem:pw:ledger}(iv). Hence
\[
 \begin{aligned}
 \sum_{\mathfrak o}\sum_{\mathbf u}\varpi_{\mathfrak o}(\mathbf u)
 &\le Y^{o(1)}\sum_{\mathfrak y}\ \prod_{(s,j)\in J}\sum_{\mathfrak u\ge1}w(\mathfrak u/S_{s,j})\\
 &\le Y^{o(1)}\,Y^{R-P_J}\prod_{(s,j)\in J}4S_{s,j}=Y^{R+o(1)} .
 \end{aligned}
\]
The diagonal part of \eqref{eq:pw:split}, summed over $\mathfrak o$, is
therefore $\le Y^{R-\theta+o(1)}$. Summing first over $\mathfrak m$ would give
$Y^{C+R-\theta}$; the divisor recovery removes $C$.

\emph{Step 5. Fourier expansion.} Fix $\mathfrak o$ and
$\mathfrak p_1\ne\mathfrak p_2$ in $\mathcal P_M$. Put
$Q=M\mathfrak p_1\mathfrak p_2$. Since $\xi(0)=0$ and $(M,UV)=1$ is in
$\varpi_{\mathfrak o}$,
\[
 \mathcal E_{\mathfrak p_1,\mathfrak p_2}(\mathfrak o)=\sum_{\mathbf u}
 \prod_{(s,j)\in J}w\bigl(\mathfrak u_{s,j}/S_{s,j}\bigr)\,
 \mathbf 1_{(UV,Q)=1}\,\mathbf 1_{M\mid V+\epsilon_{\rm pw} U}\,
 \xi\Bigl(\frac{V+\epsilon_{\rm pw} U}M\Bigr).
\]
Apply Lemma~\ref{lem:pw:FI}. Dirichlet characters are completely
multiplicative and vanish on nonunits, and $\mathbf u$ runs over a product
set. So
\begin{equation}\label{eq:pw:E}
 \mathcal E_{\mathfrak p_1,\mathfrak p_2}(\mathfrak o)=\frac{\overline{\xi(M)}}{\varphi(Q)}
 \sum_{\alpha\bmod Q}c_{\epsilon_{\rm pw}}(\alpha)\,
 \alpha(U_{\mathfrak o})\,(\tilde\xi\bar\alpha)(V_{\mathfrak o})
 \prod_{f\in\mathcal F_{\mathfrak D}}\Phi_f\bigl(\gamma_f(\alpha)\bigr),
\end{equation}
where $\gamma_f(\alpha)=\alpha$ if $f$ lies on side $u$ and
$\gamma_f(\alpha)=\tilde\xi\bar\alpha$ if $f$ lies on side $v$, and
\[
 \Phi_{\{(s,j)\}}(\chi)=F_{S_{s,j}}(\chi^j),\qquad
 \Phi_{g_s}(\chi)=G_s(\chi)=F_{S_{s,1}}(\chi)\,F_{S_{s,2}}(\chi^2).
\]
For the group this uses $\chi(\mathfrak u_1)\chi(\mathfrak u_2)^2
=\chi(\mathfrak u_1\mathfrak u_2^2)$. By \eqref{eq:pw:jac},
\begin{equation}\label{eq:pw:Ebound}
 |\mathcal E_{\mathfrak p_1,\mathfrak p_2}(\mathfrak o)|\le
\frac{\sqrt{\mathfrak p_1\mathfrak p_2}}{\varphi(Q)}
 \sum_{\alpha\bmod Q}\ \prod_{f\in\mathcal F_{\mathfrak D}}
 \bigl|\Phi_f(\gamma_f(\alpha))\bigr| .
\end{equation}
This is one family of $\varphi(Q)$ modes. Each $\gamma_f$ is a bijection of the
character group modulo $Q$.

\emph{Step 6. Torsion modes.} Here $Q=M\mathfrak p_1\mathfrak p_2$, as in
Step~5. Let
$\mathcal T_{\rm tor}=\{\alpha:\ \gamma_f(\alpha)^j=\chi_0\ \text{for some}\
f=\{(s,j)\}\in\widehat{\mathcal U}\}$. The character group modulo $Q$ is isomorphic to
$(\mathbb Z/Q)^\times$, whose $j$-torsion has at most $j^{\omega(Q)+1}$
elements. So $|\mathcal T_{\rm tor}|\le4\cdot6^{\omega(Q)+1}=Y^{o(1)}$, as
$Q\le Y^2\cdot4Y^{2\theta}=Y^{O_{d,\theta}(1)}$.
Trivially $|\Phi_f|\le\prod_{(s,j)\in f}4S_{s,j}$. By
Lemma~\ref{lem:pw:ledger}(ii), $\varphi(Q)\ge Q\,Y^{-o(1)}\ge
Y^{W+2\theta-o(1)}$. The modes in $\mathcal T_{\rm tor}$ contribute to
\eqref{eq:pw:Ebound} at most
\[
 \frac{2Y^{\theta}}{Y^{W+2\theta-o(1)}}\,Y^{o(1)}\,4^6Y^{P_J}
 =Y^{P_J-W-\theta+o(1)} .
\]
Sum over $\mathfrak o$ with \eqref{eq:pw:nout}. By
Lemma~\ref{lem:pw:ledger}(iii) the total is
\[
 Y^{C+R-W-\theta+o(1)}\le Y^{R-\theta+o(1)} .
\]

\emph{Step 7. Generic modes.} By Hölder with four exponents $4$,
\[
 \sum_{\alpha\notin\mathcal T_{\rm tor}}\prod_{f\in\mathcal F_{\mathfrak D}}\bigl|\Phi_f(\gamma_f(\alpha))\bigr|
 \le\prod_{f\in\mathcal F_{\mathfrak D}}\Bigl(\sum_{\alpha\notin\mathcal T_{\rm tor}}\bigl|\Phi_f(\gamma_f(\alpha))\bigr|^4\Bigr)^{1/4}.
\]
Singleton $f=\{(s,j)\}$: for $\alpha\notin\mathcal T_{\rm tor}$ the character
$\gamma_f(\alpha)^j$ is nonprincipal. The map $\alpha\mapsto\gamma_f(\alpha)^j$
has fibres of size $\le6^{\omega(Q)+1}$. Lemma~\ref{lem:pw:S4} gives
\[
 \sum_{\alpha\notin\mathcal T_{\rm tor}}|\Phi_f(\gamma_f(\alpha))|^4
 \le6^{\omega(Q)+1}\sum_{\chi\ne\chi_0}|F_{S_{s,j}}(\chi)|^4
 \le\varphi(Q)\,S_{s,j}^{\,2}\,Y^{o(1)} .
\]
Here $Q/\varphi(Q)=Y^{o(1)}$, since $Q\le Y^{O_{d,\theta}(1)}$.
Group $f=g_s$: $\gamma_f$ is a bijection, so by Lemma~\ref{lem:pw:G4}
\[
 \sum_{\alpha}|G_s(\gamma_f(\alpha))|^4=\sum_{\chi\bmod Q}|G_s(\chi)|^4
\le\varphi(Q)\,(S_{s,1}S_{s,2})^2\,Y^{o(1)}\Bigl(1+\frac{4^6Y^{2\mathcal L_{g_s}}}{Q}\Bigr).
\]
Here $Q\ge MY^{2\theta}\ge Y^{W+2\theta}\ge Y^{2\mathcal L_{g_s}}$ by
Lemma~\ref{lem:pw:ledger}(ii) and \eqref{eq:pw:guard}. So the bracket is
$\le1+4^6$. The guard is closed: equality in \eqref{eq:pw:guard} still leaves a
bounded number of shifts $m'\equiv m\pmod Q$. Multiply the four bounds:
\[
 \sum_{\alpha\notin\mathcal T_{\rm tor}}\prod_{f\in\mathcal F_{\mathfrak D}}\bigl|\Phi_f(\gamma_f(\alpha))\bigr|
 \le\varphi(Q)\,Y^{o(1)}\prod_{(s,j)\in J}S_{s,j}^{1/2}
 =\varphi(Q)\,Y^{P_J/2+o(1)} .
\]
In \eqref{eq:pw:Ebound} the factor $\varphi(Q)$ cancels. Since
$\sqrt{\mathfrak p_1\mathfrak p_2}\le2Y^\theta$, the generic modes contribute
$Y^{\theta+P_J/2+o(1)}$. Sum over $\mathfrak o$ with \eqref{eq:pw:nout}:
\[
 Y^{C+R-P_J+P_J/2+\theta+o(1)}=Y^{E_{\mathfrak D}+\theta+o(1)} .
\]

\emph{Step 8. Assembly.} By \eqref{eq:pw:step1}, \eqref{eq:pw:amb},
\eqref{eq:pw:split} and Steps 4, 6, 7,
\[
 \mathcal B_d\le Y^{o(1)}\bigl(Y^{R-\theta}+Y^{R-\theta}+Y^{E_{\mathfrak D}+\theta}\bigr).
\]
Take $\log_Y$.
\end{proof}

The proof uses four fourth moments in every case. With two groups,
\[
 \sum_\alpha|G_{s}G_{s'}\Phi_f\Phi_{f'}|\le
 \|G_s\|_4\|G_{s'}\|_4\|\Phi_f\|_4\|\Phi_{f'}\|_4 ,
\]
over the one mode family; no mixed correlation and no eighth moment enter.
Both guards are imposed at the same $\theta$.

\subsection{Elimination of \texorpdfstring{$\theta$}{theta}}\label{sec:pow3:elim}

\begin{lemma}\label{lem:pw:elim}
Let $0\le k\le2$, let $E,R,\Gamma_1,\dots,\Gamma_k$ be real, and let
$0<\rho_0\le1$. Put
\[
 \theta^*=\max\bigl(0,\tfrac12(R-E),\Gamma_1,\dots,\Gamma_k\bigr),\qquad
 \Theta(\rho_0)=\{m\rho_0:\ m\in\mathbb Z,\ 1\le m\le4/\rho_0\},
\]
and assume $\theta^*\le3$.
Some $\theta\in\Theta(\rho_0)$ satisfies $\theta\ge\Gamma_i$ for all $i$ and
\[
 \max(E+\theta,R-\theta)\le\max\bigl(E,\tfrac12(E+R),E+\Gamma_1,\dots,E+\Gamma_k\bigr)+\rho_0 .
\]
\end{lemma}

\begin{proof}
Let $\theta$ be the least element of $\Theta(\rho_0)$ with $\theta\ge\theta^*$.
It exists since $\max\Theta(\rho_0)>4-\rho_0\ge3$. Then
$\theta^*\le\theta\le\theta^*+\rho_0$. As $\theta\ge\frac12(R-E)$,
$R-\theta\le E+\theta$. So the left side is
$E+\theta\le E+\theta^*+\rho_0$, and $E+\theta^*$ is the displayed maximum.
\end{proof}

\begin{proposition}\label{thm:pw:rows}
Let $0<\rho_0\le1$. There is $Y_{\rm el}=Y_{\rm el}(\varepsilon_0,d,\rho_0)$ such that for
$Y\ge Y_{\rm el}$ and every descriptor $\mathfrak D$:
\begin{enumerate}
\item[(a)] if $\mathcal G=\emptyset$,
\begin{equation}\label{eq:pw:POW3}
 D\le E_{\mathfrak D}+2\rho_0\quad\text{or}\quad D\le\tfrac12(E_{\mathfrak D}+R_\sigma)+2\rho_0;
\end{equation}
\item[(b)] if $\mathcal G=\{g\}$,
\begin{equation}\label{eq:pw:G12}
 D\le E_{\mathfrak D}+2\rho_0\quad\text{or}\quad D\le\tfrac12(E_{\mathfrak D}+R_\sigma)+2\rho_0
 \quad\text{or}\quad D\le E_{\mathfrak D}+\mathcal L_g-\tfrac12W_\sigma+2\rho_0;
\end{equation}
\item[(c)] if $\mathcal G=\{g,g'\}$,
\begin{equation}\label{eq:pw:G1212}
 \begin{gathered}
 D\le E_{\mathfrak D}+2\rho_0\quad\text{or}\quad D\le\tfrac12(E_{\mathfrak D}+R_\sigma)+2\rho_0\quad\text{or}\\
 D\le E_{\mathfrak D}+\mathcal L_g-\tfrac12W_\sigma+2\rho_0\quad\text{or}\quad
 D\le E_{\mathfrak D}+\mathcal L_{g'}-\tfrac12W_\sigma+2\rho_0 .
 \end{gathered}
\end{equation}
\end{enumerate}
Here $E_{\mathfrak D}$, $P_J$, $\mathcal L_g$ are as in \eqref{eq:pw:PSL} and $\mathcal C_\sigma$, $W_\sigma$,
$R_\sigma$ as in \eqref{eq:pw:CWR}.
\end{proposition}

\begin{proof}
Write $\mathcal G=\{g^{(1)},\dots,g^{(k)}\}$ and take
$\Gamma_i=\mathcal L_{g^{(i)}}-\frac12W_\sigma$ in Lemma~\ref{lem:pw:elim}. Its hypothesis
holds: $\mathcal C_\sigma\ge0$ and $P_J\le R_\sigma$ give
$E_{\mathfrak D}\ge\frac12R_\sigma\ge0$. Also $R_\sigma\le2$, since
$r_s\le V_s\le1$ for each represented side. Thus
$\frac12(R_\sigma-E_{\mathfrak D})\le\frac14R_\sigma\le\frac12$; and
 $\Gamma_i\le\mathcal L_{g^{(i)}}\le V_s\le1$. The chosen $\theta\in\Theta(\rho_0)$ satisfies
\eqref{eq:pw:guard}. Proposition~\ref{prop:pw:master} at this $\theta$ gives
\[
 D\le\max\bigl(E_{\mathfrak D},\tfrac12(E_{\mathfrak D}+R_\sigma),E_{\mathfrak D}+\Gamma_1,\dots,E_{\mathfrak D}+\Gamma_k\bigr)+\rho_0+o(1).
\]
The grid $\Theta(\rho_0)$ and the set of descriptors are finite. So the
$o(1)$ is $\le\rho_0$ for $Y\ge Y_{\rm el}$. Since
$E_{\mathfrak D}+\max(0,\Gamma_i)=\max(E_{\mathfrak D},E_{\mathfrak D}+\Gamma_i)$, the
maximum is the largest arm. A bound by a maximum is the stated disjunction.
\end{proof}

Each clause is used together with the direct row \eqref{eq:pair},
$D\le R_\sigma$. So (a) reads
$D\le\min\bigl(R_\sigma,\max(E_{\mathfrak D},\tfrac12(E_{\mathfrak D}+R_\sigma))\bigr)$ up to the
relaxation. A clause may not be replaced by the conjunction of its arms.

\subsection{The installed descriptor families}\label{sec:pow3:families}

Proposition~\ref{thm:pw:rows} holds for every descriptor. The system of
Section~\ref{sec:cert} installs the following finite families. For brevity,
$s_j$ in a descriptor denotes the slot $(s,j)$; its mass is always $s_{s,j}$.

\begin{enumerate}
\item[$\mathrm{POW3}(\sigma)$.] $\mathcal G=\emptyset$, $\mathcal U=J$ any
four of the twelve slots $u_1,\dots,u_6,v_1,\dots,v_6$. For $\sigma=b$ and
$\sigma=c$ all $\binom{12}4=495$ sets are installed. For $\sigma=a$, install
the $486$ sets not containing $\{b_1,c_1,c_4\}$. The nine omitted sets are
valid rows, but were not installed in the finite system. Including them would
remove the displayed witness; the optimum of the enlarged system is not
computed here.
\item[$\mathrm{GROUP12c}(y)$.] $\sigma=c$, $y\in\{a,b\}$, $z$ the other
 represented side, $\mathcal G=\{g_y\}$, $\mathcal U=\{y_4,z_1,z_4\}$. Two
clauses.
\item[$\mathrm{GROUP12ab}(\sigma)$.] $\sigma\in\{a,b\}$, $z$ the represented
 side other than $c$, $\mathcal G=\{g_c\}$, $\mathcal U$ any three of
$c_3,c_4,c_5,c_6,z_1,\dots,z_6$. In total $2\binom{10}3=240$ clauses.
\item[$\mathrm{GROUP1212ab}(\sigma)$.] $\sigma\in\{a,b\}$, $z$ as before,
$\mathcal G=\{g_c,g_z\}$, $\mathcal U$ any two of
$z_3,z_4,z_5,z_6,c_3,c_4,c_5,c_6$. In total $2\binom82=56$ clauses.
\end{enumerate}

For these families, \eqref{eq:pw:CWR} and \eqref{eq:pw:PSL} read as follows.
In $\mathrm{GROUP12c}(y)$,
\[
 \begin{gathered}
 C=c-s_{c,3}-s_{c,6},\qquad R=a+b,\qquad
 P_J=s_{y,1}+s_{y,2}+s_{y,4}+s_{z,1}+s_{z,4},\\
 W=v_c+s_{c,1}+2s_{c,2}+s_{c,4}+2s_{c,5},\qquad
 \mathcal L_{g_y}=s_{y,1}+2s_{y,2}.
 \end{gathered}
\]
In $\mathrm{GROUP12ab}(\sigma)$ and $\mathrm{GROUP1212ab}(\sigma)$,
$C=\sigma-s_{\sigma,3}-s_{\sigma,6}$,
$W=v_\sigma+s_{\sigma,1}+2s_{\sigma,2}+s_{\sigma,4}+2s_{\sigma,5}$,
$R=r_z+r_c$, $\mathcal L_{g_c}=s_{c,1}+2s_{c,2}$, $\mathcal L_{g_z}=s_{z,1}+2s_{z,2}$. Here a
bare side letter is its root mass $r_s=\sum_{j\le d}s_{s,j}$ from
\eqref{eq:pw:coord}.

All clauses of this section carry the one relaxation $2\rho_0$; we take
$\rho_0=\varepsilon_0/2$. Table~\ref{tab:rho} records it.


\section{Moment rows}\label{sec:moment}

\subsection{Setting}\label{sec:mo:setting}

Fix $\varepsilon_0\le4/7$ and $d=d(\varepsilon_0)$ as in
Section~\ref{sec:fold}. Work in a nonempty box
admissible in \eqref{eq:reduction} at scale $Y$, and put
$D=\log_Y\mathcal B_d$. Fix $\sigma\in\{a,b,c\}$, and let $(y,z)$ be the
other sides in alphabetical order. Choose nonempty
$T\subseteq\{1,\dots,6,\mathrm{tail}\}$, with $\kappa_T,\Omega_T$ as in
\eqref{eq:T}. Write $\mathfrak m_i$, $\mathfrak y_i$, $\mathfrak z_i$ for the
roots of the $\sigma$-, $y$- and $z$-sides. Let $M$ be the modulus \eqref{eq:M}, and
\[
 \epsilon_{\rm mo}=-1\ (\sigma=c),\qquad \epsilon_{\rm mo}=+1\ (\sigma\ne c).
\]
Then the side values satisfy $n_y\equiv\epsilon_{\rm mo}\,n_z\pmod M$, as
$M\mid n_\sigma$ and $n_a+n_b=n_c$.
Here $M$ is the cut modulus of \eqref{eq:M}; the cubic-sieve section uses its
own local coefficient $M$. The function $\varphi$ denotes Euler's totient.

A \emph{block} is a pair $\beta=(s,j)$, where $s\in\{y,z\}$ and $1\le j\le d$.
Its root $x_\beta$ is $\mathfrak y_j$ or $\mathfrak z_j$. Let $S_\beta$ be its dyadic size, and
put $j_\beta=j$ and $h_\beta=\log_YS_\beta$. Thus $h_\beta$ is $y_j$ or $z_j$;
it is unrelated to $h_d=d(d+1)/2$.
The block is \emph{ordinary} when $j\le6$. For a set $\mathcal H$ of blocks put
$h_{\mathcal H}=\sum_{\beta\in\mathcal H}h_\beta$, and let
$\mathcal H_y,\mathcal H_z$ be its blocks on $y$ and $z$. Put
\begin{equation}\label{eq:mo:P}
 P_{\rm mo}=\sum_{\beta}h_\beta,
\end{equation}
the sum over all $2d$ blocks: the total root mass of the sides $y$ and $z$.
Tails are included in $P_{\rm mo}$.

All $o(1)$ terms below tend to $0$ as $Y\to\infty$ for fixed $\varepsilon_0$
and $d$ (and $\delta$ in Proposition~\ref{prop:mo:PS23}). They are uniform in
the box, in $(\sigma,T)$ and in $\mathcal H$.

\begin{lemma}\label{lem:mo:size}
For $Y\ge Y_0$,
$Y^{\Omega_T}\le M\le Y^2$ and $\varphi(M)\ge M/\operatorname{div}(M)\ge Y^{\Omega_T-o(1)}$.
Every root $x_\beta$ is prime to $M$. Moreover
$\Omega_T\le1$ and $h_\beta\le1/j_\beta$.
\end{lemma}

\begin{proof}
The lower bound is \eqref{eq:Msize} at scale $Y$. $M$ divides
 $\prod_j\mathfrak m_j^{\,j}<2^{d(d+1)/2}Y^{V_\sigma}\le Y^2$ for large $Y$, by
\eqref{eq:boxconstraints}. Next
$\varphi(M)=M\prod_{p\mid M}(1-1/p)\ge M2^{-\omega(M)}\ge M/\operatorname{div}(M)$,
and $\operatorname{div}(M)=Y^{o(1)}$ for $M\le Y^2$. Coprimality is
Lemma~\ref{lem:coprime}. Finally $\Omega_T\le V_\sigma\le1$ and
$j_\beta h_\beta\le V_s\le1$ by \eqref{eq:boxconstraints}.
\end{proof}

\subsection{Reduction and orthogonality}\label{sec:mo:red}

For a set $\mathcal H$ of blocks, a unit $K$ modulo $M$ and a character
$\psi$ modulo $M$ put
\begin{equation}\label{eq:mo:N}
 \begin{gathered}
 \Sigma_\beta(\psi)=\sum_{S_\beta\le x<2S_\beta}\psi(x),\\
 N_{\mathcal H}(K)=\#\Bigl\{(x_\beta)_{\beta\in\mathcal H}:\
 \begin{aligned}
 &S_\beta\le x_\beta<2S_\beta,\ (x_\beta,M)=1,\\
 &\textstyle\prod_{\beta\in\mathcal H_y}x_\beta^{\,j_\beta}\equiv
 K\prod_{\beta\in\mathcal H_z}x_\beta^{\,j_\beta}\ (\mathrm{mod}\ M)
 \end{aligned}\Bigr\}.
 \end{gathered}
\end{equation}

\begin{lemma}\label{lem:mo:red}
For every set $\mathcal H$ of blocks,
\[
 \mathcal B_d\le Y^{\kappa_T+P_{\rm mo}-h_{\mathcal H}+o(1)}\max_KN_{\mathcal H}(K),
\]
the maximum over the outer data fixed in the proof and over units $K$
modulo the resulting $M$.
\end{lemma}

\begin{proof}
\emph{Enumeration.} Enumerate the roots of $\sigma$ selected by $T$.
Its individual tail roots are among them exactly when $\mathrm{tail}\in T$.
Also enumerate the roots of all blocks outside $\mathcal H$, on $y$ and $z$. A tail block in $\mathcal H$
remains unenumerated. As $\#\{x:S\le x<2S\}=S$, there are
$Y^{\kappa_T}\cdot Y^{P_{\rm mo}-h_{\mathcal H}}$ choices. Fix one. Then $M$ is fixed.

\emph{Congruence.} Write $n_y=\Pi_y\prod_{\mathcal H_y}x_\beta^{\,j_\beta}$ and
$n_z=\Pi_z\prod_{\mathcal H_z}x_\beta^{\,j_\beta}$. The coefficients $n_{y,0},n_{z,0}$
are fixed by the box, so $\Pi_y,\Pi_z$ are fixed. They are units modulo $M$
by Lemmas~\ref{lem:mo:size} and~\ref{lem:coprime}. So
$n_y\equiv\epsilon_{\rm mo} n_z$ gives the congruence of
\eqref{eq:mo:N} with $K\equiv\epsilon_{\rm mo} \Pi_z\Pi_y^{-1}$. The sign
$\epsilon_{\rm mo}$ equals $\epsilon_{\rm cut}$ in \eqref{eq:Ksign}.
The roots $x_\beta$ are units. So the remaining tuples number at most
$N_{\mathcal H}(K)$.

\emph{Recovery.} Given all roots of $y$ and $z$, the side values $n_y,n_z$ are
fixed, hence $n_\sigma$. Divide by $n_{\sigma,0}$ and by the selected part. The
quotient $Q_\sigma=\prod_{i\in[6]\setminus T}\mathfrak m_i^{\,i}\cdot(\prod_{j\ge7}\mathfrak m_j^{\,j})^{1-\mathbf 1_{\mathrm{tail}\in T}}<2^{d^2}Y$
determines the unselected roots of $\sigma$ up to
$\operatorname{div}(Q_\sigma)^{d-1}=Y^{o(1)}$ choices by Lemma~\ref{lem:divisor}.
\end{proof}

\begin{lemma}\label{lem:mo:orth}
\[
 N_{\mathcal H}(K)=\frac1{\varphi(M)}\sum_{\chi\bmod M}\overline{\chi}(K)
 \prod_{\beta\in\mathcal H_y}\Sigma_\beta(\chi^{j_\beta})
 \prod_{\beta\in\mathcal H_z}\Sigma_\beta(\overline\chi^{\,j_\beta}).
\]
\end{lemma}

\begin{proof}
For units $g_1,g_2$,
$\varphi(M)^{-1}\sum_\chi\chi(g_1)\overline\chi(g_2)=\mathbf 1_{g_1\equiv g_2}$.
Take $g_1=\prod_{\mathcal H_y}x_\beta^{\,j_\beta}$,
$g_2=K\prod_{\mathcal H_z}x_\beta^{\,j_\beta}$. Characters vanish on nonunits,
so summing over all $x_\beta$ in the windows imposes $(x_\beta,M)=1$. Expand
by complete multiplicativity.
\end{proof}

\subsection{Moment inputs}\label{sec:mo:inputs}

\begin{lemma}\label{lem:mo:kernel}
Let $j\ge1$. Then $\#\{\chi\bmod M:\chi^j=\chi_0\}\le(2j)^{\omega(M)}$, and every
fibre of $\chi\mapsto\chi^j$ has at most $(2j)^{\omega(M)}$ elements. For
$j\le d$ and $M\le Y^2$ this is $Y^{o(1)}$.
\end{lemma}

\begin{proof}
The character group is isomorphic to $(\mathbb Z/M)^\times$, and the fibres of a
homomorphism are cosets of its kernel. By the Chinese remainder theorem count
$\eta^j=1$ locally. Odd $p^e$: cyclic, at most $\gcd(j,\varphi(p^e))\le j$.
$2^e$: the group is $1$, $C_2$ or $C_2\times C_{2^{e-2}}$, at most
$\gcd(j,2)\gcd(j,2^{e-2})\le2j$. Finally
$\omega(M)\ll\log M/\log\log M$.
\end{proof}

We use the 2009 author-preprint formulation of \cite[Theorem~1]{CS}; the
published article is listed in the bibliography. The displayed bound is
meaningful for $m\ge3$, where $\log\log m>0$. For $m=1,2$ the nonprincipal
character set is empty, so those cases contribute zero.

\begin{theoremA}[{\cite[Theorem~1]{CS}}]\label{thm:mo:CS}
For every integer $m\ge3$ and integers $a,B$ with $B>0$ we have
\[
 \frac1{\varphi(m)}\sum_{\chi\ne\chi_0}\Bigl|\sum_{x=a+1}^{a+B}\chi(x)\Bigr|^4
 \ll8^{\omega(m)}\operatorname{div}(m)(\log m)^3(\log\log m)^7B^2.
\]
\end{theoremA}

Here $\chi$ runs over all characters modulo $m$, including imprimitive
characters. The published reference is J. Number Theory~130 (2010), no.~3,
767--785; the frozen source is the 2009 author preprint.
The window $S\le x<2S$ is $a=S-1$, $B=S$.

\begin{lemma}\label{lem:mo:elem}
Let $r,S,M\in\N$ with $r\ge1$, $S\ge1$ and $S^r\le M$. Then
\[
 \sum_{\chi\bmod M}\Bigl|\sum_{S\le x<2S}\chi(x)\Bigr|^{2r}
 \le(2^r+1)\,\varphi(M)\,S^r\max_{n\le(2S)^r}\operatorname{div}(n)^{r-1}.
\]
\end{lemma}

\begin{proof}
Expand the $2r$-th power and apply orthogonality. The left side is
$\varphi(M)$ times the number of $(\mathbf x,\mathbf x')\in[S,2S)^{2r}$ with
all coordinates prime to $M$ and
$x_1\cdots x_r\equiv x'_1\cdots x'_r\pmod M$. Fix $\mathbf x'$: $S^r$ choices.
Then $n=x_1\cdots x_r<(2S)^r\le2^rM$ lies in a fixed class modulo $M$: at most
$2^r+1$ values. Each $n$ has at most $\operatorname{div}(n)^{r-1}$
factorizations by Lemma~\ref{lem:divisor}.
\end{proof}

\begin{lemma}\label{lem:mo:power}
Let $\beta$ be a block, $j=j_\beta$, $S=S_\beta$, and
$\mathcal E_j=\{\chi\bmod M:\chi^j=\chi_0\}$.
\begin{enumerate}
\item[(i)] $\displaystyle\sum_{\chi\notin\mathcal E_j}|\Sigma_\beta(\chi^j)|^4
 \le\varphi(M)\,S^2\,Y^{o(1)}$.
\item[(ii)] If $r_{\max}\in\N$, $1\le r\le r_{\max}$ and $S^r\le M$, then
 $\displaystyle\sum_{\chi\bmod M}|\Sigma_\beta(\chi^j)|^{2r}\le\varphi(M)\,S^r\,Y^{o(1)}$,
 the $o(1)$ depending on $r_{\max}$.
\end{enumerate}
\end{lemma}

\begin{proof}
(i) For $\chi\notin\mathcal E_j$, $\psi=\chi^j$ is nonprincipal, and each
$\psi$ has at most $(2j)^{\omega(M)}$ preimages
(Lemma~\ref{lem:mo:kernel}). So the sum is at most
$(2j)^{\omega(M)}\sum_{\psi\ne\psi_0}|\Sigma_\beta(\psi)|^4$. For $M\le2$ this
vanishes. For $M\ge3$ apply Theorem~\ref{thm:mo:CS} with $m=M$, $a=S-1$, $B=S$.
Its factor is at most $\operatorname{div}(M)^4(2\log Y)^{10}=Y^{o(1)}$.
(ii) The same fibre bound reduces to all $\psi$; apply
Lemma~\ref{lem:mo:elem}. The total factor is
$(2j)^{\omega(M)}(2^{r_{\max}}+1)\max_{n\le(2S)^r}\operatorname{div}(n)^{r-1}$,
which is $Y^{O_{d,r_{\max}}(1/\log\log Y)}$.
\end{proof}

\begin{lemma}\label{lem:mo:holder}
Let $\mathcal X$ be a set of characters modulo $M$ and $\mathcal H$ a finite
set. For $\beta\in\mathcal H$ let $f_\beta:\mathcal X\to\mathbb C$,
$S_\beta\ge1$, $C_{\rm H}\ge1$ and integers $r_\beta\ge1$ satisfy
\[
 |f_\beta|\le S_\beta,\qquad
 \sum_{\chi\in\mathcal X}|f_\beta(\chi)|^{2r_\beta}\le C_{\rm H}\,\varphi(M)\,S_\beta^{\,r_\beta}.
\]
Let $0\le\vartheta_\beta\le1$ with $\sum_\beta\vartheta_\beta/(2r_\beta)\le1$. Then
\[
 \frac1{\varphi(M)}\sum_{\chi\in\mathcal X}\prod_{\beta\in\mathcal H}|f_\beta(\chi)|
 \le C_{\rm H}^{|\mathcal H|}\prod_{\beta\in\mathcal H}S_\beta^{\,1-\vartheta_\beta/2}.
\]
\end{lemma}

\begin{proof}
$|f_\beta|\le S_\beta^{1-\vartheta_\beta}|f_\beta|^{\vartheta_\beta}$. Apply Hölder to
the functions $|f_\beta|^{\vartheta_\beta}$ $(\vartheta_\beta>0)$ with exponents
$2r_\beta/\vartheta_\beta$, and to the constant $1$ with the complementary
exponent. With $|\mathcal X|\le\varphi(M)$,
\[
 \begin{aligned}
 \sum_{\mathcal X}\prod_\beta|f_\beta|^{\vartheta_\beta}
 &\le\varphi(M)^{1-\sum\vartheta_\beta/(2r_\beta)}\\
 &\quad\cdot\prod_\beta\bigl(C_{\rm H}\varphi(M)S_\beta^{\,r_\beta}\bigr)^{\vartheta_\beta/(2r_\beta)}\\
 &\le C_{\rm H}^{|\mathcal H|}\varphi(M)\prod_\beta S_\beta^{\,\vartheta_\beta/2}.
 \end{aligned}
\]
Multiply by $\prod_\beta S_\beta^{1-\vartheta_\beta}$.
\end{proof}

\subsection{The rows}\label{sec:mo:rows}

\begin{proposition}[CS4]\label{prop:mo:CS4}
Let $\mathcal H$ be a set of at most four blocks, in total over $y$ and $z$.
Then
\begin{equation}\label{eq:mo:CS4}
 D\le \kappa_T+P_{\rm mo}-\min\Bigl(\Omega_T,\frac{h_{\mathcal H}}2\Bigr)+o(1).
\end{equation}
\end{proposition}

\begin{proof}
If $\mathcal H=\emptyset$, then
$N_{\emptyset}(K)=\mathbf 1_{K\equiv1\pmod M}\le1$, and the assertion is
immediate. Write $h=h_{\mathcal H}$. By Lemma~\ref{lem:mo:red} it suffices to
show $N_{\mathcal H}(K)\le Y^{\max(h-\Omega_T,\,h/2)+o(1)}$ uniformly in $K$.
Expand by Lemma~\ref{lem:mo:orth}. Let
$\mathcal E=\bigcup_{\beta\in\mathcal H}\mathcal E_{j_\beta}$, the characters
for which some held power is principal. It contains $\chi_0$.

\emph{Torsion.} $|\mathcal E|\le4(2d)^{\omega(M)}=Y^{o(1)}$ by
Lemma~\ref{lem:mo:kernel}. On $\mathcal E$ use $|\Sigma_\beta|\le S_\beta$. By
Lemma~\ref{lem:mo:size} these characters contribute at most
\[
 \frac{|\mathcal E|}{\varphi(M)}\prod_\beta S_\beta\le Y^{h-\Omega_T+o(1)} .
\]

\emph{Generic.} Let $\mathcal X$ be the complement of $\mathcal E$. For
$\chi\in\mathcal X$ and every $\beta$, $\chi^{j_\beta}$ and
$\overline\chi^{\,j_\beta}$ are nonprincipal, and
$|\Sigma_\beta(\overline\chi^{\,j_\beta})|=|\Sigma_\beta(\chi^{j_\beta})|$. Lemma~\ref{lem:mo:power}(i)
gives the hypothesis of Lemma~\ref{lem:mo:holder} with $r_\beta=2$,
 $C_{\rm H}=Y^{o(1)}$. Take $\vartheta_\beta=1$; then
$\sum_\beta\vartheta_\beta/4=|\mathcal H|/4\le1$. The generic characters
contribute at most $Y^{o(1)}\prod_\beta S_\beta^{1/2}=Y^{h/2+o(1)}$.

\emph{Total.} $N_{\mathcal H}(K)\le Y^{o(1)}(Y^{h-\Omega_T}+Y^{h/2})$. With
Lemma~\ref{lem:mo:red},
\[
 D\le \kappa_T+P_{\rm mo}-h+\max\Bigl(h-\Omega_T,\frac h2\Bigr)+o(1)
 =\kappa_T+P_{\rm mo}-\min\Bigl(\Omega_T,\frac h2\Bigr)+o(1).
\]
\end{proof}

No case distinction in $\Omega_T$ occurs. The proof holds for every $M\ge1$; for
$M\le2$ the generic sum is empty. If $\mathcal H=\emptyset$ the row is
trivial.

\begin{proposition}[PS23]\label{prop:mo:PS23}
Let $0<\delta\le1/2$. Then
\begin{equation}\label{eq:mo:PS23}
 D\le \kappa_T+P_{\rm mo}-\min\Bigl(\frac{P_{\rm mo}}2,\frac{2\Omega_T}3\Bigr)+d\delta+o(1),
\end{equation}
the $o(1)$ depending on $\delta$.
\end{proposition}

\begin{proof}
\emph{Paid blocks.} Let $\mathcal H=\{\beta:h_\beta\ge\delta\}$ and
$h=h_{\mathcal H}$. The other blocks number at most $2d$, so
$P_{\rm mo}-h<2d\delta$.

\emph{Orders.} For $\beta\in\mathcal H$ put
\[
 r_\beta=2\quad(h_\beta>\Omega_T/2),\qquad
 r_\beta=\lfloor\Omega_T/h_\beta\rfloor\quad(h_\beta\le\Omega_T/2).
\]
In the second case $2\le r_\beta\le1/\delta$, by Lemma~\ref{lem:mo:size}, and
$S_\beta^{\,r_\beta}=Y^{r_\beta h_\beta}\le Y^{\Omega_T}\le M$. In both cases
\begin{equation}\label{eq:mo:ratio}
 r_\beta h_\beta\ge\frac{2\Omega_T}3 .
\end{equation}
Indeed $2h_\beta>\Omega_T$ in the first case. In the second, put $t=\Omega_T/h_\beta\ge2$:
for $t<3$, $\lfloor t\rfloor/t=2/t>2/3$; for $t\ge3$,
$\lfloor t\rfloor>t-1\ge2t/3$.

\emph{Moments.} Expand $N_{\mathcal H}(K)$ by Lemma~\ref{lem:mo:orth}. Let
$\mathcal E=\bigcup_{\beta\in\mathcal H}\mathcal E_{j_\beta}$ and $\mathcal X$
its complement. On $\mathcal X$,
\[
 \sum_{\chi\in\mathcal X}|\Sigma_\beta(\chi^{j_\beta})|^{2r_\beta}\le\varphi(M)\,
 S_\beta^{\,r_\beta}\,Y^{o(1)}\qquad(\beta\in\mathcal H),
\]
by Lemma~\ref{lem:mo:power}(i) if $h_\beta>\Omega_T/2$, and by
Lemma~\ref{lem:mo:power}(ii) with $r_{\max}=\lfloor1/\delta\rfloor$ if
$h_\beta\le\Omega_T/2$.

\emph{Budget.} Put $b_{\rm H}=\sum_{\beta\in\mathcal H}1/(2r_\beta)$ and
$\vartheta=\min(1,1/b_{\rm H})$, with $\vartheta=1$ if
$\mathcal H=\emptyset$. Then $\sum_\beta\vartheta/(2r_\beta)\le1$.
By \eqref{eq:mo:ratio},
\[
 \frac\vartheta2\,h=\vartheta\sum_{\beta}\frac{r_\beta h_\beta}{2r_\beta}
 \ge\frac{2\Omega_T}3\,\vartheta b_{\rm H} .
\]
If $b_{\rm H}\le1$, then $\vartheta=1$ and the saving $\vartheta h/2$
equals $h/2$. If $b_{\rm H}>1$, then
$\vartheta b_{\rm H}=1$ and $\vartheta h/2\ge2\Omega_T/3$. In both cases
$\vartheta h/2\ge\min(h/2,2\Omega_T/3)$.

\emph{Generic.} Lemma~\ref{lem:mo:holder} with $\vartheta_\beta=\vartheta$ and
$|\mathcal H|\le2d$ bounds the generic characters by
$Y^{h-\vartheta h/2+o(1)}\le Y^{h-\min(h/2,2\Omega_T/3)+o(1)}$.

\emph{Torsion.} $|\mathcal E|\le2d(2d)^{\omega(M)}=Y^{o(1)}$. As in
Proposition~\ref{prop:mo:CS4} these characters contribute
$Y^{h-\Omega_T+o(1)}$. Since $\min(h/2,2\Omega_T/3)\le\Omega_T$, this is dominated.

\emph{Total.} By Lemma~\ref{lem:mo:red},
\[
 D\le \kappa_T+P_{\rm mo}-\min\Bigl(\frac h2,\frac{2\Omega_T}3\Bigr)+o(1).
\]
Finally $\min(h/2,2\Omega_T/3)\ge\min(P_{\rm mo}/2,2\Omega_T/3)-(P_{\rm mo}-h)/2$
and $(P_{\rm mo}-h)/2<d\delta$.
\end{proof}

If $h_\beta\le \Omega_T/2$, then $S_\beta^{\,2}\le M$, and
Lemma~\ref{lem:mo:power}(ii) with $r=2$ gives the bound of
Lemma~\ref{lem:mo:power}(i). Thus this proof invokes
Theorem~\ref{thm:mo:CS} only for blocks with $h_\beta>\Omega_T/2$. When every
block is at most $\Omega_T/2$, the elementary moment suffices.

For larger blocks the elementary moment with Lemma~\ref{lem:mo:holder}
falls short. Take $\Omega_T=1$ and four block masses
\[
 \frac35,\quad \frac{17}{50},\quad \frac{17}{50},\quad \frac1{50}.
\]
The target saving is $\min(P_{\rm mo}/2,2\Omega_T/3)=13/20$. With the
elementary moment, the available orders satisfy
$r_\beta\le\lfloor\Omega_T/h_\beta\rfloor$, giving
$r=(1,2,2,50)$. The best saving with
per-block exponents is the fractional-knapsack value
\[
 \max_{\substack{0\le\vartheta_i\le1\\\sum_i\vartheta_i/(2r_i)\le1}}
 \frac12\sum_i\vartheta_i h_i=\frac{161}{250},
 \qquad (\vartheta_1,\vartheta_2,\vartheta_3,\vartheta_4)
 =\left(\frac{49}{50},1,1,1\right).
\]
The greedy order is by the ratio $r_i h_i$. A uniform exponent instead gives
$\vartheta=100/101$ and saving $65/101$. Even the per-block optimum is short by
\[
 \frac{13}{20}-\frac{161}{250}=\frac3{500}.
\]
The masses and $\Omega_T$ are homogeneous. Rescale them into the admissible
region $L\le1$. This concerns only that method. It is no counterexample to
Proposition~\ref{prop:mo:PS23}.

\begin{corollary}\label{cor:mo:rows}
Put $\delta=\varepsilon_0/(2d)$.
Fix $Y\ge Y_0$, a pair $(\sigma,T)$, and a set $\mathcal H$ of at most four
ordinary blocks. Then at least one arm of each displayed disjunction holds:
\begin{gather}
 D\le \kappa_T+P_{\rm mo}-\Omega_T+\rho_{\mathrm{CS}}\quad\text{or}\quad
 D\le \kappa_T+P_{\rm mo}-\tfrac12h_{\mathcal H}+\rho_{\mathrm{CS}},\label{eq:mo:CS4or}\\
 D\le \kappa_T+\tfrac12P_{\rm mo}+\rho_{\mathrm{PS}}\quad\text{or}\quad
 D\le \kappa_T+P_{\rm mo}-\tfrac23\Omega_T+\rho_{\mathrm{PS}}\label{eq:mo:PS23or}
\end{gather}
with $\rho_{\mathrm{CS}}=O_d(1/\log\log Y)\le\varepsilon_0/2$ and
$\rho_{\mathrm{PS}}=\varepsilon_0/2+O_{\varepsilon_0,d}(1/\log\log Y)\le\varepsilon_0$.
\end{corollary}

\begin{proof}
Propositions~\ref{prop:mo:CS4} and~\ref{prop:mo:PS23} apply, since
$d\delta=\varepsilon_0/2$. The divisor, prime-factor and logarithmic factors
are each $Y^{O_d(1/\log\log Y)}$. The exponents depend on $d$ and $1/\delta$.
The constants contribute $O(1/\log Y)$. The families are finite. A maximum
gives the disjunction of its arms.
\end{proof}

The installed families are: $\mathrm{CS4}(\sigma,T,\mathcal H)$ with
$|\mathcal H|=4$ ordinary blocks, $3\cdot127\cdot\binom{12}4=188\,595$
descriptors; and $\mathrm{PS23}(\sigma,T)$, $3\cdot127=381$ descriptors. Table~\ref{tab:rho}
records the relaxation of each arm. The certificate fixes the actual held
set $\mathcal H$ of each $\mathrm{CS4}$ clause.

\section{The finite system}\label{sec:cert}

The variables are the $28$ nonnegative exponents
\begin{equation}\label{eq:vars}
 (a,b,c;\,a_1,\dots,a_6,b_1,\dots,b_6,c_1,\dots,c_6;\,v_a,v_b,v_c;\,
 f_a,f_b,f_c;\,D).
\end{equation}
Here $L$, $p_i$, $\tau_s$ and $V_s$ of \eqref{eq:V} are linear forms. The
auxiliary variables $f_s$ encode \eqref{eq:fourier}:
\begin{equation}\label{eq:fourierE}
 \begin{gathered}
 f_s\ge s_2+s_4+s_6,\qquad f_s\ge s_3+s_6,\qquad f_s\ge s_5\qquad(s=a,b,c),\\
 f_a+f_b+f_c\le4L-6D .
 \end{gathered}
\end{equation}
A point satisfies \eqref{eq:fourier} if and only if some $f_s$ satisfy
\eqref{eq:fourierE}; take $f_s$ the maximum. Since
$\mathsf A_{\rm cut}+\pos{\mathsf B_{\rm cut}}
=\max(\mathsf A_{\rm cut},\mathsf A_{\rm cut}+\mathsf B_{\rm cut})$, the
limiting cut \eqref{eq:cut} is the clause
\begin{equation}\label{eq:cutor}
 D\le\mathsf A_{\rm cut}\quad\text{or}\quad
 D\le\mathsf A_{\rm cut}+\mathsf B_{\rm cut} .
\end{equation}

Let $\mathcal S$ be the following system.
\begin{itemize}
\item $29$ linear rows: \eqref{eq:L}, \eqref{eq:root}, \eqref{eq:63},
 \eqref{eq:pair}, the last two rows of \eqref{eq:tailrows} per side, the ten
 rows of \eqref{eq:fourierE}, \eqref{eq:HB3} and \eqref{eq:ARCHorder}.
\item The clauses of Table~\ref{tab:sys}, each with all its arms.
\end{itemize}
All rows are in the limiting form of Table~\ref{tab:rho}: every entry in its
last column is dropped. A clause holds if one of
its arms holds. It is never replaced by the conjunction of its arms.
Coefficients, dyadic bounds, tail enumeration and recovery costs are inside
the rows. So the theorems below are finite statements about \eqref{eq:vars}.

\begin{table}[ht]
\centering\small
\caption{The clauses of $\mathcal S$.}\label{tab:sys}
\begin{tabular}{@{}llrl@{}}
\hline clause & arms & number & source\\\hline
(G1)--(G13) & $2$ & $13$ & Appendix~\ref{app:geom}\\
\eqref{eq:ARCH} & $2$ & $1$ & Corollary~\ref{cor:arch}\\
\eqref{eq:SPor} & $2$ & $30$ & Corollary~\ref{cor:spacing}\\
\eqref{eq:cutor} & $2$ & $1\,560\,195$ & Proposition~\ref{prop:cut}\\
\eqref{eq:mo:CS4or} & $2$ & $188\,595$ & Corollary~\ref{cor:mo:rows}\\
\eqref{eq:mo:PS23or} & $2$ & $381$ & Corollary~\ref{cor:mo:rows}\\
\eqref{eq:pw:POW3} & $2$ & $1476$ & Proposition~\ref{thm:pw:rows}(a)\\
\eqref{eq:pw:G12} & $3$ & $242$ & Proposition~\ref{thm:pw:rows}(b)\\
\eqref{eq:pw:G1212} & $4$ & $56$ & Proposition~\ref{thm:pw:rows}(c)\\
\hline
\end{tabular}
\end{table}

The descriptor sets are those of Sections~\ref{sec:cut}, \ref{sec:spacing},
\ref{sec:pow3:families} and~\ref{sec:moment}. The $242$ one-group clauses are
the two $\mathrm{GROUP12c}$ clauses and the $240$ $\mathrm{GROUP12ab}$
clauses. In total $\mathcal S$ has $29$ linear rows and $1\,750\,989$ clauses.
The certificate below uses only these families.

\subsection{The certificate}

\begin{theorem}[Computer-assisted]\label{thm:cert}
Every nonnegative solution of $\mathcal S$ satisfies $D\le20/37$.
\end{theorem}

\begin{proof}
Write $x$ for the vector \eqref{eq:vars}. Let $e_D$ be the unit vector in the
$D$ coordinate. The certificate is a finite rooted trie. Each internal node
carries one clause of $\mathcal S$ and all its arms as children. Each leaf
$t$ carries a set $I(t)$ of rows. These are the linear rows of
$\mathcal S$, the rows $-x_k\le0$ and the arms on the path from the root to
$t$. No other row is admissible. Write each row as
 $\langle g_i,x\rangle\le b_i$. The certificate gives $\eta_i\ge0$ in
$\mathbb Q$ on $I(t)$. The checker verifies at every leaf
\begin{equation}\label{eq:farkas}
 \sum_{i\in I(t)}\eta_i g_i=e_D,
 \qquad
 \sum_{i\in I(t)}\eta_i b_i\le\frac{20}{37}.
\end{equation}
Let $x\ge0$ satisfy $\mathcal S$. At each internal node $x$ satisfies one arm
of its clause. Descend into that child. So $x$ reaches a leaf $t$, and
every row of $I(t)$ holds at $x$. Hence
\begin{equation}\label{eq:chain}
 D=\langle e_D,x\rangle
 =\sum_{i\in I(t)}\eta_i\langle g_i,x\rangle
 \le\sum_{i\in I(t)}\eta_i b_i
 \le\frac{20}{37}. \qedhere
\end{equation}
\end{proof}

The checker also verifies that no leaf path is a prefix of another. The
clause \eqref{eq:ARCH} may occur at any depth; every \eqref{eq:ARCH} node has
both children. Put
\begin{equation}\label{eq:Lambda}
 \Lambda_U=\max_t\sum_{\substack{i\in I(t)\\ \text{not }-x_k\le0}}\eta_i .
\end{equation}
Each row is encoded in the normalisation of its displayed source and of
Table~\ref{tab:rows}. In that normalisation its right-hand side rises by at
most the entry of Table~\ref{tab:rho}. The checker gives $\Lambda_U=\certLambda$, attained at leaf
$\certLambdaLeaf$. Tables~\ref{tab:cert} and~\ref{tab:fam} give the
statistics. The combined row for the two grouped families in Table~\ref{tab:fam}
records their certificate census together.

\begin{table}[ht]
\centering\small
\caption{The certificate of Theorem~\ref{thm:cert}.}\label{tab:cert}
\begin{tabular}{@{}lr@{}}
\hline leaves (all of kind \texttt{upper}) & $\certLeaves$\\
internal nodes & $\certInternal$\\
nodes & $\certNodes$\\
depth & $\certDepth$\\
infeasible leaves & $\certInfeasible$\\
$\Lambda_U$ & $\certLambda$\\
worst certificate value $\sum_i\eta_i b_i$ & $20/37$\\
leaves with certificate value exactly $20/37$ & $\certTight$\\
largest number of nonzero multipliers on a leaf & $\certMaxTerms$\\
distinct cuts \eqref{eq:cutor} on paths & $\certUNused$\\
\hline
\end{tabular}
\end{table}

\begin{table}[ht]
\centering\small
\caption{Census of the certificate by family. A row of a clause family is
one of its arms.}\label{tab:fam}
\begin{tabular}{@{}lrrrr@{}}
\hline family & \shortstack[r]{internal\\nodes} & \shortstack[r]{distinct clauses\\on paths}
 & \shortstack[r]{distinct rows\\with $\eta_i>0$} & \shortstack[r]{pairs (leaf, row)\\with $\eta_i>0$}\\\hline
\certFamilyTable
\hline
\end{tabular}
\end{table}

\subsection{The optimum}\label{sec:witness}
Let $x^*$ be the point
\begin{equation}\label{eq:wit}
 \begin{gathered}
 a=\tfrac{17}{37},\ a_1=\tfrac7{37},\ a_3=\tfrac{10}{37};\qquad
 b=\tfrac7{37},\ b_1=\tfrac2{37},\ v_b=\tfrac{35}{37};\\
 c=\tfrac{13}{37},\ c_1=\tfrac5{37},\ c_4=\tfrac8{37};\qquad
 f_a=\tfrac{10}{37},\ f_c=\tfrac8{37};\qquad D=\tfrac{20}{37},
 \end{gathered}
\end{equation}
with every unlisted coordinate zero.

\begin{proposition}\label{prop:wit}
$x^*$ satisfies $\mathcal S$. Thus $20/37$ is the exact optimum of the
installed system $\mathcal S$.
\end{proposition}

\begin{proof}
The witness verifier of Appendix~\ref{app:ver} checks every row and every arm
at $x^*$ in exact arithmetic. Table~\ref{tab:wit} gives the result. We check
the binding rows by hand.

\emph{Linear rows.} $L=1$, and
\[
 V_a=\tfrac7{37}+3\cdot\tfrac{10}{37}=1,\qquad
 V_b=\tfrac2{37}+\tfrac{35}{37}=1,\qquad
 V_c=\tfrac5{37}+4\cdot\tfrac8{37}=1 .
\]
So \eqref{eq:ARCHorder} and $\sum_{i\le6}is_i+v_s\le1$ hold with equality.
Also $\tau_a=\tau_c=0$ and $\tau_b=5/37$, so $v_b=7\tau_b$. Then
\eqref{eq:63} holds with equality on all three sides. The pair rows give
$b+c=20/37=D$. For \eqref{eq:fourierE}, $f_a=a_3$, $f_b=0$, $f_c=c_4$, and
$f_a+f_b+f_c=18/37\le28/37=4L-6D$. For \eqref{eq:HB3},
$L-p_3=27/37$.

\emph{Clauses.} One arm suffices. Both arms of \eqref{eq:ARCH} hold. For
(G2), $P_2=24/37$, $Q_2=17/37$ and $L-Q_2=20/37$. For \eqref{eq:SPor} with
$(e,y,z,\ell)=(a,b,c,2)$, $\mathsf B_{\rm sp}=V_a-V_c+c_2=0$ and
$\mathsf A_{\rm sp}=b+c-c_2=20/37$. For \eqref{eq:cutor} with $\sigma=c$,
$T=\{1,4\}$, $R_y=\{1,3\}$, $R_z=\{1\}$: $q=1$,
\[
 \begin{gathered}
 \mathsf A_{\rm cut}=(a-a_1-a_3)+(b-b_1)+c_1+c_4=\tfrac{18}{37},\\
 \mathsf B_{\rm cut}=(a_1+3a_3+b_1)-(c_1+4c_4)=\tfrac{39}{37}-1=\tfrac2{37}.
 \end{gathered}
\]
For \eqref{eq:mo:PS23or} with $\sigma=b$, $T=\{\mathrm{tail}\}$:
$\kappa_T=\tau_b=5/37$, $\Omega_T=v_b$, $P_{\rm mo}=a+c=30/37$, and
$\kappa_T+P_{\rm mo}/2=20/37$. For
\eqref{eq:mo:CS4or} with $\sigma=a$, $T=\{1\}$,
$\mathcal H=\{b_1,b_2,c_1,c_4\}$: $\kappa_T=\Omega_T=7/37$, and
$\kappa_T+P_{\rm mo}-\Omega_T=b+c=20/37$. For \eqref{eq:pw:POW3} with $\sigma=c$,
$J=\{a_1,a_2,a_3,b_1\}$: $C_c=13/37$, $R_c=24/37$, $P_J=19/37$, and
$E_{\mathfrak D}=55/74>20/37$. This is the smallest bound within
$\mathrm{POW3}(c)$.
\end{proof}

\begin{table}[ht]
\centering\small
\caption{Smallest bound at $x^*$ per family. The bound of a clause is its
largest arm. For \eqref{eq:pw:POW3}--\eqref{eq:pw:G1212} it is taken before
intersection with \eqref{eq:pair}. The other linear rows hold.}\label{tab:wit}
\begin{tabular}{@{}lrr@{}}
\hline family & entries & smallest bound\\\hline
\eqref{eq:pair} & $3$ & $20/37$\\
\eqref{eq:HB3} & $1$ & $27/37$\\
(G1)--(G13) & $13$ & $20/37$\\
\eqref{eq:SPor} & $30$ & $20/37$\\
\eqref{eq:cutor} & $1\,560\,195$ & $20/37$\\
\eqref{eq:mo:CS4or} & $188\,595$ & $20/37$\\
\eqref{eq:mo:PS23or} & $381$ & $20/37$\\
$\mathrm{POW3}(a)$, $(b)$, $(c)$ & $486$, $495$, $495$ & $41/74$, $26/37$, $55/74$\\
$\mathrm{GROUP12c}$, group on $a$; on $b$ & $1$; $1$ & $33/37$; $65/74$\\
$\mathrm{GROUP12ab}(a)$; $(b)$ & $120$; $120$ & $21/37$; $26/37$\\
$\mathrm{GROUP1212ab}(a)$; $(b)$ & $28$; $28$ & $21/37$; $26/37$\\
\hline
\end{tabular}
\end{table}

The nine $\mathrm{POW3}(a)$ descriptors with $J\supseteq\{b_1,c_1,c_4\}$ are
valid rows. They are not in $\mathcal S$. At $x^*$ they have $P_J=15/37$,
$E_{\mathfrak D}=39/74<20/37$ and
$\tfrac12(E_{\mathfrak D}+R_a)=79/148<20/37$. Both arms fail at $x^*$.
Leaving out valid rows keeps
Theorem~\ref{thm:cert} true. The point $x^*$ lies in the installed $\mathcal S$.
Adding the nine valid rows removes this witness; the optimum of that enlarged
system is not asserted here. The point is not an $abc$ triple and gives no
lower bound for $N_1(X)$.

\section{Stability and the endpoint}\label{sec:end}

\begin{lemma}\label{lem:stab}
Let $\rho\ge0$ and $x\ge0$. Suppose that, with right-hand sides raised by
$\rho$,
\begin{itemize}
\item every linear row of $\mathcal S$ holds at $x$;
\item every clause of $\mathcal S$ has an arm that holds at $x$.
\end{itemize}
The rows $-x_k\le0$ are not relaxed. Then
\begin{equation}\label{eq:stab}
 D\le\frac{20}{37}+\Lambda_U\rho .
\end{equation}
\end{lemma}

\begin{proof}
Descend the trie as in the proof of Theorem~\ref{thm:cert}, with relaxed
arms. At the leaf, \eqref{eq:chain} with relaxed right-hand sides adds
\[
 \rho\sum_{i\ \text{not }-x_k\le0}\eta_i\le\Lambda_U\rho . \qedhere
\]
\end{proof}

\begin{lemma}\label{lem:endpoint}
Let $0<\mu\le1$. Let $Y_0=Y_0(\varepsilon_0,d)$ be the largest of
$Y_{\rm red}$ of Section~\ref{sec:red}, $\exp(\exp(1))$,
$Y_{\rm sv}(8,\theta)$ for $\theta\in\Theta(\varepsilon_0/2)$,
$Y_{\rm el}(\varepsilon_0,d,\varepsilon_0/2)$ and two further thresholds. The
first is the point beyond which the statements made for large $Y$ in
Sections~\ref{sec:arch}--\ref{sec:moment} hold. The second is the least $Y$
beyond which each displayed vanishing loss, including $\log_Y2^{1+h_d}$, is at
most $\varepsilon_0/2$. It depends only on
$\varepsilon_0$ and $d$. Let $Y\ge Y_0$, and take a
nonempty box admissible in \eqref{eq:reduction} at scale $Y$. Then every
linear row of $\mathcal S$ holds in a $\mu$-free form. So does at least one
arm of every clause. Each right-hand side rises by at
most the amount in Table~\ref{tab:rho}. In
particular \eqref{eq:L} becomes $L\le1+3\varepsilon_0/4$. Consequently
\begin{equation}\label{eq:rho}
 \rho=18\varepsilon_0+o_{\varepsilon_0,d}(1)
\end{equation}
is admissible in Lemma~\ref{lem:stab}, and $\rho\le19\varepsilon_0$ for
$\mu\le1$ and $Y\in[Y_0,X]$.
\end{lemma}

\begin{table}[ht]
\centering\small
\caption{The relaxation carried by each family of $\mathcal S$. Here
$\varepsilon_1=\varepsilon_0$, $\rho_0=\varepsilon_0/2$ and
$\delta=\varepsilon_0/(2d)$.}\label{tab:rho}
\begin{tabular}{@{}p{0.31\textwidth}p{0.33\textwidth}r@{}}
\hline family & source & raised by\\\hline
\eqref{eq:L} & \cite[(6.1)]{BBLT} & $3\varepsilon_0/4$\\
\eqref{eq:root}, \eqref{eq:63}, \eqref{eq:tailrows}, the $f_s$ rows of \eqref{eq:fourierE} & definitions and \eqref{eq:boxconstraints} & $0$\\
\eqref{eq:pair} & Lemma~\ref{lem:app:pair} & $O_d(1/\log\log Y)$\\
Appendix~\ref{app:geom} & Corollary~\ref{cor:app:geom} & $6\varepsilon_0+o_{\varepsilon_0,d}(1)$\\
\eqref{eq:fourier}, the last row of \eqref{eq:fourierE} & Lemma~\ref{lem:app:fourier}, $D$-coefficient $6$ & $18\varepsilon_0+o_{\varepsilon_0,d}(1)$\\
\eqref{eq:cut}, \eqref{eq:cutor} & Proposition~\ref{prop:cut} & $o_{\varepsilon_0,d}(1)$\\
\eqref{eq:ARCHorder} & Lemma~\ref{lem:archorder} & $\varepsilon_0/2$\\
\eqref{eq:ARCH}, each arm & Lemma~\ref{lem:archcover} & $3\varepsilon_0/2+\log_Y2^{1+h_d}$\\
\eqref{eq:HB3} & Proposition~\ref{prop:HB3} & $\varepsilon_0+O_{\varepsilon_0}(1/\log Y)$\\
\eqref{eq:SPor}, each arm & Proposition~\ref{prop:spacing} & $\varepsilon_0+O_d(1/\log\log Y)$\\
\eqref{eq:pw:POW3}--\eqref{eq:pw:G1212}, each arm & Proposition~\ref{thm:pw:rows} & $2\rho_0=\varepsilon_0$\\
\eqref{eq:mo:CS4or}, each arm & Corollary~\ref{cor:mo:rows} & $O_d(1/\log\log Y)$\\
\eqref{eq:mo:PS23or}, each arm & Corollary~\ref{cor:mo:rows} & $\varepsilon_0/2+O_{\varepsilon_0,d}(1/\log\log Y)$\\
\hline
\end{tabular}
\end{table}

\begin{proof}
Run Section~\ref{sec:red} at scale $Y$ with $\lambda$ replaced by $\mu$ and
$\log_Y$ in place of $\log_X$ in \eqref{eq:coord}. In
\eqref{eq:boxconstraints}, $\mu$ occurs only in the radical constraint. Its
upper half is
\begin{equation}\label{eq:Lmu}
 L\le\mu+3\varepsilon_0/4\le1+3\varepsilon_0/4,
\end{equation}
the one row of $\mathcal S$ that depends on $\mu$. Its lower half is
\begin{equation}\label{eq:muL}
 \mu\le L+O_d\Bigl(\frac1{\log Y}\Bigr).
\end{equation}
The proof of \cite[Theorem~1.3]{BBLT} uses the weaker radical bound
\[
 \prod_{j\le d}A_jB_jC_j\le Y^{\mu+3\varepsilon_0}\le Y.
\]
This needs $\varepsilon_0\le(1-\mu)/3$, false at $\mu=1$ for every
$\varepsilon_0>0$. It is the only use of $\mu<1$ in that proof, and the
visible source of its $\mu$-dependent constant. We
keep \eqref{eq:Lmu} as a relaxed row instead.

By \eqref{eq:muL}, $\mu$ leaves the other rows. The other constraints of
\eqref{eq:boxconstraints} are $\mu$-free. So \eqref{eq:root}, \eqref{eq:63},
\eqref{eq:tailrows} and Lemma~\ref{lem:tail} are $\mu$-free and exact. So are
the rows $f_s\ge\cdots$ of \eqref{eq:fourierE}, which define $f_s$. The
constants of the reduction and of Appendix~\ref{app:windows} depend only on
$\varepsilon_0$ and $d$, not on $\mu$. The condition
$\varepsilon_0\le(1-\mu)/3$ is needed for the separate step $ABC\le Y$ in the
proof of \cite[Theorem~1.3]{BBLT}.

\emph{Pair row.} \cite[Proposition~2.3]{BBLT} has no rate in $Y$.
Lemma~\ref{lem:app:pair} at scale $Y$ gives \eqref{eq:pair} with the factor
$\Delta=Y^{O_d(1/\log\log Y)}$. So \eqref{eq:pair} costs
$O_d(1/\log\log Y)$, as in Table~\ref{tab:rho}.

\emph{Cut.} The proof of Proposition~\ref{prop:cut} uses only the dyadic
ranges, \eqref{eq:box} and the fixed-$d$ lemmas. The same argument at scale
$Y$ applies for every $0<\mu\le1$. Its divisor and root-class factors cost
$O_d(1/\log\log Y)$.

\emph{Fourier and geometry.} Lemma~\ref{lem:app:fourier} at scale $Y$
contains no $\mu$; its $\varepsilon_0$-exponent is $3\varepsilon_0$. Our
Fourier row is $\kappa_F$ times the logarithm of that bound, with
\[
 \kappa_F=\frac{[D]_{\text{row}}}{[D]_{\text{lemma}}}=\frac61
 =\frac{[L]_{\text{row}}}{[L]_{\text{lemma}}}=\frac{4}{2/3}=\frac1{1/6}=6
\]
on the $D$, $L$ and root terms. So $3\varepsilon_0$ becomes $18\varepsilon_0$.
Corollary~\ref{cor:app:geom} has $\varepsilon_0$-exponents $4\varepsilon_0$,
$6\varepsilon_0$, and $\mu$ appears as $Y^{\mu}$ and $Y^{\mu-1}$. By
\eqref{eq:muL}, $\mu\to L$ costs $O_d(1/\log Y)$. The rows (G1)--(G13) have
coefficient $1$ on $D$ and $L$, so they rise by at most $6\varepsilon_0$.

\emph{New rows.} Sections~\ref{sec:arch}--\ref{sec:moment} work at scale $Y$
with $\log_Y$ coordinates. They use \eqref{eq:box}, the dyadic ranges, the
coefficient bound and the $\mu$-free constraints of
\eqref{eq:boxconstraints}. So their rows are $\mu$-free. Their errors are
those of Corollary~\ref{cor:arch}, Proposition~\ref{prop:HB3} with
$\varepsilon_1=\varepsilon_0$, Corollary~\ref{cor:spacing},
Proposition~\ref{thm:pw:rows} with $\rho_0=\varepsilon_0/2$, and
Corollary~\ref{cor:mo:rows} with $\delta=\varepsilon_0/(2d)$. For
$Y\ge Y_0$ each is at most $2\varepsilon_0$. The cubic-sieve proofs use only
$M\ge Y^{W_\sigma}$; the coefficient $n_{\sigma,0}\ge1$ is the only
coefficient property used in that bound. The factor $Q/\varphi(Q)$ is
accounted for in Step~7 of Proposition~\ref{prop:pw:master}. The factor six above multiplies
only the printed $3\varepsilon_0$ of the Fourier row. No other row carries it.

These are the entries of Table~\ref{tab:rho}; the largest is \eqref{eq:rho}.
All constants depend only on $\varepsilon_0$ and $d$, not on $\mu$. The
$Y$-dependent errors are the displayed $O_d(1/\log Y)$,
$O_d(1/\log\log Y)$ and $o_{\varepsilon_0,d}(1)$ terms. They include the
constants of the reduction in Section~\ref{sec:red}, of
Lemma~\ref{lem:app:fourier} and of Corollary~\ref{cor:app:geom}.
The $Y_0$ threshold absorbs them.
\end{proof}

By \eqref{eq:rho},
\begin{equation}\label{eq:Lambdarho}
 \Lambda_U\rho=\certLambdaFrac\bigl(18\varepsilon_0+o_{\varepsilon_0,d}(1)\bigr)
 =\certLambdaRho\,\varepsilon_0+o_{\varepsilon_0,d}(1).
\end{equation}
The certificate has no infeasible leaf, so no infeasibility margin enters.

Lemma~\ref{lem:endpoint} removes the restriction $\lambda<1$ of
\cite[Theorem~1.3]{BBLT}. At $\mu=1$ the row \eqref{eq:L} is lost, and
Lemma~\ref{lem:stab} prices the relaxation $\rho$ at $\Lambda_U\rho$ in the
exponent. The constant does not depend on $\mu$. $\mathcal S$, its
certificate and $\Lambda_U$ are free of $\varepsilon_0$ and $d$. So choosing
$\varepsilon_0$ from $\Lambda_U$ is not circular. The parameters are chosen
in this order:
\[
 \begin{gathered}
 \varepsilon\ \longrightarrow\ \text{freeze certificate}\ \longrightarrow\ \varepsilon_0
 \ \longrightarrow\ d=d(\varepsilon_0),\\
 \delta=\frac{\varepsilon_0}{2d},\ \rho_0=\frac{\varepsilon_0}2,
 \ \varepsilon_1=\varepsilon_0\ \longrightarrow\ X.
 \end{gathered}
\]
The grid $\Theta(\rho_0)$ of Lemma~\ref{lem:pw:elim} is finite once $\rho_0$
is fixed. Proposition~\ref{prop:mo:PS23} needs no tail cutoff: $d$ is fixed
before $\delta$, and the paid blocks cost $d\delta=\varepsilon_0/2$.

\section{Proof of Theorem~\ref{thm:main}}\label{sec:proof}

Fix $\varepsilon>0$. Choose $\varepsilon_0>0$ with
$\certBudget\,\varepsilon_0<\varepsilon/2$ and with folded index at least
seven (Section~\ref{sec:fold}). This fixes $d=d(\varepsilon_0)$, then
$\delta$, $\rho_0$ and $\varepsilon_1=\varepsilon_0$ as above. The factors
$8^{\omega(M)}$, $(2j)^{\omega(M)}$ and $\operatorname{div}(N)^{O_d(1)}$ are
$Y^{o(1)}$ for fixed $d$ (Lemmas~\ref{lem:qroot}, \ref{lem:divisor},
\ref{lem:mo:kernel}); they need no budget. Dyadic constants and the count of
boxes are inside \eqref{eq:reduction}.

Let $\mu\le1$ and $Y\le X$. If $\mu\le0$, then $N^*_\mu(Y)=0$, as
$\rad(abc)\ge1\ge c^\mu$. Let $0<\mu\le1$. If $Y<Y_0(\varepsilon_0,d)$, then
$N^*_\mu(Y)\le Y^3\ll_{\varepsilon_0,d}1$. If $Y\ge Y_0$, then for a nonempty
box Lemma~\ref{lem:endpoint} raises every row and arm of $\mathcal S$ by at
most $\rho\le19\varepsilon_0$. So the stability and reduction losses are at
most
\[
 19\Lambda_U\varepsilon_0+\varepsilon_0
 =\certBudget\,\varepsilon_0
 <\frac\varepsilon2.
\]
Thus Theorem~\ref{thm:cert} and Lemma~\ref{lem:stab} give
\begin{equation}\label{eq:boxbound}
 \mathcal B_d\le Y^{20/37+\certLambdaNineteen\,\varepsilon_0} .
\end{equation}
By \eqref{eq:reduction} at scale $Y$, at a further cost $Y^{\varepsilon_0}$,
$N^*_\mu(Y)\ll_\varepsilon Y^{20/37+\varepsilon/2}$. The constant depends
only on $\varepsilon$. Finally \eqref{eq:dyadic} with $\lambda=1$ absorbs
$(\log X)^2$ and the maxima into the remaining $\varepsilon/2$.
\qed

\appendix

\section{BBLT estimates in the half-open convention}\label{app:windows}

We prove the pair, Fourier and geometry bounds of
\cite[Propositions~2.3, 3.1 and~4.1]{BBLT} for the boxes \eqref{eq:box}.
The reduction \cite[Proposition~2.1]{BBLT} is rerun in
Section~\ref{sec:red}. Work at scale $X$ with exponent $\lambda$, in a box
satisfying \eqref{eq:boxconstraints}. Windows are $\mathcal I(S)=[S,2S)\cap\N$
with $S\ge1$ a power of two, so $\#\mathcal I(S)=S$ and $\mathcal I(1)=\{1\}$.
Roots equal to $1$ are allowed. Put
\[
 A=\prod_{i\le d}A_i,\qquad B=\prod_{i\le d}B_i,\qquad C=\prod_{i\le d}C_i,
\]
so $ABC=X^L$. For a side $s$ with roots $\mathfrak s_i\in\mathcal I(S_{s,i})$
put $N_s=\prod_{i\le d}\mathfrak s_i^{\,i}$, so the side value is
$n_{s,0}N_s$. Since $\mathfrak s_i<2S_{s,i}$ and
$\prod_iS_{s,i}^{\,i}\le X$,
\begin{equation}\label{eq:app:size}
 1\le N_s<2^{h_d}X .
\end{equation}
Put
\[
 \Delta=\max_{1\le N<2^{h_d}X}\operatorname{div}(N)^d .
\]
By the maximal order of $\operatorname{div}$,
$\Delta=X^{O_d(1/\log\log X)}$ for $X\ge3$. So $\Delta\ll_{d,\delta'}X^{\delta'}$
for every fixed $\delta'>0$. By Lemma~\ref{lem:divisor}, a product
$\prod_{i\in I}\mathfrak s_i^{\,i}<2^{h_d}X$ with $I\subseteq[d]$ determines
the roots $(\mathfrak s_i)_{i\in I}$ up to $\Delta$ choices. The constants
below depend only on $d$ and $\varepsilon_0$. They do not depend on
$\lambda$, the coefficients or the box.

\begin{lemma}\label{lem:app:power}
Let $j\ge2$ and $m\in\Z\setminus\{0\}$. Then
\[
 \#\{(x,y)\in\N^2:x^j-y^j=m\}\le2(j-1)\operatorname{div}(|m|).
\]
\end{lemma}

\begin{proof}
Put $t=x-y$. Then $t\ne0$, and $t\mid m$ since
$x^j-y^j=t\sum_{i<j}x^{j-1-i}y^i$. So $t$ takes at most
$2\operatorname{div}(|m|)$ values. For fixed $t$,
$x^j-(x-t)^j$ is a polynomial in $x$ of degree $j-1\ge1$ with leading
coefficient $jt\ne0$. So it equals $m$ for at most $j-1$ values of $x$, and
$y=x-t$.
\end{proof}

\begin{lemma}[pair bound]\label{lem:app:pair}
$\mathcal B_d\le\Delta\min(AB,AC,BC)$.
\end{lemma}

\begin{proof}
Fix the $a$- and $b$-roots: exactly $AB$ choices. The equation in
\eqref{eq:box} fixes $c_0N_c$, hence $N_c$. By \eqref{eq:app:size}, the
$c$-roots take at most $\Delta$ values. The pairs $(a,c)$ and $(b,c)$ are
the same, since $b_0N_b=c_0N_c-a_0N_a$ and $a_0N_a=c_0N_c-b_0N_b$.
\end{proof}

\begin{lemma}[Fourier bound]\label{lem:app:fourier}
Let $2\le j,k,\ell\le d$ and put
\[
 \Pi_F=\prod_{j\mid i}A_i\prod_{k\mid i}B_i\prod_{\ell\mid i}C_i,
\]
the products over $i\le d$. Then
\[
 \mathcal B_d\le(3d)^{1/2}\Delta^{5/2}\,\frac{(ABC)^{2/3}}{\Pi_F^{1/6}}
 \ll_{\varepsilon_0,d}X^{3\varepsilon_0}\,\frac{(ABC)^{2/3}}{\Pi_F^{1/6}} .
\]
\end{lemma}

\begin{proof}
\emph{Orthogonality.} Write $e(t)=e^{2\pi it}$. For a side $s$ put
\[
 \Sigma_s(\alpha)=\sum e(\alpha\,n_{s,0}N_s)\qquad(\alpha\in\mathbb R),
\]
the sum over all root tuples of $s$ in the windows of the box. The integral
\[
 \int_0^1\Sigma_a(\alpha)\Sigma_b(\alpha)\overline{\Sigma_c(\alpha)}\,d\alpha
\]
counts the root tuples with $a_0N_a+b_0N_b=c_0N_c$. Dropping the gcd
condition, this is at least $\mathcal B_d$, and at most
$\int_0^1|\Sigma_a\Sigma_b\Sigma_c|$. By H\"older with exponents
$3,3,3$ and Cauchy--Schwarz,
\begin{equation}\label{eq:app:holder}
 \mathcal B_d\le\prod_{s=a,b,c}\Bigl(\int_0^1|\Sigma_s|^3\Bigr)^{1/3},
 \qquad
 \int_0^1|\Sigma_s|^3\le\Bigl(\int_0^1|\Sigma_s|^2\Bigr)^{1/2}
 \Bigl(\int_0^1|\Sigma_s|^4\Bigr)^{1/2}.
\end{equation}
We treat side $a$ with $j$. Sides $b$ and $c$, with $k$ and $\ell$, are the
same. Write $\mathbf a=(\mathfrak a_i)_{i\le d}$ and $N(\mathbf a)=N_a$.
As $a_0\ge1$, $\int_0^1e(\alpha a_0m)\,d\alpha=\mathbf 1_{m=0}$ for $m\in\Z$.

\emph{Second moment.} $\int_0^1|\Sigma_a|^2$ counts pairs
$(\mathbf a,\mathbf a')$ with $N(\mathbf a)=N(\mathbf a')$. Choose
$\mathbf a$ in $A$ ways; then $\mathbf a'$ takes at most $\Delta$ values. So
\[
 \int_0^1|\Sigma_a|^2\le A\Delta .
\]

\emph{Fourth moment.} Let $J=\{i\le d:j\mid i\}$, so $j\in J$. Split
$\mathbf a=(\mathbf a^\circ,\mathbf a^\bullet)$ into the roots with
$i\notin J$ and $i\in J$, and put
\[
 u(\mathbf a^\circ)=\prod_{i\notin J}\mathfrak a_i^{\,i},\qquad
 W(\mathbf a^\bullet)=\prod_{i\in J}\mathfrak a_i^{\,i}.
\]
The tuple $\mathbf a^\circ$ runs over exactly
$A^\circ=\prod_{i\notin J}A_i=A/\prod_{j\mid i}A_i$ values. Cauchy--Schwarz
over $\mathbf a^\circ$ gives
\[
 |\Sigma_a(\alpha)|^2\le A^\circ\,\mathcal T(\alpha),\qquad
 \mathcal T(\alpha)=\sum_{\mathbf a^\circ}\
 \sum_{\mathbf a^\bullet,\mathbf a'^\bullet}
 e\bigl(\alpha a_0\,u(\mathbf a^\circ)
 (W(\mathbf a^\bullet)-W(\mathbf a'^\bullet))\bigr)\ge0 .
\]
Expand $|\Sigma_a|^2$ over pairs $(\mathbf a'',\mathbf a''')$ and
integrate:
\[
 \int_0^1|\Sigma_a|^4\le A^\circ\int_0^1\mathcal T|\Sigma_a|^2
 =A^\circ\,\mathcal N,
\]
Here $\mathcal N$ counts the tuples
$(\mathbf a^\circ,\mathbf a^\bullet,\mathbf a'^\bullet,\mathbf a'',\mathbf a''')$
with
\begin{equation}\label{eq:app:4eq}
 u(\mathbf a^\circ)\bigl(W(\mathbf a^\bullet)-W(\mathbf a'^\bullet)\bigr)
 =N(\mathbf a''')-N(\mathbf a'').
\end{equation}
\emph{Zero case.} If the right side of \eqref{eq:app:4eq} is $0$, then
$W(\mathbf a^\bullet)=W(\mathbf a'^\bullet)$. Choose
$(\mathbf a^\circ,\mathbf a^\bullet)$ and $\mathbf a''$ in $A^2$ ways. Then
$\mathbf a'^\bullet$ and $\mathbf a'''$ take at most $\Delta$ values each.
This gives at most $A^2\Delta^2$ tuples.

\emph{Nonzero case.} Choose $\mathbf a'',\mathbf a'''$ in $A^2$ ways, with
right side $m_0\ne0$. By \eqref{eq:app:size}, $|m_0|<2^{h_d}X$. Each
$\mathfrak a_i$ with $i\notin J$ divides $m_0$, so $\mathbf a^\circ$ takes at
most $\operatorname{div}(|m_0|)^d\le\Delta$ values. Then
$m=m_0/u(\mathbf a^\circ)\ne0$ is fixed, and $|m|<2^{h_d}X$. As $j\mid i$ for
$i\in J$,
\[
 W(\mathbf a^\bullet)=x^j,\quad W(\mathbf a'^\bullet)=y^j,\qquad
 x=\prod_{i\in J}\mathfrak a_i^{\,i/j},\quad
 y=\prod_{i\in J}\mathfrak a_i'^{\,i/j},
\]
with $x,y\in\N$ and $x^j-y^j=m$. Lemma~\ref{lem:app:power} gives at most
$2(j-1)\operatorname{div}(|m|)\le2d\Delta$ pairs $(x,y)$. Since
$x\le x^j<2^{h_d}X$ and $y\le y^j<2^{h_d}X$, Lemma~\ref{lem:divisor} recovers
$\mathbf a^\bullet$ from $x$ and $\mathbf a'^\bullet$ from $y$, with at most
$\Delta$ values each. This gives at most $2dA^2\Delta^4$ tuples.

Hence $\mathcal N\le3dA^2\Delta^4$ and
\[
 \int_0^1|\Sigma_a|^4\le3d\,\Delta^4\,\frac{A^3}{\prod_{j\mid i}A_i}.
\]

\emph{Combination.} By \eqref{eq:app:holder},
\[
 \int_0^1|\Sigma_a|^3\le(3d)^{1/2}\Delta^{5/2}\,
 \frac{A^2}{\bigl(\prod_{j\mid i}A_i\bigr)^{1/2}} .
\]
Multiply the three sides and take the cube root. Finally
$\Delta^{5/2}\ll_{\varepsilon_0,d}X^{3\varepsilon_0}$.
\end{proof}

Take $\log_X$, multiply by $6$, keep only $i\le6$ and choose
$j,k,\ell\in\{2,3,5\}$. Since all coordinates are nonnegative, this gives
the $27$ rows of \eqref{eq:fourier}, each raised by
$18\varepsilon_0+O_{\varepsilon_0,d}(1/\log X)$.

For the geometry bound we import the following lattice-point count, in the
form used in the proof of \cite[Proposition~4.1]{BBLT}.

\begin{theoremA}[{\cite[Lemma~3]{HB}}]\label{thm:app:HB}
Let $a',b',c'$ be nonzero integers with $\gcd(a',b',c')=1$, and let
$H_a,H_b,H_c\ge1$. The number of
$(x,y,z)\in\Z^3$ with $\gcd(x,y,z)=1$, $a'x+b'y=c'z$ and $|x|\le H_a$,
$|y|\le H_b$, $|z|\le H_c$ is
\[
 \ll1+\frac{H_aH_bH_c}{\max\{|a'|H_a,|b'|H_b,|c'|H_c\}},
\]
with an absolute implied constant $C_{\rm HB}$.
\end{theoremA}

\begin{lemma}[geometry bound]\label{lem:app:geom}
Let $I,I',I''\subseteq[d]$ and put
\[
 \Pi_1=\prod_{i\in I}A_i\prod_{i\in I'}B_i\prod_{i\in I''}C_i,\qquad
 \Pi_2=\prod_{i\in I}A_i^{i-1}\prod_{i\in I'}B_i^{i-1}\prod_{i\in I''}C_i^{i-1}.
\]
Then
\[
 \mathcal B_d\le C_{\rm HB}\,\Delta^3\Bigl(\frac{ABC}{\Pi_1}
 +2^{3h_d}X^{-1+\varepsilon_0/2}\,ABC\,\Pi_2\Bigr).
\]
\end{lemma}

\begin{proof}
\emph{Variables.} For a tuple counted by $\mathcal B_d$ put
\[
 \begin{aligned}
 a'&=a_0\prod_{i\notin I}\mathfrak a_i^{\,i},&
 b'&=b_0\prod_{i\notin I'}\mathfrak b_i^{\,i},&
 c'&=c_0\prod_{i\notin I''}\mathfrak c_i^{\,i},\\
 x&=\prod_{i\in I}\mathfrak a_i^{\,i},&
 y&=\prod_{i\in I'}\mathfrak b_i^{\,i},&
 z&=\prod_{i\in I''}\mathfrak c_i^{\,i}.
 \end{aligned}
\]
Then $a'x+b'y=c'z$ and $a',b',c'\ge1$. The integers $a',b',c'$ divide the
three side values, and so do $x,y,z$. The side values are pairwise coprime by
Lemma~\ref{lem:coprime}. So $\gcd(a',b',c')=1$ and $\gcd(x,y,z)=1$.

\emph{Heights.} Put
\[
 H_a=2^{h_d}\prod_{i\in I}A_i^{\,i},\qquad
 H_b=2^{h_d}\prod_{i\in I'}B_i^{\,i},\qquad
 H_c=2^{h_d}\prod_{i\in I''}C_i^{\,i}.
\]
These are at least $1$. Since $\mathfrak a_i<2A_i$, $0<x<H_a$; likewise
$y<H_b$ and $z<H_c$. As $c_0\ge1$ and $\mathfrak c_i\ge C_i$,
\[
 c'H_c\ge\prod_{i\notin I''}C_i^{\,i}\prod_{i\in I''}C_i^{\,i}
 =\prod_{i\le d}C_i^{\,i}\ge X^{1-\varepsilon_0/2}
\]
by \eqref{eq:boxconstraints}.

\emph{Count.} Fix the complementary roots $\mathfrak a_i$ $(i\notin I)$,
$\mathfrak b_i$ $(i\notin I')$ and $\mathfrak c_i$ $(i\notin I'')$: exactly
$ABC/\Pi_1$ choices. They fix $a',b',c'$. By Theorem~\ref{thm:app:HB}, the
number of $(x,y,z)$ is at most
\[
 C_{\rm HB}\Bigl(1+\frac{H_aH_bH_c}{c'H_c}\Bigr)
 \le C_{\rm HB}\Bigl(1+2^{3h_d}X^{-1+\varepsilon_0/2}
 \prod_{i\in I}A_i^{\,i}\prod_{i\in I'}B_i^{\,i}\prod_{i\in I''}C_i^{\,i}\Bigr).
\]
Each of $x,y,z$ is less than $2^{h_d}X$. So it fixes its selected roots up to
$\Delta$ values. Multiply, and use
\[
 \frac{ABC}{\Pi_1}\prod_{i\in I}A_i^{\,i}\prod_{i\in I'}B_i^{\,i}
 \prod_{i\in I''}C_i^{\,i}=ABC\,\Pi_2 .\qedhere
\]
\end{proof}

\begin{corollary}\label{cor:app:geom}
For $I,I',I''\subseteq[d]$,
\[
 \mathcal B_d\ll_{\varepsilon_0,d}\frac{X^{\lambda+4\varepsilon_0}}{\Pi_1}
 +X^{\lambda-1+6\varepsilon_0}\,\Pi_2 .
\]
\end{corollary}

\begin{proof}
By \eqref{eq:boxconstraints}, $ABC\le X^{\lambda+3\varepsilon_0/4}$, and
$\Delta^3\ll_{\varepsilon_0,d}X^{\varepsilon_0/4}$. Lemma~\ref{lem:app:geom}
gives
\[
 \mathcal B_d\ll_{\varepsilon_0,d}\frac{X^{\lambda+\varepsilon_0}}{\Pi_1}
 +X^{\lambda-1+3\varepsilon_0/2}\,\Pi_2 ,
\]
which is stronger, as $X\ge1$.
\end{proof}

This is \cite[Proposition~4.1]{BBLT} in our convention, with the same
exponents. Lemmas~\ref{lem:app:pair}--\ref{lem:app:geom} use the windows,
\eqref{eq:box}, $n_{s,0}\ge1$ and the $\lambda$-free constraints of
\eqref{eq:boxconstraints}. Corollary~\ref{cor:app:geom} also uses
$ABC\le X^{\lambda+3\varepsilon_0/4}$. At scale $Y$ with $\lambda:=\mu$ they
give the entries of Table~\ref{tab:rho} for \eqref{eq:pair},
\eqref{eq:fourier} and (G1)--(G13).

\section{The thirteen min-disjunctions}\label{app:geom}

We work at scale $X$ with exponent $\lambda$; Lemma~\ref{lem:endpoint} uses
the same derivation at scale $Y$ with $\mu$. Put $p_i=a_i+b_i+c_i$,
$L=a+b+c$. Each row is Corollary~\ref{cor:app:geom}
for one triple $I,I',I''\subseteq[6]$. A sum of two terms is at most twice
the larger, so
\[
D\le\max(\lambda+4\varepsilon_0-P,\lambda+6\varepsilon_0-Q)
+O_{\varepsilon_0,d}(1/\log X),
\]
where
$P=\sum_{i\in I}a_i+\sum_{i\in I'}b_i+\sum_{i\in I''}c_i$ and
$Q=1-\sum_{i\in I}(i-1)a_i-\sum_{i\in I'}(i-1)b_i
 -\sum_{i\in I''}(i-1)c_i$. With $C_{\varepsilon_0,d}$ the constant of
Corollary~\ref{cor:app:geom}, the constant and the factor $2$ cost
$\log(2C_{\varepsilon_0,d})/\log X=O_{\varepsilon_0,d}(1/\log X)$. By the lower
bound in \eqref{eq:boxconstraints} and \eqref{eq:coord},
\[
 2^{-O_d(1)}X^\lambda\le\prod_{j\le d}A_jB_jC_j=X^L,
 \qquad \lambda\le L+O_{\varepsilon_0,d}(1/\log X).
\]
Hence
$\min(P,Q)\le L-D+6\varepsilon_0+O_{\varepsilon_0,d}(1/\log X)$.
The certificate uses the limiting form below, free of $\varepsilon_0$; the
slack $6\varepsilon_0+O_{\varepsilon_0,d}(1/\log X)$ is charged in
Lemma~\ref{lem:endpoint} and Table~\ref{tab:rho}. The choices, from the proof
of \cite[Theorem~1.3]{BBLT}:
\[
\begin{array}{ll}
 \text{row }1: & I=I'=I''=\{1,2,4\};\\
 \text{rows }2\text{--}4: & I=I'=\{1,2\},\ I''=\{1,2,3\}\ \text{and permutations};\\
 \text{rows }5\text{--}7: & I=\{1,2\},\ I'=I''=\{1,2,3\}\ \text{and permutations};\\
 \text{rows }8\text{--}10: & I=I'=\{1,2\},\ I''=\{1,3\}\ \text{and permutations};\\
 \text{rows }11\text{--}13: & I=\{1,2\},\ I'=I''=\{1,3\}\ \text{and permutations}.
\end{array}
\]
Written out:
\begin{align}
 P_1&=p_1+p_2+p_4,& Q_1&=1-p_2-3p_4
 \tag{G1}\label{eq:G1}\\
 P_2&=p_1+p_2+a_3,& Q_2&=1-p_2-2a_3
 \tag{G2}\label{eq:G2}\\
 P_3&=p_1+p_2+b_3,& Q_3&=1-p_2-2b_3
 \tag{G3}\label{eq:G3}\\
 P_4&=p_1+p_2+c_3,& Q_4&=1-p_2-2c_3
 \tag{G4}\label{eq:G4}\\
 P_5&=p_1+p_2+a_3+b_3,& Q_5&=1-p_2-2a_3-2b_3
 \tag{G5}\label{eq:G5}\\
 P_6&=p_1+p_2+b_3+c_3,& Q_6&=1-p_2-2b_3-2c_3
 \tag{G6}\label{eq:G6}\\
 P_7&=p_1+p_2+c_3+a_3,& Q_7&=1-p_2-2c_3-2a_3
 \tag{G7}\label{eq:G7}\\
 P_8&=p_1+a_3+b_2+c_2,& Q_8&=1-2a_3-b_2-c_2
 \tag{G8}\label{eq:G8}\\
 P_9&=p_1+b_3+a_2+c_2,& Q_9&=1-2b_3-a_2-c_2
 \tag{G9}\label{eq:G9}\\
 P_{10}&=p_1+c_3+a_2+b_2,& Q_{10}&=1-2c_3-a_2-b_2
 \tag{G10}\label{eq:G10}\\
 P_{11}&=p_1+a_2+b_3+c_3,& Q_{11}&=1-a_2-2b_3-2c_3
 \tag{G11}\label{eq:G11}\\
 P_{12}&=p_1+b_2+a_3+c_3,& Q_{12}&=1-b_2-2a_3-2c_3
 \tag{G12}\label{eq:G12}\\
 P_{13}&=p_1+c_2+a_3+b_3,& Q_{13}&=1-c_2-2a_3-2b_3
 \tag{G13}\label{eq:G13}
\end{align}
and each row is used as the disjunction
$D\le L-P_j$ or $D\le L-Q_j$.


\section{Computer checks}\label{app:ver}

The paper comes with external files. They are the certificate report of
Theorem~\ref{thm:cert} and the
checking programs of Table~\ref{tab:sha}. The rows, witness
\eqref{eq:wit}, and verification procedures are specified in the text. All
arithmetic is exact. No checking program imports a numerical package or an
LP solver.

\subsection{The report}
The file \texttt{final-report.raw.json.xz} decompresses to the report.
\texttt{normalize\_completed\_report.py} recomputes its derived trie
statistics; on this report its output equals its input byte for byte. The
report is one JSON object with fields \texttt{target} ($\texttt{"20/37"}$),
\texttt{system}, \texttt{search\_flags}, \texttt{complete}, \texttt{stats},
\texttt{certificates}, \texttt{survivors}, \texttt{open\_stack} and
\texttt{elapsed\_s}. The lists \texttt{survivors} and \texttt{open\_stack}
are empty. Each certificate has exactly four fields.
\begin{itemize}
\item \texttt{kind}: always \texttt{"upper"}.
\item \texttt{path}: the list of arm labels from the root to the leaf.
\item \texttt{terms}: a list of \texttt{[row-label,"p/q"]}, multipliers
 positive.
\item \texttt{rhs}: the value $\sum_i\eta_i b_i$.
\end{itemize}
The paths form the trie of Theorem~\ref{thm:cert}. A label names one row or
arm of $\mathcal S$ (Table~\ref{tab:rows}). The verifier rebuilds each row
from its label and reads no stored coefficient vector. The metadata string
\texttt{system} also contains the label \texttt{DET}. It names a row
family outside $\mathcal S$. The verifier builds no row from it,
and no path or multiplier in this certificate uses such a row.

\begin{center}\small
\begin{tabular}{@{}>{\raggedright\arraybackslash}p{0.30\textwidth}>{\raggedright\arraybackslash}p{0.60\textwidth}@{}}
\hline program & what it checks\\\hline
\texttt{verify\_tree\_}\penalty0\texttt{report\_clean.py} &
 Standard library only. Strict JSON, duplicate keys rejected; target, system
 string and flags. At each leaf: every label well formed and in
 $\mathcal S$; every term a linear row, a row $-x_k\le0$, or an arm on the
 path; \eqref{eq:farkas} exactly, that is $\sum_i\eta_i g_i=e_D$ and
 $\sum_i\eta_i b_i\le20/37$; the declared \texttt{rhs}. The trie: paths
 distinct and prefix-free; the children of each internal node are exactly
 the arms of one clause, \eqref{eq:ARCH} included at any depth. It
 recomputes Tables~\ref{tab:cert} and~\ref{tab:fam} and $\Lambda_U$. It
 rejects $\certMut$ corrupted reports, among them a missing arm, a missing
 \eqref{eq:ARCH} arm, a tampered multiplier, a row outside its path, a
 forbidden $\mathrm{POW3}(a)$ descriptor and a prefix leaf.\\
\texttt{verify\_tree\_}\penalty0\texttt{report.py} &
 a second replay of \eqref{eq:farkas} and of the node counts\\
\texttt{verify\_witness\_}\penalty0\texttt{clean.py} &
 $x^*$ against every clause of $\mathcal S$ (Table~\ref{tab:wit}); it
 rejects $\certWitMut$ corrupted constructors\\
\texttt{audit\_clause\_}\penalty0\texttt{compiler.py},
\texttt{compiler\_}\penalty0\texttt{differential\_used.py} &
 every row used by the certificate, and every $\mathrm{POW3}$ and grouped
 descriptor, rebuilt by the search program's compiler and compared with the
 verifier's row\\
\hline
\end{tabular}
\end{center}

The driver runs these programs and writes their output:
\begin{center}\small
\texttt{run\_verification.sh}\quad$\longrightarrow$\quad
\texttt{FROZEN\_VERIFICATION.stdout}.
\end{center}
Soundness uses the displayed row definitions and the label map in
Table~\ref{tab:rows}. It also uses the exact Farkas identities at every leaf
and the admissibility and trie checks of the first program. The compiler
differential also compares the verifier's row builder with the search
compiler. The coefficient formulas and derivations in
Sections~\ref{sec:red}--\ref{app:geom} identify those rows with the system
$\mathcal S$.

Table~\ref{tab:rows} gives the labels. Here $s\in\{a,b,c\}$; $\sigma$ is the
modulus or powered side; $(y,z)$ is as in Section~\ref{sec:cut}. A set
$T\subseteq\{1,\dots,6,\mathrm{tail}\}$ is written as the integer
$\sum_{t\in T}2^{t-1}$, with the tail counted as $t=7$. Free sets $R_y,R_z$
are written likewise, in $[1,63]$ or $0$. Slots are written \texttt{a1}, \dots,
\texttt{c6}.

\begin{table}[ht]
\centering\scriptsize
\caption{Verifier labels and their sources.}\label{tab:rows}
\renewcommand{\arraystretch}{1.15}
\begin{tabular}{@{}>{\raggedright\arraybackslash}p{0.42\textwidth}@{\quad}>{\raggedright\arraybackslash}p{0.50\textwidth}@{}}
\hline verifier label & row or arm\\\hline
\texttt{L} & \eqref{eq:L}\\
\texttt{roots<=s}, \texttt{rad-s} & \eqref{eq:root}, \eqref{eq:63}\\
\texttt{tail7-s}, \texttt{val-s} & $v_s\ge7\tau_s$, $\sum_{i\le6}is_i+v_s\le1$ of \eqref{eq:tailrows}\\
\texttt{pair-xy} & \eqref{eq:pair}\\
\texttt{fourier-s-j}, $j=0,1,2$; \texttt{fourier-sum} & the rows $f_s\ge s_2+s_4+s_6$, $s_3+s_6$, $s_5$; the last row of \eqref{eq:fourierE}\\
\texttt{HB3} & \eqref{eq:HB3}\\
\texttt{ARCH:a<=c}, \texttt{ARCH:b<=c} & \eqref{eq:ARCHorder}\\
\texttt{ARCH:max-a}, \texttt{ARCH:max-b} & the arms of \eqref{eq:ARCH}\\
\texttt{G0}; \texttt{G1-s}; \texttt{G2-st}; \texttt{G3-s}; \texttt{G4-s};\newline suffix \texttt{:0}, \texttt{:1} & (G1); (G2)--(G4); (G5)--(G7), $st=ab,bc,ca$; (G8)--(G10); (G11)--(G13); the $P$ arm, the $Q$ arm\\
\texttt{SP:e:y:z:k:0}, \texttt{:1} & The field $k$ equals the spacing exponent $\ell$; the two arms of \eqref{eq:SPor}\\
\texttt{UN:}$\sigma$\texttt{:}$R_y$\texttt{:}$R_z$\texttt{:}$T$\texttt{:00}, \texttt{:10} & $D\le\mathsf A_{\rm cut}$, $D\le\mathsf A_{\rm cut}+\mathsf B_{\rm cut}$ of \eqref{eq:cutor}\\
\texttt{CS4:}$\sigma$\texttt{:}$T$\texttt{:}$\mathcal H$\texttt{:W}, \texttt{:h2} & the two arms of \eqref{eq:mo:CS4or}; $T$ is the integer mask above, and $\mathcal H$ is the four comma-separated, lexicographically sorted held slots\\
\texttt{PS:}$\sigma$\texttt{:}$T$\texttt{:half}, \texttt{:2W3} & the two arms of \eqref{eq:mo:PS23or}\\
\texttt{POW3}$\sigma$\texttt{:}$J$\texttt{:S}, \texttt{:avg} & $D\le E_{\mathfrak D}$, $D\le\frac12(E_{\mathfrak D}+R_\sigma)$ of \eqref{eq:pw:POW3}\\
\texttt{GROUP12c:}\dots\newline \texttt{GROUP12a:c1c2sq+}\dots\newline \texttt{GROUP12b:c1c2sq+}\dots\newline suffix \texttt{:S}, \texttt{:avg}, \texttt{:gate} & the three arms of \eqref{eq:pw:G12}; after \texttt{c1c2sq+} come three canonical singleton slots separated by \texttt{+}\\
\texttt{GROUP1212a:b1b2sq+c1c2sq+}\dots\newline \texttt{GROUP1212b:a1a2sq+c1c2sq+}\dots\newline suffix \texttt{:S}, \texttt{:avg}, \texttt{:gate-}$s$ & the four arms of \eqref{eq:pw:G1212}; after the two fixed squared groups come two canonical singleton slots separated by \texttt{+}\\
\texttt{nonneg:x} & $-x\le0$\\
\hline
\end{tabular}
\end{table}

\begin{table}[ht]
\centering\small
\caption{SHA-256 of the files this paper depends on.}\label{tab:sha}
\begin{tabular}{@{}>{\raggedright\arraybackslash}p{0.44\textwidth}@{\quad}l@{}}
\hline file & SHA-256\\\hline
\texttt{MANIFEST.sha256} (20/37 package) & \texttt{\shaManifestI}\\
 & \texttt{\shaManifestII}\\
\texttt{final-report.raw.json.xz} & \texttt{\shaReportXzI}\\
 & \texttt{\shaReportXzII}\\
report (decompressed, normalized) & \texttt{\shaReportI}\\
 & \texttt{\shaReportII}\\
\texttt{normalize\_completed\_report.py} & \texttt{\shaNormI}\\
 & \texttt{\shaNormII}\\
\texttt{verify\_tree\_report\_clean.py} & \texttt{\shaCleanI}\\
 & \texttt{\shaCleanII}\\
\texttt{verify\_tree\_report.py} & \texttt{\shaTreeI}\\
 & \texttt{\shaTreeII}\\
\texttt{verify\_witness\_clean.py} & \texttt{\shaWitI}\\
 & \texttt{\shaWitII}\\
\texttt{audit\_clause\_compiler.py} & \texttt{\shaAuditI}\\
 & \texttt{\shaAuditII}\\
\texttt{compiler\_differential\_used.py} & \texttt{\shaDiffI}\\
 & \texttt{\shaDiffII}\\
\texttt{run\_verification.sh} & \texttt{\shaRunI}\\
 & \texttt{\shaRunII}\\
\texttt{FROZEN\_VERIFICATION.stdout} & \texttt{\shaStdoutI}\\
 & \texttt{\shaStdoutII}\\
\hline
\end{tabular}
\end{table}
\clearpage



\section*{Statements and declarations}

\textbf{Competing interests.} The author has no competing interests.

\textbf{Funding.} No funding was received for this work.

\textbf{Data availability.} The certificates, the checking programs of
Appendix~\ref{app:ver} and their outputs accompany this paper. Their SHA-256
values are in Table~\ref{tab:sha}.

\begin{thebibliography}{10}

\bibitem{Be}
C.~Bernert, \emph{The exceptional set in the abc conjecture}, preprint (2025),
arXiv:2506.13364v1.

\bibitem{BBLT}
C.~Bernert, T.~Browning, J.~D. Lichtman, and J.~Ter\"av\"ainen,
\emph{Bounds on the exceptional set in the abc conjecture}, to appear in
C.~R.~Math.; arXiv:2410.12234v2, revised 9 May 2026.

\bibitem{BLT}
T.~Browning, J.~D. Lichtman, and J.~Ter\"av\"ainen,
\emph{Bounds on the exceptional set in the abc conjecture}, preprint (2024),
arXiv:2410.12234v1; earlier version of \cite{BBLT}.

\bibitem{CS}
Todd Cochrane and Sanying Shi, \emph{The congruence $x_1x_2\equiv x_3x_4\pmod m$
and mean values of character sums}, J. Number Theory \textbf{130} (2010),
no.~3, 767--785,
\url{https://doi.org/10.1016/j.jnt.2009.08.006}.

\bibitem{dB}
N.~G. de Bruijn, \emph{On the number of integers $\le x$ whose prime factors
divide $n$}, Illinois J. Math. \textbf{6} (1962), 137--141,
\url{https://doi.org/10.1215/ijm/1255631814}.

\bibitem{GMRR}
O.~Gorodetsky, K.~Matom\"aki, M.~Radziwi\l\l, and B.~Rodgers, \emph{On the
variance of squarefree integers in short intervals and arithmetic
progressions}, Geom. Funct. Anal. \textbf{31} (2021), 111--149,
\url{https://doi.org/10.1007/s00039-021-00557-5}.

\bibitem{HB}
D.~R. Heath-Brown, \emph{Diophantine approximation with square-free numbers},
Math. Z. \textbf{187} (1984), 335--344,
\url{https://doi.org/10.1007/BF01161951}.

\bibitem{HBcube}
D.~R. Heath-Brown, \emph{Counting square-full solutions to $x+y=z$},
preprint (2026), arXiv:2601.07817v1.

\bibitem{IR}
K.~Ireland and M.~Rosen, \emph{A classical introduction to modern number
theory}, 2nd ed., Grad. Texts in Math. \textbf{84}, Springer, New York,
1990, \url{https://doi.org/10.1007/978-1-4757-2103-4}.

\bibitem{Ka}
D.~M. Kane, \emph{On the number of ABC solutions with restricted radical sizes},
J. Number Theory \textbf{154} (2015), 32--43,
\url{https://doi.org/10.1016/j.jnt.2014.09.021}.

\bibitem{Li}
R.~Li, \emph{On the exceptional set in the abc conjecture}, preprint (2025),
arXiv:2507.02885v1.

\end{thebibliography}

\end{document}
